\documentclass[10pt,a4paper]{article}
\usepackage[utf8]{inputenc}
\usepackage[T1]{fontenc}
\IfFileExists{ngerman.ldf}{\usepackage[ngerman]{babel}}{}
\IfFileExists{lmodern.sty}{\usepackage{lmodern}}{}
\usepackage{amsmath,amssymb,mathtools}
\usepackage[a4paper,margin=2.2cm]{geometry}
\usepackage{booktabs,tabularx,array,longtable}
\usepackage[shortlabels]{enumitem}
\usepackage{xcolor}
\usepackage[most]{tcolorbox}
\usepackage[colorlinks=true,linkcolor=blue!50!black,citecolor=blue!50!black,urlcolor=blue!50!black]{hyperref}

\tcbset{colback=blue!3,colframe=blue!50!black,fonttitle=\bfseries\small,fontupper=\small,
        left=4pt,right=4pt,top=2pt,bottom=2pt}
\newcommand{\tp}[1]{[\![#1]\!]}
\newcommand{\Li}{\tp{1/\varepsilon}^{-1}}
\providecommand{\enquote}[1]{\glqq#1\grqq}
\providecommand{\glqq}{,,}\providecommand{\grqq}{``}
\setlength{\parindent}{0pt}\setlength{\parskip}{3pt}

\title{\textbf{FourierGrating}\\[4pt]\large Implementierungsplan: Fourier-Modal-Methode mit
Differentialmethode (implizites Runge--Kutta) und RCWA\\ für 1D- und 2D-periodische, isotrope und
anisotrope Gitter, inkl.\ Nahfeld\\[4pt]\normalsize C++20/CUDA, AlmaLinux 10, NVIDIA RTX 6000 Ada}
\author{Spezifikation für die Implementierung durch einen KI-Agenten}
\date{Version 2.15 (Teile 5--6: Erweiterungen M12--M19; Teil 7: Zylinderkoordinaten M20--M21)}

\begin{document}
\maketitle
\tableofcontents
\clearpage
%=====================================================================
\section{Zweck, Umfang und Lesehinweise}\label{sec:scope}
%=====================================================================

Dieses Dokument ist die verbindliche technische Spezifikation für die Implementierung von
\textbf{FourierGrating}, einem Löser für die Beugung elektromagnetischer Wellen an periodischen
Strukturen nach der Fourier-Modal-Methode (FMM). Es richtet sich an einen KI-Agenten (oder
Entwickler), der die Implementierung ohne weitere Rückfragen beginnen soll. Alle Formeln sind
hergeleitet oder auf geprüfte Quellen zurückgeführt; Stellen, an denen die Vorlage
(Nevière/Popov, \emph{Light Propagation in Periodic Media}, 2003) fehlerhaft ist, sind in
Abschnitt~\ref{sec:errata} aufgeführt und im Text bereits korrigiert.

\subsection{Funktionsumfang}
\begin{center}
\begin{tabular}{lll}
\toprule
Variante & Periodizität & Material \\
\midrule
1D-isotrop    & $x$ (eine Fouriertransformation, FT) & $\varepsilon$ skalar, $\mu=1$ \\
1D-anisotrop  & $x$                                  & $\overline{\overline\varepsilon}$ voller komplexer $3\times3$-Tensor \\
2D-isotrop    & $x$ und $z$ (zwei FT)               & $\varepsilon$ skalar \\
2D-anisotrop  & $x$ und $z$                          & $\overline{\overline\varepsilon}$ voller Tensor \\
\bottomrule
\end{tabular}
\end{center}
Für jede Variante gibt es zwei Löser für den modulierten Bereich:
\begin{enumerate}[nosep]
  \item \textbf{RK}: Differentialmethode, numerische Integration der Propagationsgleichung
        $dF/dy=M(y)F$ mit einem \emph{impliziten} Runge--Kutta-Verfahren (Gauß--Legendre, optional
        Radau IIA), stabilisiert durch die S-Matrix.
  \item \textbf{RCWA}: Zerlegung in Scheiben mit $y$-unabhängigem $M$ (Treppenapproximation),
        Eigenwertzerlegung je Scheibe, stabile Scheiben-S-Matrix.
\end{enumerate}
Beide Löser verwenden dieselbe Fourierfaktorisierung (FFF, \enquote{Fast Fourier Factorization}
mit Normalenvektorfeldern), dieselben Randbedingungen, dieselbe S-Matrix-Verkettung und dieselbe
Nahfeldrekonstruktion. Das Nahfeld wird jeweils in einem \textbf{horizontalen} ($y=\text{const}$)
und einem \textbf{vertikalen} Schnitt ($z=\text{const}$, Ebene $x$--$y$) ausgegeben.

Zielsystem: AlmaLinux 10 (x86\_64, GCC 14), C++20, CUDA (NVIDIA RTX 6000 Ada,
Compute Capability 8.9, 48\,GB). Ein reiner CPU-Pfad ist Pflicht (Referenz, Tests, CI).

\subsection{Nicht-Ziele (vorerst)}
Magnetische Anisotropie ($\overline{\overline\mu}\neq\mathbb 1$) wird in den Formeln mitgeführt, aber erst
nach Meilenstein M9 implementiert. Nichtlineare Optik, Photonische-Kristall-Bandstrukturen,
Zylinderkoordinaten und Polologie (Buchkapitel X--XV) sind nicht Teil dieses Plans.

\subsection{Lesehinweise für den Agenten}
\begin{itemize}[nosep]
  \item Abschnitte \ref{sec:conv}--\ref{sec:field} enthalten die Mathematik, Abschnitt
        \ref{sec:arch} die Software-Architektur, Abschnitt \ref{sec:validation} die Tests mit
        Referenzwerten, Abschnitt \ref{sec:milestones} die Reihenfolge der Arbeit.
  \item Jede Formel, die implementiert wird, bekommt einen Unit-Test (Abschnitt~\ref{sec:validation}).
  \item Bei Widerspruch zwischen diesem Dokument und dem Buch gilt dieses Dokument.
  \item Gleichungsnummern des Buches werden als \enquote{Buch (IV.21)} zitiert.
\end{itemize}

%=====================================================================
\section{Konventionen und normierte Einheiten}\label{sec:conv}
%=====================================================================

\subsection{Zeitabhängigkeit, Koordinaten, Schichtung}
\begin{itemize}[nosep]
  \item Zeitabhängigkeit $e^{-i\omega t}$; Felder sind komplexe Amplituden.
  \item Periodizität in $x$ (Periode $d_x$) und -- im 2D-Fall -- in $z$ (Periode $d_z$). Die
        Schichtung (Stapelrichtung) ist $y$. Bei 1D-Gittern ist die Struktur in $z$ invariant.
  \item Das \textbf{Substrat} (Medium 1) belegt $y<0$, das \textbf{Superstrat} (Medium $M$) belegt
        $y>y_{\text{top}}$. Dazwischen liegen Schichten $j=2,\dots,M-1$ mit Grenzen
        $0=y_1<y_2<\dots<y_{M-1}=y_{\text{top}}$. Licht fällt von oben ein.
  \item \enquote{Abwärts} (down, Amplitude $a$) heißt Ausbreitung bzw.\ Abklingen in $-y$,
        \enquote{aufwärts} (up, Amplitude $b$) in $+y$.
\end{itemize}

\subsection{Normierung}
Mit $k_0=\omega/c=2\pi/\lambda_0$ und $Z_0=\sqrt{\mu_0/\varepsilon_0}$ werden eingeführt:
\begin{equation}
  \tilde{\mathbf H}=Z_0\mathbf H,\qquad
  \hat x=k_0x,\ \hat y=k_0y,\ \hat z=k_0z,\qquad
  \hat\alpha=\alpha/k_0,\ \hat\beta=\beta/k_0,\ \hat\gamma=\gamma/k_0 .
  \label{eq:norm}
\end{equation}
Mit relativen Materialtensoren $\overline{\overline\varepsilon}_r$, $\overline{\overline\mu}_r$ lauten die
Maxwellgleichungen dann
\begin{equation}
  \widehat{\operatorname{rot}}\,\mathbf E=i\,\overline{\overline\mu}_r\tilde{\mathbf H},\qquad
  \widehat{\operatorname{rot}}\,\tilde{\mathbf H}=-i\,\overline{\overline\varepsilon}_r\mathbf E .
  \label{eq:maxwell_norm}
\end{equation}
Herleitung: $\operatorname{rot}\mathbf E=i\omega\mu_0\mu_r\mathbf H=i\omega\sqrt{\mu_0\varepsilon_0}\,\mu_r\tilde{\mathbf H}
=ik_0\mu_r\tilde{\mathbf H}$ und
$\operatorname{rot}\tilde{\mathbf H}=Z_0(-i\omega\varepsilon_0\varepsilon_r\mathbf E)=-ik_0\varepsilon_r\mathbf E$;
anschließend $\partial/\partial x=k_0\,\partial/\partial\hat x$.
\textbf{Im Code werden ausschließlich normierte Größen verwendet} (Hut wird im Folgenden weggelassen;
alle Längen in Einheiten von $1/k_0$). Benutzereingaben (µm, Grad) werden an der Schnittstelle
umgerechnet. Damit entsprechen alle Buchformeln dem Fall $\omega=1$, $\mu_0=\varepsilon_0=1$.

\subsection{Wellenzahlen, Wurzelzweig}
\begin{equation}
  \alpha_n=\alpha_0+nK_x,\quad K_x=\frac{2\pi}{k_0d_x},\qquad
  \gamma_m=\gamma_0+mK_z,\quad K_z=\frac{2\pi}{k_0d_z}
  \label{eq:alpha_gamma}
\end{equation}
(im 1D-Fall $\gamma_m\equiv\gamma_0$). In einem isotropen homogenen Medium mit $\varepsilon$:
\begin{equation}
  \beta_{nm}=\sqrt{\varepsilon-\alpha_n^2-\gamma_m^2},\qquad
  \operatorname{Im}\beta\ge0,\ \text{und}\ \operatorname{Re}\beta\ge0\ \text{falls}\ \operatorname{Im}\beta=0 .
  \label{eq:beta_branch}
\end{equation}
\textbf{Implementierung:}\\ \verb!b = sqrt(z); if (b.imag()<0 || (b.imag()==0 && b.real()<0)) b = -b;!
\textbf{Rayleigh-Anomalie:} Falls $|\beta|<\eta_{\text{RA}}$ (Standard $10^{-10}$), wird
$\beta\to\beta+i\eta_{\text{RA}}$ gesetzt und eine Warnung protokolliert.

\subsection{Einfallende Welle}
$\Phi$ = Winkel zwischen $\mathbf k$ und der $y$-Achse, $\Psi$ = Azimut der Projektion auf die
$xz$-Ebene gegen $x$, $\delta$ = Polarisationswinkel (Buch Kap.~IX, korrigiert):
\begin{equation}
  \alpha_0=\nu_M\sin\Phi\cos\Psi,\qquad \beta_0=\nu_M\cos\Phi,\qquad \gamma_0=\nu_M\sin\Phi\sin\Psi,
  \qquad \mathbf k^{(i)}=(\alpha_0,-\beta_0,\gamma_0),
  \label{eq:incidence}
\end{equation}
($\nu_M$: reelle Brechzahl des Superstrats; Buch (IX.1) hat fälschlich $\cos\Psi$ in $\beta_0$),
\begin{equation}
  \mathbf A=\begin{pmatrix}\sin\Psi\cos\delta+\sin\delta\cos\Phi\cos\Psi\\ \sin\delta\sin\Phi\\
  \cos\Phi\sin\Psi\sin\delta-\cos\delta\cos\Psi\end{pmatrix},\qquad |\mathbf A|=1,\ \mathbf A\cdot\mathbf k^{(i)}=0 .
  \label{eq:polarization}
\end{equation}
$\delta=0$: $\mathbf A\perp$ Einfallsebene (\enquote{s}); $\delta=90^\circ$: in der Einfallsebene (\enquote{p}).
Bei Normaleinfall gilt $\Psi=\pi$, dann ist $\delta$ der Winkel zwischen $\mathbf A$ und $z$.
1D-in-plane: $\Psi=0$ bzw.\ $\pi$, $\gamma_0=0$; TE $\equiv\delta=0$ ($\mathbf E\parallel z$), TM $\equiv\delta=90^\circ$.

%=====================================================================
\section{Grundgleichungen im Fourierraum}\label{sec:equations}
%=====================================================================

\subsection{Komponentenform}
Aus \eqref{eq:maxwell_norm} (Buch (V.14)--(V.19) mit $\omega=1$):
\begin{align}
  \partial_yE_z-\partial_zE_y&=iB_x, & \partial_yH_z-\partial_zH_y&=-iD_x,\label{eq:mx1}\\
  \partial_zE_x-\partial_xE_z&=iB_y, & \partial_zH_x-\partial_xH_z&=-iD_y,\label{eq:mx2}\\
  \partial_xE_y-\partial_yE_x&=iB_z, & \partial_xH_y-\partial_yH_x&=-iD_z,\label{eq:mx3}
\end{align}
mit $\mathbf D=\overline{\overline\varepsilon}_r\mathbf E$, $\mathbf B=\overline{\overline\mu}_r\mathbf H$ (hier stets
$\mathbf H\equiv\tilde{\mathbf H}$).

\subsection{Pseudoperiodische Entwicklung}
Wegen der Pseudoperiodizität (Buch I.2.3) gilt für jede Feldkomponente
\begin{equation}
  u(x,y,z)=\sum_{(n,m)\in\mathcal N}u_{nm}(y)\,e^{i(\alpha_nx+\gamma_mz)},
  \label{eq:pseudo}
\end{equation}
mit der Abschneidemenge $\mathcal N$ (Abschnitt~\ref{sec:fourier}). Mit $\alpha=\operatorname{diag}(\alpha_n)$,
$\gamma=\operatorname{diag}(\gamma_m)$ (im Einzelindexraum, Abschnitt~\ref{sec:indexmap}) wird
$\partial_x\to i\alpha$, $\partial_z\to i\gamma$. Produkte $\varepsilon E$ werden durch die
Faktorisierung aus Abschnitt~\ref{sec:fff} ersetzt:
\begin{equation}
  [\mathbf D]=Q_\varepsilon[\mathbf E],\qquad [\mathbf B]=Q_\mu[\mathbf H],
  \label{eq:Qdef}
\end{equation}
mit $3\times3$-Blockmatrizen $Q_\varepsilon=(Q_{\varepsilon,ij})_{i,j\in\{x,y,z\}}$, jeder Block
$n_f\times n_f$ ($n_f=|\mathcal N|$). Ohne magnetische Anisotropie ist $Q_\mu=\mathbb 1$.

\subsection{Elimination von $E_y$, $H_y$}
Aus \eqref{eq:mx2} (rechte Spalte): $i\gamma[H_x]-i\alpha[H_z]=-i[D_y]$, also
$[D_y]=\alpha[H_z]-\gamma[H_x]$, und mit \eqref{eq:Qdef}
\begin{equation}
  [E_y]=Q_{\varepsilon,yy}^{-1}\big(\alpha[H_z]-\gamma[H_x]-Q_{\varepsilon,yx}[E_x]-Q_{\varepsilon,yz}[E_z]\big),
  \label{eq:Ey}
\end{equation}
und analog aus \eqref{eq:mx2} (links): $i[B_y]=i\gamma[E_x]-i\alpha[E_z]$:
\begin{equation}
  [H_y]=Q_{\mu,yy}^{-1}\big(\gamma[E_x]-\alpha[E_z]-Q_{\mu,yx}[H_x]-Q_{\mu,yz}[H_z]\big).
  \label{eq:Hy}
\end{equation}

\subsection{Propagationsgleichung}
Mit dem Zustandsvektor $F=([E_x],[E_z],[H_x],[H_z])^T\in\mathbb C^{4n_f}$ gilt
\begin{equation}
  \frac{dF}{dy}=M(y)\,F .
  \label{eq:ode}
\end{equation}
Herleitung der Zeilen: $\partial_yE_x=i\alpha E_y-iB_z$ \eqref{eq:mx3}, $\partial_yE_z=i\gamma E_y+iB_x$
\eqref{eq:mx1}, $\partial_yH_x=i\alpha H_y+iD_z$ \eqref{eq:mx3}, $\partial_yH_z=i\gamma H_y-iD_x$
\eqref{eq:mx1}; Einsetzen von \eqref{eq:Ey}, \eqref{eq:Hy} und
$B_x=Q_{\mu,xx}H_x+Q_{\mu,xy}H_y+Q_{\mu,xz}H_z$ usw. ergibt
(mit $\mathcal E\equiv Q_{\varepsilon,yy}^{-1}$, $\mathcal U\equiv Q_{\mu,yy}^{-1}$, Buch (V.22') mit $\omega=1$):
\begin{equation}
\begin{aligned}
  M_{11}&=-i\alpha\mathcal EQ_{\varepsilon,yx}-iQ_{\mu,zy}\mathcal U\gamma, &
  M_{12}&=-i\alpha\mathcal EQ_{\varepsilon,yz}+iQ_{\mu,zy}\mathcal U\alpha,\\
  M_{13}&=iQ_{\mu,zy}\mathcal UQ_{\mu,yx}-i\alpha\mathcal E\gamma-iQ_{\mu,zx}, &
  M_{14}&=iQ_{\mu,zy}\mathcal UQ_{\mu,yz}+i\alpha\mathcal E\alpha-iQ_{\mu,zz},\\
  M_{21}&=-i\gamma\mathcal EQ_{\varepsilon,yx}+iQ_{\mu,xy}\mathcal U\gamma, &
  M_{22}&=-i\gamma\mathcal EQ_{\varepsilon,yz}-iQ_{\mu,xy}\mathcal U\alpha,\\
  M_{23}&=-iQ_{\mu,xy}\mathcal UQ_{\mu,yx}-i\gamma\mathcal E\gamma+iQ_{\mu,xx}, &
  M_{24}&=-iQ_{\mu,xy}\mathcal UQ_{\mu,yz}+i\gamma\mathcal E\alpha+iQ_{\mu,xz},\\
  M_{31}&=-iQ_{\varepsilon,zy}\mathcal EQ_{\varepsilon,yx}+i\alpha\mathcal U\gamma+iQ_{\varepsilon,zx}, &
  M_{32}&=-iQ_{\varepsilon,zy}\mathcal EQ_{\varepsilon,yz}-i\alpha\mathcal U\alpha+iQ_{\varepsilon,zz},\\
  M_{33}&=-i\alpha\mathcal UQ_{\mu,yx}-iQ_{\varepsilon,zy}\mathcal E\gamma, &
  M_{34}&=-i\alpha\mathcal UQ_{\mu,yz}+iQ_{\varepsilon,zy}\mathcal E\alpha,\\
  M_{41}&=iQ_{\varepsilon,xy}\mathcal EQ_{\varepsilon,yx}+i\gamma\mathcal U\gamma-iQ_{\varepsilon,xx}, &
  M_{42}&=iQ_{\varepsilon,xy}\mathcal EQ_{\varepsilon,yz}-i\gamma\mathcal U\alpha-iQ_{\varepsilon,xz},\\
  M_{43}&=-i\gamma\mathcal UQ_{\mu,yx}+iQ_{\varepsilon,xy}\mathcal E\gamma, &
  M_{44}&=-i\gamma\mathcal UQ_{\mu,yz}-iQ_{\varepsilon,xy}\mathcal E\alpha .
\end{aligned}
\label{eq:Mgeneral}
\end{equation}
Für $Q_\mu=\mathbb 1$ ($\mathcal U=\mathbb 1$, $Q_{\mu,ij}=0$ für $i\neq j$) vereinfacht sich das zu
\begin{equation}
\begin{aligned}
  M_{11}&=-i\alpha\mathcal EQ_{\varepsilon,yx}, & M_{12}&=-i\alpha\mathcal EQ_{\varepsilon,yz},
  & M_{13}&=-i\alpha\mathcal E\gamma, & M_{14}&=i\alpha\mathcal E\alpha-i\mathbb 1,\\
  M_{21}&=-i\gamma\mathcal EQ_{\varepsilon,yx}, & M_{22}&=-i\gamma\mathcal EQ_{\varepsilon,yz},
  & M_{23}&=-i\gamma\mathcal E\gamma+i\mathbb 1, & M_{24}&=i\gamma\mathcal E\alpha,
\end{aligned}
\label{eq:Mmu1a}
\end{equation}
\begin{equation}
\begin{aligned}
  M_{31}&=i\big(Q_{\varepsilon,zx}-Q_{\varepsilon,zy}\mathcal EQ_{\varepsilon,yx}+\alpha\gamma\big), &
  M_{32}&=i\big(Q_{\varepsilon,zz}-Q_{\varepsilon,zy}\mathcal EQ_{\varepsilon,yz}-\alpha^2\big),\\
  M_{33}&=-iQ_{\varepsilon,zy}\mathcal E\gamma, & M_{34}&=iQ_{\varepsilon,zy}\mathcal E\alpha,\\
  M_{41}&=i\big(Q_{\varepsilon,xy}\mathcal EQ_{\varepsilon,yx}-Q_{\varepsilon,xx}+\gamma^2\big), &
  M_{42}&=i\big(Q_{\varepsilon,xy}\mathcal EQ_{\varepsilon,yz}-Q_{\varepsilon,xz}-\gamma\alpha\big),\\
  M_{43}&=iQ_{\varepsilon,xy}\mathcal E\gamma, & M_{44}&=-iQ_{\varepsilon,xy}\mathcal E\alpha .
\end{aligned}
\label{eq:Mmu1b}
\end{equation}
\begin{tcolorbox}[title=Prüfpunkt (Unit-Test \texttt{test\_propagation\_matrix})]
In einem homogenen isotropen Medium ($Q_{\varepsilon,ij}=\varepsilon\delta_{ij}\mathbb 1$) müssen die
analytischen Moden aus Abschnitt~\ref{sec:homog} Eigenvektoren von $M$ zu den Eigenwerten
$\lambda=\pm i\beta_{nm}$ sein. Achtung: Buch (VIII.3) lässt in $M_{32}$ den Term $-i\alpha^2$ weg
(falsch); \eqref{eq:Mmu1b} ist korrekt.
\end{tcolorbox}

\subsection{Spezialfälle}
\paragraph{1D, in-plane ($\gamma=0$), isotrop, $N_z=0$:} Dann $Q_{\varepsilon,xz}=Q_{\varepsilon,zx}=Q_{\varepsilon,yz}=
Q_{\varepsilon,zy}=0$, $Q_{\varepsilon,zz}=\tp{\varepsilon}$, und \eqref{eq:ode} zerfällt in
\begin{align}
  \text{TE: }&\frac{d}{dy}\begin{pmatrix}[E_z]\\{}[H_x]\end{pmatrix}=
   \begin{pmatrix}0&i\mathbb 1\\ i(\tp{\varepsilon}-\alpha^2)&0\end{pmatrix}\begin{pmatrix}[E_z]\\{}[H_x]\end{pmatrix},
   \label{eq:TE}\\
  \text{TM: }&\frac{d}{dy}\begin{pmatrix}[E_x]\\{}[H_z]\end{pmatrix}=
   \begin{pmatrix}-i\alpha\mathcal EQ_{\varepsilon,yx}& i\alpha\mathcal E\alpha-i\mathbb 1\\
   i(Q_{\varepsilon,xy}\mathcal EQ_{\varepsilon,yx}-Q_{\varepsilon,xx})&-iQ_{\varepsilon,xy}\mathcal E\alpha\end{pmatrix}
   \begin{pmatrix}[E_x]\\{}[H_z]\end{pmatrix}.\label{eq:TM}
\end{align}
Mit $Q_{\varepsilon,xx}=A+\Li$, $Q_{\varepsilon,yy}=\tp{\varepsilon}-A$, $Q_{\varepsilon,xy}=Q_{\varepsilon,yx}=B$
($\Delta=\tp\varepsilon-\Li$, $A=\Delta\tp{c^2}$, $B=\Delta\tp{cs}$, $c=\cos\theta$, $s=\sin\theta$,
$\tan\theta=g'$) ist \eqref{eq:TM} identisch mit Buch (IV.21), (IV.22). Diese beiden
$2n_f$-Systeme sind der erste Implementierungsschritt (Meilenstein M2/M3).
\paragraph{1D konisch:} volles $4n_f$-System mit $\gamma=\gamma_0\mathbb 1$.
\paragraph{2D:} volles $4n_f$-System mit $n_f=(2N_x+1)(2N_z+1)$.

%=====================================================================
\section{Fourier-Darstellung und Toeplitz-Matrizen}\label{sec:fourier}
%=====================================================================

\subsection{Abschneidemenge und Einzelindex}\label{sec:indexmap}
1D: $\mathcal N=\{n:\ -N_x\le n\le N_x\}$, $n_f=2N_x+1$, Index $r=n+N_x$ (0-basiert).
2D (rechteckige Abschneidung): $\mathcal N=\{(n,m):|n|\le N_x,|m|\le N_z\}$,
\begin{equation}
  r=(m+N_z)+(2N_z+1)(n+N_x),\qquad
  n=-N_x+\Big\lfloor\frac{r}{2N_z+1}\Big\rfloor,\quad m=r-(2N_z+1)(n+N_x)-N_z
  \label{eq:indexmap}
\end{equation}
(0-basiert; das Buch (IX.14) dividiert fälschlich durch $N+1$). Optional später: kreisförmige
Abschneidung $(n/N_x)^2+(m/N_z)^2\le1$; dann ist die Abbildung eine Tabelle. Der Code darf nur
über die Klasse \verb|IndexMap| auf $(n,m)\leftrightarrow r$ zugreifen.

\subsection{Fourierkoeffizienten}
Für eine $d_x$-periodische (bzw.\ $(d_x,d_z)$-periodische) Funktion $f$ gilt
\begin{equation}
  f(x,z)=\sum_{p,q}f_{pq}\,e^{i(pK_xx+qK_zz)},\qquad
  f_{pq}=\frac{1}{d_xd_z}\iint_{\text{Zelle}}f(x,z)\,e^{-i(pK_xx+qK_zz)}\,dx\,dz .
  \label{eq:fcoef}
\end{equation}
Für Toeplitz-Matrizen werden $|p|\le2N_x$, $|q|\le2N_z$ benötigt.

\paragraph{Intervall (1D).} $f=\chi_{[x_1,x_2]}$ (Periode $d$, $K=2\pi/d$):
\begin{equation}
  f_0=\frac{x_2-x_1}{d},\qquad f_p=\frac{e^{-ipKx_2}-e^{-ipKx_1}}{-ipKd}\quad(p\ne0).
  \label{eq:interval}
\end{equation}
Stückweise konstante Funktionen $=\sum_j\varepsilon_j\chi_{I_j}$ (Linearität).

\paragraph{Rechteck (2D).} Produkt zweier Intervallformeln.

\paragraph{Polygon (2D, exakt).} Für ein einfaches Polygon $P$ mit Ecken $\mathbf v_1,\dots,\mathbf v_L$
(gegen den Uhrzeigersinn), Kanten $\mathbf t_j=\mathbf v_{j+1}-\mathbf v_j$, Länge $\ell_j$, Mittelpunkt
$\mathbf c_j$, äußerer Normale $\mathbf n_j$ und $\mathbf k=(pK_x,qK_z)\neq0$:
\begin{equation}
  \iint_Pe^{-i\mathbf k\cdot\mathbf r}d^2r=\frac{i}{|\mathbf k|^2}\sum_{j=1}^{L}(\mathbf k\cdot\mathbf n_j)\,\ell_j\,
  e^{-i\mathbf k\cdot\mathbf c_j}\,\operatorname{sinc}\Big(\frac{\mathbf k\cdot\mathbf t_j}{2}\Big),\qquad
  \operatorname{sinc}u=\frac{\sin u}{u},
  \label{eq:polygon}
\end{equation}
für $\mathbf k=0$ die Fläche. Herleitung: Gaußscher Satz mit
\[
  e^{-i\mathbf k\cdot\mathbf r}=\frac{i}{|\mathbf k|^2}\nabla\cdot\big(\mathbf k\,e^{-i\mathbf k\cdot\mathbf r}\big),\qquad
  \int_0^1e^{-i\mathbf k\cdot(\mathbf v_j+s\mathbf t_j)}\ell_j\,ds=\ell_j\,e^{-i\mathbf k\cdot\mathbf c_j}\operatorname{sinc}(\mathbf k\cdot\mathbf t_j/2).
\]
Kreise/Ellipsen: analytisch über Bessel-Funktion ($J_1$) oder als feines Polygon.

\paragraph{Allgemeine Funktionen.} Abtastung auf einem feinen Gitter ($\ge 8\times$ die höchste
benötigte Ordnung) und FFT. \textbf{Nur für stetige Funktionen} (z.\,B.\ Normalenvektorkomponenten);
für unstetige Funktionen (Permittivität) möglichst analytische Formeln, sonst Aliasing-Fehler
$O(1/N_{\text{grid}})$.

\subsection{Toeplitz-Matrizen}
\begin{equation}
  \tp{f}_{rr'}=f_{n-n',\,m-m'},\qquad r\leftrightarrow(n,m),\ r'\leftrightarrow(n',m').
  \label{eq:toeplitz}
\end{equation}
1D: gewöhnliche Toeplitz-Matrix; 2D: block-Toeplitz mit Toeplitz-Blöcken. Rechenregeln
(Buch Anh.~II; im abgeschnittenen Raum nur näherungsweise):
$\tp{1}=\mathbb 1$, $\tp{f+g}=\tp f+\tp g$, $\tp{fg}\approx\tp f\tp g$ falls keine gleichzeitigen Sprünge.
\begin{tcolorbox}[title=Numerischer Befund (Kapitelprüfung)]
$\tp f\tp g-\tp{fg}$ ist im Inneren der Matrix auf Maschinengenauigkeit null, an den Rändern
(höchste Ordnungen) bleibt ein Fehler $O(1)$, der mit $N$ \emph{nicht} abnimmt. Deshalb nie
$\tp f\tp g$ durch $\tp{fg}$ ersetzen, wenn eine der Funktionen unstetig ist, außer die Spezifikation
verlangt es ausdrücklich.
\end{tcolorbox}

%=====================================================================
\section{Faktorisierung (FFF) und die Matrizen $Q_\varepsilon$}\label{sec:fff}
%=====================================================================

\subsection{Regeln von L.~Li}
Für stückweise glatte, beschränkte, periodische $f,g$ und $h=fg$:
\begin{enumerate}[nosep]
  \item keine gleichzeitigen Sprünge: $[h]=\tp f[g]$ (Laurent);
  \item nur komplementäre Sprünge ($h$ stetig): $[h]=\tp{1/f}^{-1}[g]$ (invers);
  \item gleichzeitige, nicht komplementäre Sprünge: keine der beiden Regeln.
\end{enumerate}
Die FFF zerlegt $\mathbf E$ in stetige Größen, sodass nur Fall 1 und 2 auftreten.

\subsection{Normalen- und Tangentenfelder}
$\mathbf N(x,z)$ ist ein reelles Einheitsvektorfeld, das auf jeder Materialgrenze der Schicht mit der
Flächennormalen übereinstimmt und \emph{an den Unstetigkeitsstellen von $\varepsilon$ stetig} ist
(anderswo darf es beliebig sein). Es hängt nicht von $y$ ab, wenn die Schicht durch ein Profil
$y=g(x,z)$ beschrieben wird (Normale in $(x,y,z)$ parallel zur Normale im Profilpunkt darunter).
\begin{align}
  \text{Profil }y=g(x)\ \text{(1D)}:\quad &\mathbf N=\frac{(-g',\,1,\,0)}{\sqrt{1+g'^2}},\label{eq:N1d}\\
  \text{Profil }y=g(x,z)\ \text{(2D)}:\quad &\mathbf N=\frac{(-\partial_xg,\,1,\,-\partial_zg)}{\sqrt{1+(\partial_xg)^2+(\partial_zg)^2}},\label{eq:N2d}\\
  \text{lamellar 1D}:\quad &N_x^2=1,\ N_y=N_z=0\ \text{(nur $\mathbf N\mathbf N^T$ wird gebraucht)},\label{eq:Nlam}\\
  \text{$y$-invarianter Querschnitt (2D)}:\quad &\mathbf N=(N_x,0,N_z)\ \text{aus Algorithmus \ref{alg:nfield}}.\notag
\end{align}
\begin{tcolorbox}[title=Algorithmus NF: Normalenfeld für $y$-invariante 2D-Querschnitte (v1.5)\label{alg:nfield}]
\begin{enumerate}[nosep]
  \item Raster der Einheitszelle mit $G_x\times G_z$ Punkten $x_i=d_xi/G_x$, $z_k=d_zk/G_z$,
        $G_{x,z}=\max(G,8(2N_{x,z}+1))$, Standard $G=512$.
  \item \textbf{Materialgrenzen} exakt aus den Polygonen: jede Kante wird in 64 Stücke geteilt; ein Stück ist
        Grenze, wenn die Materialien beidseits (periodische Punktsuche, Abstand $10^{-9}d$) verschieden sind.
        Damit entfallen Kanten auf dem Zellrand, die an ein periodisches Bild desselben Materials stoßen,
        innere Kanten aneinanderstoßender Einschlüsse gleichen Materials und Einschlüsse aus Hintergrundmaterial.
        Für jeden Rasterpunkt $\mathbf r$: nächster Grenzpunkt $\mathbf p$ (Grenzstücke mit Bildern $\pm d_x,\pm d_z$),
        Richtung $\mathbf n=(\mathbf r-\mathbf p)/|\mathbf r-\mathbf p|$ (auf der Grenze: Kantennormale).
        Gleich weit entfernte Kandidaten (Mittelachse, Toleranz $10^{-12}d$) werden gemittelt (Symmetrie).
  \item \textbf{Glättung} mit periodischem Gaußfilter gleicher \emph{physikalischer} Breite in $x$ und $z$
        ($\sigma$ Rasterpunkte entlang der kürzeren Periode, Standard 3).
  \item \textbf{Tensorvariante (Standard):} geglättet wird das Projektorfeld $\mathbf n\mathbf n^T$; das Ergebnis
        ist direkt $\tp{\cdot}$-Quelle für $N_iN_j$ (symmetrisch, Spur 1, stetig, an geraden Grenzen $=\mathbf n\mathbf n^T$).
        \textbf{Vektorvariante} (Fassung v1.4): geglättet wird $\nabla\phi$ der vorzeichenbehafteten Abstandsfunktion
        ($\phi>0$ im Hintergrund), $\mathbf N=\mathbf g/|\mathbf g|$ wo $|\mathbf g|>\tau\max|\mathbf g|$ ($\tau=10^{-3}$),
        sonst $N_xN_x=N_zN_z=\tfrac12$, $N_xN_z=0$.
  \item Fourierkoeffizienten von $N_x^2$, $N_xN_z$, $N_z^2$ per DFT des Rasters (glatte Funktionen).
\end{enumerate}
Literatur: Schuster et al.\ (JOSA A 24, 2880, 2007); Götz et al.\ (Opt.\ Express 16, 17295, 2008).
\end{tcolorbox}
\textbf{Begründung der Tensorvariante (Befund M7).} Für die FFF-Herleitung \eqref{eq:Qiso} genügt
\emph{irgendeine} stetige Zerlegung $\mathbf E=(\mathbb 1-P)\mathbf E+P\mathbf E$, deren Matrixfeld $P$ an den
Unstetigkeitsstellen von $\varepsilon$ gleich $\mathbf n\mathbf n^T$ ist; $P$ muss weder Projektor noch aus einem
Vektorfeld gebildet sein. An Mittelachsen zwischen gegenüberliegenden Wänden (Lamellenmitte, Steg zwischen zwei
Löchern) kehrt $\nabla\phi$ das Vorzeichen um; die Vektorglättung läuft dort durch null und fällt auf $\tfrac12$
zurück. Liegt die Mittelachse auf Rasterpunkten, weicht $N_x^2$ einer $z$-invarianten Lamelle in den Koeffizienten um
$\approx2\cdot10^{-3}$ von 1 ab (Unit-Test \texttt{VectorMethodFallbackOnMedialAxis}); V9 wäre nicht exakt. Die
Tensorvariante ist vorzeichenfrei, singularitätsfrei und reproduziert Lamellen exakt ($10^{-14}$). Sie ist
auch für die Projektorform \eqref{eq:QaniB} (Variante B, M8) direkt verwendbar. An konvexen Ecken liefert sie
gemischte (Rang-2-)Tensoren statt gedrehter Normalen; beide Varianten weichen dort von $\mathbf n\mathbf n^T$ ab.
Ein stetiges \emph{vorzeichenbehaftetes} Feld existiert für Löcher nicht; \texttt{FFF\_rig} ist für
\texttt{Polygon2D} daher nicht zulässig.
Tangentenvektoren für Variante~A (unten) mit $\mathbf T_2\times\mathbf T_1=\mathbf N$:
\begin{align}
  \text{Profilschichten }(N_y>0):\quad&\mathbf T_1=\frac{(N_y,-N_x,0)}{\sqrt{N_x^2+N_y^2}},\quad
  \mathbf T_2=\mathbf T_1\times\mathbf N=\frac{(-N_xN_z,\,-N_yN_z,\,N_x^2+N_y^2)}{\sqrt{N_x^2+N_y^2}},\label{eq:T_profile}\\
  \text{$y$-invariante Schichten }(N_y=0):\quad&\mathbf T_1=\hat{\mathbf y},\quad \mathbf T_2=\hat{\mathbf y}\times\mathbf N=(N_z,0,-N_x).\label{eq:T_yinv}
\end{align}
(Prüfung: $(N_z,0,-N_x)\times(0,1,0)=(N_x,0,N_z)=\mathbf N$. Die Buchwahl \eqref{eq:T_profile} ist für
$N_x=N_y=0$ singulär und daher für $y$-invariante Schichten ungeeignet.)

\subsection{Isotrope Materialien}
Mit $\Delta=\tp{\varepsilon}-\Li$ (Buch (IV.32), (IX.11)):
\begin{equation}
  Q_{\varepsilon}=\tp{\varepsilon}\otimes\mathbb 1_3-\Delta\,\tp{\mathbf N\mathbf N^T},\qquad
  Q_{\varepsilon,ij}=\delta_{ij}\tp{\varepsilon}-\Delta\tp{N_iN_j}.
  \label{eq:Qiso}
\end{equation}
Herleitung: $\mathbf D=\varepsilon\mathbf E_T+\varepsilon\mathbf E_N$; $\varepsilon\mathbf E_T$: Laurent
($\mathbf E_T$ stetig), $\varepsilon\mathbf E_N=\mathbf D_N$ stetig: invers; mit $\mathbf E_N=\mathbf N\mathbf N^T\mathbf E$
und $\mathbf E_T=\mathbf E-\mathbf E_N$ folgt \eqref{eq:Qiso} (Laurent für $\mathbf N\mathbf N^T\mathbf E$, da $\mathbf N$ stetig).
Grenzfälle (Unit-Tests): $N\to\pm\hat x$ überall $\Rightarrow$ Li-Formulierung für Lamellen
($Q_{xx}=\Li$, $Q_{yy}=Q_{zz}=\tp\varepsilon$); $\mathbf N=\hat y$ überall $\Rightarrow$ klassische Laurent-Formulierung
($Q_{yy}=\Li$, sonst $\tp\varepsilon$, vgl.\ Buch IV.2.3: $a/d\to0$).

\subsection{Varianten der isotropen FFF (Befund aus M3)}\label{sec:fffvariants}
Die Buchformel \eqref{eq:Qiso} entsteht aus einer exakteren Form durch zwei Näherungen im
abgeschnittenen Raum. Aus $D_i=\varepsilon E_{T,i}+N_iD_N$ mit $E_{T,i}=E_i-N_iE_N$ (stetig), $D_N$ (stetig),
$E_N=\sum_jN_jE_j$ folgt mit den Li-Regeln (Laurent für $\varepsilon\cdot$stetig und $N\cdot$beliebig,
invers für $D_N=\varepsilon E_N$):
\begin{equation}
  Q^{(\mathrm{rig})}_{ij}=\delta_{ij}\tp{\varepsilon}-\tp{\varepsilon}\tp{N_i}\tp{N_j}+\tp{N_i}\,\Li\,\tp{N_j}.
  \label{eq:Qrig}
\end{equation}
Das Buch ersetzt $\tp{N_i}\tp{N_j}\to\tp{N_iN_j}$ und vertauscht $\tp{N_i}$ mit $\Li$; das ergibt
\eqref{eq:Qiso}. Eine dritte Variante symmetrisiert das Produkt:
\begin{equation}
  Q^{(\mathrm{sym})}_{ij}=\delta_{ij}\tp{\varepsilon}-\tfrac12\big(\Delta\tp{N_iN_j}+\tp{N_iN_j}\Delta\big).
  \label{eq:Qsym}
\end{equation}
\textbf{Numerischer Befund} (RK, TM, $\lambda=0{,}6328\,\mu$m; Referenzen aus der C-Methode):
\begin{center}
\small
\begin{tabular}{lcc}
\toprule
& Gold-Sinus $d=1$, $a=0{,}25\,\mu$m, $20^\circ$, $e_{-1}$ ($N=8/16/24$) & Al-Sinus V4, $e_{-1}$ ($N=20$) \\
\midrule
Referenz & $0{,}26614$ & $0{,}8696$ \\
\verb|FFF_iso| \eqref{eq:Qiso} & falsch ($0{,}327$ / $0{,}217$; $\sum e>1$) & $0{,}86951$ \\
\verb|FFF_sym| \eqref{eq:Qsym} & $0{,}2681$ / $0{,}2672$ / $0{,}2665$ & $0{,}86565$ \\
\verb|FFF_rig| \eqref{eq:Qrig} & $0{,}26609$ / $0{,}26621$ / -- & $0{,}87075$ \\
\bottomrule
\end{tabular}
\end{center}
\begin{itemize}[nosep]
  \item Die Buchvariante versagt bei schwach absorbierenden Metallen ($\operatorname{Re}\varepsilon<0$, kleines
        $\operatorname{Im}\varepsilon$). Vermuteter Mechanismus: $1/\varepsilon$ läuft fast durch null, $\Li$ wird groß,
        und damit auch der Fehler der Vertauschung von $\tp{N_i}$ und $\Li$. Der Löser gibt dann eine Warnung aus.
  \item In einer $y$-invarianten Scheibe aus isotropem Material müssen die Eigenwerte als $\pm\lambda$-Paare
        auftreten. Laurent und \verb|FFF_sym| erhalten diese Paarung exakt, \verb|FFF_iso| und \verb|FFF_rig|
        nicht. Da \verb|FFF_rig| trotzdem am genauesten ist, ist die Paarung \emph{kein} Gütekriterium.
  \item \verb|FFF_rig| braucht das \emph{vorzeichenbehaftete} Normalenfeld; springt es (Ecken von Dreiecks-
        und Echelette-Profilen), versagt die Variante ($\sum e>1$). Dort ist \verb|FFF_sym| zu verwenden.
  \item Keine der FFF-Varianten erhält bei verlustfreien Medien die Energie exakt (Abweichung
        $10^{-5}$--$5\cdot10^{-4}$), Laurent schon. Tests prüfen die Energie bei Profilen daher mit Toleranz.
\end{itemize}
\textbf{2D (M7):} Für \texttt{Polygon2D} (Querschnitte mit Löchern/Säulen) ist \verb|FFF_sym| Standard;
\verb|FFF_iso| liefert beim Ag-Lochgitter (V10, $\varepsilon=-79{,}9+1{,}79i$) bei $N=3$ und $N=6$ $R>1$
(dieselbe Ursache wie bei Gold in 1D), \verb|FFF_rig| ist mangels stetigem Vorzeichenfeld nicht anwendbar.
\texttt{Profile2D} verwendet wie in 1D \verb|FFF_rig| (analytische Formen) bzw.\ \verb|FFF_sym| (Funktionen);
$\mathbf N$ aus \eqref{eq:N2d} ist dort stetig und vorzeichenbehaftet.

\textbf{Standard (\texttt{Factorization::Auto})} für isotrope Profilschichten: \verb|FFF_rig| bei glatten
Profilen (Sinus), \verb|FFF_sym| bei Profilen mit Ecken (stückweise linear) und bei benutzerdefinierten
Funktionen. \verb|FFF_iso| bleibt als Buchvariante wählbar. Beide neuen Varianten sind Eigenherleitungen
dieses Projekts; ob sie so in der Literatur stehen, ist noch zu prüfen. Offener Punkt: Für Gold-Dreiecke
gibt es noch keine belastbare Referenz (V14).

\subsection{Anisotrope Materialien, Variante A (Buch Kap.~V)}
Stetiger Vektor $\mathbf F_\varepsilon=(E_{T_1},D_N,E_{T_2})=A_\varepsilon\mathbf E$ mit
\begin{equation}
  A_\varepsilon=\begin{pmatrix}\mathbf T_1^T\\ \mathbf v^T\\ \mathbf T_2^T\end{pmatrix},\quad
  \mathbf v^T=\mathbf N^T\overline{\overline\varepsilon}\ \ \big(v_j=\textstyle\sum_iN_i\varepsilon_{ij}\big),\qquad
  \xi=\mathbf N^T\overline{\overline\varepsilon}\mathbf N,
  \label{eq:Aeps}
\end{equation}
\begin{equation}
  C_\varepsilon=A_\varepsilon^{-1}=\frac{1}{\xi}\Big(\ \mathbf v\times\mathbf T_2\ \Big|\ \mathbf N\ \Big|\ -\mathbf v\times\mathbf T_1\ \Big)
  \quad\text{(Spalten)},\qquad \det A_\varepsilon=\xi .
  \label{eq:Ceps}
\end{equation}
(Buchtext S.\,102 nennt die Spaltenreihenfolge falsch; (V.11) selbst ist richtig.)
$\mathbf D=\overline{\overline\varepsilon}C_\varepsilon\mathbf F_\varepsilon$: Laurent; $\mathbf F_\varepsilon=A_\varepsilon\mathbf E$ stetig: invers.
\begin{equation}
  Q_\varepsilon=\tp{\overline{\overline\varepsilon}C_\varepsilon}\,\tp{C_\varepsilon}^{-1}
  \label{eq:QaniA}
\end{equation}
($3\times3$ Blöcke; $\tp{X}$ für eine matrixwertige Funktion $X$ ist die Blockmatrix der Toeplitz-Matrizen
der Einträge). Speziell 1D-Profil ($\mathbf T_2=\hat z$, $\mathbf T_1=(N_y,-N_x,0)$):
\begin{equation}
  C_\varepsilon=\frac1\xi\begin{pmatrix}v_y&N_x&-v_zN_x\\-v_x&N_y&-v_zN_y\\0&0&\xi\end{pmatrix},\qquad
  \xi=v_xN_x+v_yN_y
  \label{eq:C1d}
\end{equation}
(Buch (VIII.2) hat in $C_{23}$ ein falsches Vorzeichen). \textbf{Voraussetzung:} $\xi\ne0$ überall
(bei indefiniten, \enquote{hyperbolischen} Tensoren nicht garantiert; dann Warnung und Variante B
oder Laurent). \textbf{Achtung:} $C_\varepsilon$ hängt über $\overline{\overline\varepsilon}(x,y,z)$ von $y$ ab, außer bei
$y$-invarianten Schichten -- $\tp{C_\varepsilon}$ muss für jede Höhe neu berechnet werden (Buch S.\,108
behauptet Gegenteiliges).

\paragraph{Lamellen mit Variante A.} Für 1D-Lamellen genügt das \emph{konstante} Feld $\mathbf N=\hat x$:
An jeder Wand ist es bis auf das Vorzeichen die Normale, und das Vorzeichen ist überall gleich, also ist
$\mathbf F_\varepsilon$ stetig. Mit isotropem Tensor reproduziert Variante A dann exakt die Li-Lamellenformel
(V8, $10^{-10}$). Standard für anisotrope Lamellen und Profile ist \verb|FFF_A|.
\paragraph{Fourierkoeffizienten.} $C_\varepsilon(x)$ und $\overline{\overline\varepsilon}C_\varepsilon(x)$ sind stückweise glatt
(Tensor stückweise konstant, $\mathbf N$ glatt). Koeffizienten per zusammengesetzter 8-Punkt-Gauß-Quadratur je
Intervall mit $|pK|h\le1$ (\verb|piecewiseSmoothCoefficients|); für konstante Funktionen exakt bis auf Rundung.

\subsection{Anisotrope Materialien, Variante B (vorzeichenunabhängig)}
Variante A verlangt ein stetiges $\mathbf N$ \emph{mit Vorzeichen}. Bei 2D-Querschnitten (Löcher, Säulen)
gibt es kein solches Feld; dort würden $E_{T_1}$, $D_N$ im Inneren eines homogenen Gebiets springen. Variante B
verwendet nur die Projektoren $P_N=\mathbf N\mathbf N^T$, $P_T=\mathbb 1-P_N$ und die stetigen
\emph{Vektoren} $\mathbf e_t=P_T\mathbf E$, $\mathbf d_n=P_N\mathbf D$:
aus $\mathbf N^T\mathbf D=\mathbf N^T\overline{\overline\varepsilon}\mathbf e_t+\xi\,\mathbf N^T\mathbf E$ folgt
\begin{equation}
  \mathbf E=\Big(\mathbb 1-\tfrac1\xi P_N\overline{\overline\varepsilon}\Big)\mathbf e_t+\tfrac1\xi\mathbf d_n,\qquad
  \mathbf D=\overline{\overline\varepsilon}\Big(\mathbb 1-\tfrac1\xi P_N\overline{\overline\varepsilon}\Big)\mathbf e_t+\tfrac1\xi\overline{\overline\varepsilon}\,\mathbf d_n .
  \label{eq:projector}
\end{equation}
Alle Koeffizientenfunktionen sind invariant unter $\mathbf N\to-\mathbf N$. Laurent auf beide
Zeilen, $[\mathbf e_t]=\tp{P_T}[\mathbf E]$ (Laurent), Elimination von $[\mathbf d_n]$:
\begin{equation}
  Q_\varepsilon^{(B)}=\tp{G_1}\tp{P_T}+\tp{G_2}\,\tp{\xi^{-1}}^{-1}\big(\mathbb 1-\tp{H_1}\tp{P_T}\big),
  \label{eq:QaniB}
\end{equation}
\[
  G_1=\overline{\overline\varepsilon}\big(\mathbb 1-\xi^{-1}P_N\overline{\overline\varepsilon}\big),\quad
  G_2=\xi^{-1}\overline{\overline\varepsilon},\quad H_1=\mathbb 1-\xi^{-1}P_N\overline{\overline\varepsilon},
\]
mit $\tp{\xi^{-1}}$ = blockdiagonal ($3\times$ dieselbe skalare Toeplitz-Matrix). Für isotropes $\varepsilon$
reduziert sich \eqref{eq:QaniB} bis auf Abschneideeffekte ($\tp{P_T}\tp{P_T}\approx\tp{P_T}$) auf
\eqref{eq:Qiso}. \textbf{Test:} Variante A und B müssen auf Profilschichten mit $N$ gegen dieselben
Effizienzen konvergieren; Variante B mit isotropem Tensor gegen \eqref{eq:Qiso}.
\emph{Hinweis:} Variante B ist eine Eigenherleitung dieses Plans (nicht im Buch). M5-Befund: am Kobalt-Gitter
(V7) stimmt sie mit Variante A auf 1--3\,\% (TM-Anteile) bzw.\ 2--15\,\% (sehr kleine TE-Anteile) überein; ein
Literaturvergleich (Li, J.\ Mod.\ Opt.\ 45, 1313 (1998)) steht noch aus. Sie ist für 2D-Querschnitte (M8) gedacht.

\paragraph{Variante B für \texttt{Polygon2D} (M8, v1.6).} $P_N$ muss an den Grenzflächen ein Projektor
$\mathbf N\mathbf N^T$ sein und $\xi=\mathbf N^T\overline{\overline\varepsilon}\mathbf N$ bleibt nur für Rang-1-Felder
sinnvoll; das Tensorfeld von Algorithmus NF (Rang 2 an Ecken) ist hier ungeeignet (für Metalle könnte
$\mathbb 1-P+P\overline{\overline\varepsilon}$ singulär werden). Verwendet wird das \emph{Rang-1-Feld}
$\mathbf N=\mathbf g/|\mathbf g|$ der geglätteten vorzeichenbehafteten Abstandsfunktion (Vektorvariante von
Algorithmus NF; Vorzeichenwechsel an Mittelachsen sind für B bedeutungslos), wo $|\mathbf g|\le\tau\max|\mathbf g|$
die Hauptachse des geglätteten Tensors (z.\,B.\ Lamellenmitte: exakt $\hat x$). Die Koeffizientenfunktionen
$F\in\{G_1,G_2,H_1,P_T,\xi^{-1}\}$ springen mit dem Material:
\[
  F=F(\overline{\overline\varepsilon}_{\text{bg}},\mathbf N)+\textstyle\sum_k\chi_k\big[F(\overline{\overline\varepsilon}_k,\mathbf N)-F(\overline{\overline\varepsilon}_{\text{bg}},\mathbf N)\big].
\]
Glatte Anteile per DFT des NF-Rasters (Rastergröße Zweierpotenz, FFT), Produkte mit $\chi_k$ per Faltung
$\sum_{p'q'}\hat D_{p'q'}\hat\chi_{p-p',q-q'}$ mit den exakten Polygonkoeffizienten und $P_f=2P$ Ordnungen des
glatten Faktors (Laurent-Regel für glatt $\times$ unstetig). Für konstantes $\mathbf N$ exakt: ein
$z$-invarianter Streifen liefert mit isotropem Tensor die Li-Formel ($10^{-10}$, V8-2D). Voraussetzung:
Einschlüsse liegen im Hintergrund (das Vorzeichen von $\phi$ ist sonst nicht definiert); sich berührende
Einschlüsse verschiedener Materialien nur mit Laurent. Standard: anisotrope \texttt{Polygon2D} \verb|FFF_B|,
anisotrope \texttt{Profile2D} \verb|FFF_A| ($\mathbf N$ aus \eqref{eq:N2d} stetig und vorzeichenbehaftet,
$\mathbf T_1=(N_y,-N_x,0)/\sqrt{N_x^2+N_y^2}$, $\mathbf T_2=\mathbf T_1\times\mathbf N$).

\subsection{Weitere Faktorisierungsoptionen (für Vergleiche)}
\verb|Laurent| ($Q_{ij}=\tp{\varepsilon_{ij}}$), \verb|LiLamellar| (nur 1D-Lamellen, Buch (IV.5)/(IV.6)),
\verb|FFF_iso| \eqref{eq:Qiso}, \verb|FFF_sym| \eqref{eq:Qsym}, \verb|FFF_rig| \eqref{eq:Qrig},
\verb|FFF_A| \eqref{eq:QaniA}, \verb|FFF_B| \eqref{eq:QaniB}.
Die Option wird je Schicht gewählt; Standard: isotrope Lamellen \verb|LiLamellar|, isotrope Profile (1D und 2D) siehe
Abschnitt~\ref{sec:fffvariants}, isotrope \texttt{Polygon2D} \verb|FFF_sym|, anisotrop Profil (1D und 2D) \verb|FFF_A|,
anisotrop $y$-invariant (Lamellen: \verb|FFF_A| mit $\mathbf N=\hat x$; \texttt{Polygon2D}) \verb|FFF_B|.

%=====================================================================
\section{Homogene Medien: Modenbasen}\label{sec:homog}
%=====================================================================
In homogenen Medien (Substrat, Superstrat, homogene Zwischenschichten, Referenzmedium der
dünnen Schichten) ist $M$ \emph{ordnungsweise} entkoppelt: $M$ zerfällt in $n_f$ unabhängige
$4\times4$-Matrizen $M^{(r)}$ (eine je Ordnung $r\leftrightarrow(n,m)$).

\subsection{Isotrop: analytische Basis der tangentialen E-Komponenten}
Für eine ebene Welle $e^{i(\alpha x+k_yy+\gamma z)}$ ($k_y=+\beta$ auf, $-\beta$ ab) gilt in normierten
Einheiten $\mathbf H=\mathbf k\times\mathbf E$ ($\mu=1$), $\mathbf k\cdot\mathbf E=0$. Mit
$E_y=-(\alpha E_x+\gamma E_z)/k_y$ und $\alpha^2+k_y^2+\gamma^2=\varepsilon$:
\begin{equation}
  \begin{pmatrix}H_x\\H_z\end{pmatrix}=\frac{1}{k_y}K\begin{pmatrix}E_x\\E_z\end{pmatrix},\qquad
  K=\begin{pmatrix}\alpha\gamma&\varepsilon-\alpha^2\\-(\varepsilon-\gamma^2)&-\alpha\gamma\end{pmatrix},\qquad
  \det K=\varepsilon\beta^2 .
  \label{eq:Kmat}
\end{equation}
Amplitudenvektor $V=(\mathbf a;\mathbf b)$ mit $\mathbf a=(a_x,a_z)$ (ab), $\mathbf b=(b_x,b_z)$ (auf) = tangentiale
E-Komponenten der Teilwellen am Bezugspunkt. Mit $\mathcal V=K\beta^{-1}$ (ordnungsweise $2\times2$):
\begin{equation}
  F=\Psi V,\qquad \Psi=\begin{pmatrix}\mathbb 1&\mathbb 1\\-\mathcal V&\mathcal V\end{pmatrix},\qquad
  \Psi^{-1}=\frac12\begin{pmatrix}\mathbb 1&-\mathcal V^{-1}\\\mathbb 1&\mathcal V^{-1}\end{pmatrix},\qquad
  \mathcal V^{-1}=\frac{1}{\varepsilon\beta}\begin{pmatrix}-\alpha\gamma&-(\varepsilon-\alpha^2)\\\varepsilon-\gamma^2&\alpha\gamma\end{pmatrix}.
  \label{eq:Psi_iso}
\end{equation}
Die Spalten sind Eigenvektoren von $M^{(r)}$ zu $\lambda=-i\beta$ (ab) und $\lambda=+i\beta$ (auf).
(Diese Basis ist der in Buch Anh.~IV vorzuziehen: dort wird durch $k^2-\gamma_m^2$ dividiert, was im 2D-Fall
verschwinden kann; hier tritt nur $\beta$ im Nenner auf.) Kontrollen: 1D-TE ($\gamma=0$, $E_x=0$):
$H_x=k_yE_z$; 1D-TM: $E_x=-k_yH_z/\varepsilon$.

\subsection{Anisotrop: $4\times4$-Eigenzerlegung je Ordnung (Berreman)}
$M^{(r)}$ aus \eqref{eq:Mgeneral} mit $Q_{\varepsilon,ij}=\varepsilon_{ij}$ (Skalare) und $\alpha_n,\gamma_m$.
Eigenwerte $\lambda_k$, Eigenvektoren $w_k$ ($k=1..4$). Klassifikation (Feld $\propto e^{\lambda y}$):
\begin{equation}
  \text{auf (b)}:\ \operatorname{Re}\lambda<-\tau_\lambda\ \ \text{oder}\ \ \big(|\operatorname{Re}\lambda|\le\tau_\lambda\ \text{und}\ S_y(w)>0\big),
  \qquad\text{sonst ab (a)},
  \label{eq:classify}
\end{equation}
mit dem zeitgemittelten Poynting-Fluss der Mode
\begin{equation}
  S_y(w)=\tfrac12\operatorname{Re}\big(E_zH_x^*-E_xH_z^*\big)
  \label{eq:Sy}
\end{equation}
($E_x,E_z,H_x,H_z$ = Komponenten von $w$). Es müssen genau zwei auf- und zwei abwärts
klassifizierte Moden entstehen, sonst Fehler. Normierung: $\Psi$-Spalten so skalieren, dass
$\max(|E_x|,|E_z|)=1$. Für isotrope Medien muss dieses Verfahren \eqref{eq:Psi_iso} (bis auf
Skalierung/Mischung entarteter Moden) reproduzieren (Unit-Test).

\subsection{Leistungsfluss und Effizienzen}
In einem homogenen Medium sind die Ordnungen bezüglich des Flusses orthogonal; die von der Ordnung
$r$ in $\pm y$ transportierte Leistung (pro Periode, normiert) ist \eqref{eq:Sy} mit den Fourierkoeffizienten
dieser Ordnung. Einfallende Leistung $S_{\text{inc}}=|S_y(\text{einf.\ Welle})|$. Effizienzen:
\begin{equation}
  e^{(r)}_{nm}=\frac{S_y^{\uparrow}(r)}{S_{\text{inc}}},\qquad e^{(t)}_{nm}=\frac{|S_y^{\downarrow}(r)|}{S_{\text{inc}}}
  \quad(\text{nur propagierende Ordnungen in verlustfreien Medien}).
  \label{eq:eff}
\end{equation}
Geschlossene Form für isotrope verlustfreie Medien (Buch (IX.20), korrigiert; normiert,
$k^2\to\varepsilon$, $|\mathbf A|=1$):
\begin{equation}
  e_{nm}=\frac{(\varepsilon-\gamma_m^2)|B_x|^2+(\varepsilon-\alpha_n^2)|B_z|^2+2\alpha_n\gamma_m\operatorname{Re}(B_x\overline{B_z})}
  {\beta_{00}^{(M)}\,\beta_{nm}},
  \label{eq:eff_closed}
\end{equation}
wobei $\varepsilon$ und $\beta_{nm}$ zu dem Medium gehören, in dem die Ordnung läuft (für die
Transmission also $\varepsilon_1$, $\beta^{(1)}_{nm}$; das Buch vergisst $k_2\to k_1$), und
$\beta^{(M)}_{00}=\beta_0$. Bei absorbierendem Substrat nur Reflexion und Absorption $1-\sum e^{(r)}$.
Für die Ausgabe zusätzlich s/p-Zerlegung je Ordnung: mit $\kappa=(\alpha_n,\gamma_m)$,
$\hat{\mathbf s}=(\kappa_z,0,-\kappa_x)/|\kappa|$ (für $\kappa=0$: $\hat{\mathbf s}=\hat z$),
$\hat{\mathbf p}=\hat{\mathbf k}\times\hat{\mathbf s}$; Amplituden $\mathbf E\cdot\hat{\mathbf s}$, $\mathbf E\cdot\hat{\mathbf p}$
mit $E_y$ aus $\mathbf k\cdot\mathbf E=0$.

%=====================================================================
\section{S-Matrix-Formalismus}\label{sec:smatrix}
%=====================================================================

\subsection{Konventionen}
Zwischen zwei benachbarten Schichten liegt jeweils eine \emph{unendlich dünne Schicht aus dem
Referenzmedium} (Standard: Superstrat; konfigurierbar). Dort werden die Felder durch
$V=(\mathbf a;\mathbf b)$ in der Basis $\Psi_{\text{ref}}$ \eqref{eq:Psi_iso} beschrieben, mit dem Phasenbezug
\emph{lokal an dieser Grenzfläche}. Ein Block (Schicht, Schichtfolge) mit unterer Grenze \enquote{u}
und oberer Grenze \enquote{o} hat die S-Matrix
\begin{equation}
  \begin{pmatrix}\mathbf b_{\text{o}}\\\mathbf a_{\text{u}}\end{pmatrix}=
  \begin{pmatrix}S_{11}&S_{12}\\S_{21}&S_{22}\end{pmatrix}
  \begin{pmatrix}\mathbf b_{\text{u}}\\\mathbf a_{\text{o}}\end{pmatrix}
  \label{eq:Sdef}
\end{equation}
(auslaufend = $S\cdot$einlaufend; $S_{11}$: Transmission aufwärts, $S_{12}$: Reflexion von oben,
$S_{21}$: Reflexion von unten, $S_{22}$: Transmission abwärts; Buch (III.28)). Blockgröße
$n_h=2n_f$ (zwei Polarisationskomponenten je Ordnung; im 1D-in-plane-Fall $n_h=n_f$).

\subsection{Transfermatrix $\to$ S-Matrix}
Ist $T$ die Transfermatrix $(\mathbf a_{\text{o}};\mathbf b_{\text{o}})=T(\mathbf a_{\text{u}};\mathbf b_{\text{u}})$, so gilt
(Buch (III.30), (III.31), ergänzt):
\begin{equation}
  S_{11}=T_{22}-T_{21}T_{11}^{-1}T_{12},\quad S_{12}=T_{21}T_{11}^{-1},\quad
  S_{21}=-T_{11}^{-1}T_{12},\quad S_{22}=T_{11}^{-1}.
  \label{eq:T2S}
\end{equation}
Nur für dünne Scheiben verwenden (Abschnitt~\ref{sec:rk}); $T_{11}$ mit LU faktorisieren, nicht invertieren.

\subsection{Redheffer-Sternprodukt}
Für einen unteren Block $A$ und einen darüberliegenden Block $B$ (gemeinsame innere Grenzfläche mit
$\mathbf a_m,\mathbf b_m$) folgt aus $\mathbf b_m=A_{11}\mathbf b_{\text{in}}+A_{12}\mathbf a_m$,
$\mathbf a_m=B_{21}\mathbf b_m+B_{22}\mathbf a_{\text{in}}$:
\begin{equation}
\begin{aligned}
  D_1&=(\mathbb 1-A_{12}B_{21})^{-1}, & D_2&=(\mathbb 1-B_{21}A_{12})^{-1},\\
  C_{11}&=B_{11}D_1A_{11}, & C_{12}&=B_{12}+B_{11}D_1A_{12}B_{22},\\
  C_{21}&=A_{21}+A_{22}D_2B_{21}A_{11}, & C_{22}&=A_{22}D_2B_{22},
\end{aligned}
\qquad C=A\star B .
\label{eq:star}
\end{equation}
Neutrales Element: $S_{11}=S_{22}=\mathbb 1$, $S_{12}=S_{21}=0$. Das Buch-Rekursionsschema
(III.44)--(III.47) ist der Spezialfall \enquote{Stapel $\star$ Scheibe aus $T$}; es enthält in (III.45) einen Fehler
(Buch: $S_{21}$ statt $S_{22}$ im zweiten Term) -- im Code ausschließlich \eqref{eq:star} verwenden.
\textbf{Test:} Sternprodukt beliebiger Folgen = S aus dem Produkt der $T$-Matrizen (für gut konditionierte $T$).

\subsection{Grenzflächen und Gesamtstapel}
Übergang von einem homogenen Medium mit Basis $\Psi_X$ in die dünne Referenzschicht (Dicke 0):
$T=\Psi_{\text{ref}}^{-1}\Psi_X$ bzw.\ umgekehrt; S per \eqref{eq:T2S} (bei Dicke 0 unkritisch).
Diese Matrix ersetzt die Buchformeln (III.24)/(III.26) (dort sind $C$ und $C'$ vertauscht) und behebt
das Buchproblem $A'^{(1)}=A^{(1)}$ (gilt nur bei $\nu_1=\nu_M$). Gesamtstapel:
\begin{equation}
  S_{\text{tot}}=S_{\text{Substrat}\to\text{ref}}\star S_{\text{Schicht }2}\star\cdots\star S_{\text{Schicht }M-1}\star S_{\text{ref}\to\text{Superstrat}} .
  \label{eq:Stot}
\end{equation}
Mit Einfall nur von oben ($\mathbf b_{\text{Substrat,in}}=0$) und $\mathbf a_{\text{in}}=(A_x,A_z)$ in der Ordnung $(0,0)$:
\begin{equation}
  \mathbf b_{\text{Superstrat}}=S_{\text{tot},12}\,\mathbf a_{\text{in}},\qquad
  \mathbf a_{\text{Substrat}}=S_{\text{tot},22}\,\mathbf a_{\text{in}} .
  \label{eq:solution}
\end{equation}
Die Spalte der Ordnung $(0,0)$ genügt (\enquote{zentrale Spalte}); für Einfall aus mehreren Richtungen
(z.\,B.\ Reziprozitätstest) beliebige Spalten.

\subsection{Interne Amplituden (für das Nahfeld)}
Gespeichert werden für jede innere Grenzfläche $q$ die kumulierte S-Matrix des Teilstapels darunter,
$S_L^{(q)}$, und darüber, $S_U^{(q)}$. Mit den bekannten Eingängen $\mathbf b_{\text{in}}$ (unten, meist 0)
und $\mathbf a_{\text{in}}$ (oben):
\begin{equation}
  \big(\mathbb 1-S^{(q)}_{L,12}S^{(q)}_{U,21}\big)\mathbf b_q=S^{(q)}_{L,11}\mathbf b_{\text{in}}+S^{(q)}_{L,12}S^{(q)}_{U,22}\mathbf a_{\text{in}},\qquad
  \mathbf a_q=S^{(q)}_{U,21}\mathbf b_q+S^{(q)}_{U,22}\mathbf a_{\text{in}} .
  \label{eq:internal}
\end{equation}
Speicherbedarf: $2\cdot4n_h^2\cdot16$\,Byte je Grenzfläche (2D, $N=10$: $n_h=882$, ca.\ 100\,MB je
Grenzfläche). Alternative bei Speichermangel: nur $S_L$ speichern und $S_U$ in einem zweiten Durchlauf
von oben erzeugen.

%=====================================================================
\section{Schichtmodelle und Scheibenzerlegung}\label{sec:layers}
%=====================================================================

\subsection{Schichttypen}
\begin{center}
\small
\begin{tabular}{p{3.2cm}p{4.7cm}p{6.2cm}}
\toprule
Typ & Beschreibung & $\varepsilon$ bei Höhe $y$ / Normalenfeld \\
\midrule
\verb|Homogeneous| & Dicke $h$, Material & konstant; analytische Moden, keine Integration \\
\verb|Lamellar1D| & Intervalle $[x_j,x_{j+1}]$ mit Materialien, Dicke $h$ & \eqref{eq:interval}; $N_x^2=1$ \\
\verb|Profile1D| & $y=g(x)$ (Funktion oder Stützstellen), Material unten/oben, Tiefe $a$ &
 Schnittpunkte $g(x)=y$ per Nullstellensuche $\Rightarrow$ Intervalle; $\mathbf N$ aus \eqref{eq:N1d} \\
\verb|Polygon2D| & $y$-invarianter Querschnitt aus Polygonen/Rechtecken/Kreisen, Dicke $h$ &
 \eqref{eq:polygon}; $\mathbf N$ nach Algorithmus NF \\
\verb|Profile2D| & $y=g(x,z)$, Material unten/oben, Tiefe $a$ & Abtastung von $\{g>y\}$ + Polygonisierung
 (Marching Squares) oder feines FFT-Raster; $\mathbf N$ aus \eqref{eq:N2d} \\
\verb|Function| & benutzerdefiniertes $\overline{\overline\varepsilon}(x,y,z)$ & Abtastung + FFT (nur Tests) \\
\bottomrule
\end{tabular}
\end{center}
Beschichtete Profile (mehrere parallele Profile $g(x)+t_k$) werden als ein Schichtobjekt mit mehreren
Grenzflächen behandelt; $\mathbf N$ ist dann für alle Grenzflächen identisch (Buch IV.2.5).

\subsection{Scheibenzerlegung}
Jede modulierte Schicht wird in $K$ Scheiben $[y_{k-1},y_k]$ zerlegt (gleichmäßig oder nach Krümmung
von $g$). Kriterium gegen Kontamination (Buch III.6, mit Sicherheitsabstand):
\begin{equation}
  \rho_{\max}\,(y_k-y_{k-1})\le\ln(10^{8})\approx18.4,\qquad
  \rho_{\max}\approx\max_{(n,m)}\big|\sqrt{\alpha_n^2+\gamma_m^2-\varepsilon_{\min}}\big|,
  \label{eq:slicecrit}
\end{equation}
mit $\varepsilon_{\min}$ = Material mit dem kleinsten $\operatorname{Re}\varepsilon$ (Metalle: großes $|\rho|$).
\textbf{Achtung (Befund M3):} Die abgeschnittene FFF-Matrix kann Eigenwerte weit über dieser
Materialschätzung haben (Gold-Sinus: $|\lambda|\approx170$ statt $\approx6$). Der RK-Löser verwendet daher
$\rho_{\max}=\max(\text{Schätzung},\ \max_k|\lambda(M(y_k))|)$ an fünf Höhen $y_k$ und halbiert eine Scheibe
zusätzlich, solange $\max|T_F|>10^{8}$ (a-posteriori-Kontrolle).
RCWA: zusätzlich Treppenfehler $\Rightarrow$ $K$ per Konvergenztest. RK: Scheiben nur zur Stabilisierung,
die Genauigkeit regelt die Schrittweite.

\subsection{Permittivität bei beliebiger Höhe}
Der RK-Löser braucht $\tp{\varepsilon}(y)$, $\tp{1/\varepsilon}(y)$ bzw.\ $\tp{C_\varepsilon}(y)$ an den
Stufenknoten $y+c_ih$. \verb|Profile1D|: Nullstellen von $g(x)-y$ im Intervall $[0,d_x)$ (Abtasten auf
feinem Gitter, Vorzeichenwechsel, dann Brent/Newton), daraus Intervalle und exakte Koeffizienten
\eqref{eq:interval}. \verb|Profile2D|: Höhenlinie $g=y$ per Marching Squares als Polygon, dann
\eqref{eq:polygon}; Fallback FFT mit Oversampling. Cache: Koeffizienten je (Schicht, $y$) speichern, da
bei Feldrekonstruktion dieselben Knoten erneut gebraucht werden.

%=====================================================================
\section{Löser A: implizites Runge--Kutta (Differentialmethode)}\label{sec:rk}
%=====================================================================

\subsection{Verfahren}
Lineares System \eqref{eq:ode}. Ein $s$-stufiges RK-Verfahren mit Butcher-Tableau $(\mathsf A,\mathbf b,\mathbf c)$:
\begin{equation}
  K_i=M_i\Big(F+h\sum_{j=1}^{s}a_{ij}K_j\Big),\quad M_i=M(y+c_ih),\qquad
  F(y+h)=F+h\sum_{i=1}^{s}b_iK_i .
  \label{eq:rk}
\end{equation}
Für lineare Systeme ergibt sich ein \textbf{Schrittpropagator} (unabhängig von $F$):
\begin{equation}
  \big(\mathbb 1_{sn}-h\,\mathfrak A\big)X=\mathfrak M,\qquad
  \mathfrak A=\big(a_{ij}M_i\big)_{i,j=1}^{s},\quad \mathfrak M=\begin{pmatrix}M_1\\\vdots\\M_s\end{pmatrix},\qquad
  P_h=\mathbb 1_n+h\sum_{i=1}^{s}b_iX_i ,
  \label{eq:rkprop}
\end{equation}
mit $n=4n_f$ (bzw.\ $2n_f$ in 1D-in-plane), $X=(X_1;\dots;X_s)$, $X_i\in\mathbb C^{n\times n}$. Kosten pro
Schritt: eine LU-Zerlegung der Größe $sn$ und eine Lösung mit $n$ rechten Seiten.
Die Scheiben-Transfermatrix im Feldraum ist $T_F=\prod_kP_{h_k}$ (in Reihenfolge von unten nach oben).

\paragraph{Gauß--Legendre (Standard), $s=2$, Ordnung 4:}
\[
  \mathbf c=\Big(\tfrac12-\tfrac{\sqrt3}6,\ \tfrac12+\tfrac{\sqrt3}6\Big),\quad
  \mathsf A=\begin{pmatrix}\tfrac14&\tfrac14-\tfrac{\sqrt3}6\\\tfrac14+\tfrac{\sqrt3}6&\tfrac14\end{pmatrix},\quad
  \mathbf b=\big(\tfrac12,\tfrac12\big).
\]
\paragraph{Gauß--Legendre, $s=3$, Ordnung 6:}
\[
  \mathbf c=\Big(\tfrac12-\tfrac{\sqrt{15}}{10},\tfrac12,\tfrac12+\tfrac{\sqrt{15}}{10}\Big),\quad
  \mathsf A=\begin{pmatrix}\tfrac5{36}&\tfrac29-\tfrac{\sqrt{15}}{15}&\tfrac5{36}-\tfrac{\sqrt{15}}{30}\\
  \tfrac5{36}+\tfrac{\sqrt{15}}{24}&\tfrac29&\tfrac5{36}-\tfrac{\sqrt{15}}{24}\\
  \tfrac5{36}+\tfrac{\sqrt{15}}{30}&\tfrac29+\tfrac{\sqrt{15}}{15}&\tfrac5{36}\end{pmatrix},\quad
  \mathbf b=\big(\tfrac5{18},\tfrac49,\tfrac5{18}\big).
\]
\paragraph{Radau IIA, $s=3$, Ordnung 5 (optional, L-stabil):}
\[
  \mathbf c=\Big(\tfrac{4-\sqrt6}{10},\tfrac{4+\sqrt6}{10},1\Big),\quad
  \mathsf A=\begin{pmatrix}\tfrac{88-7\sqrt6}{360}&\tfrac{296-169\sqrt6}{1800}&\tfrac{-2+3\sqrt6}{225}\\
  \tfrac{296+169\sqrt6}{1800}&\tfrac{88+7\sqrt6}{360}&\tfrac{-2-3\sqrt6}{225}\\
  \tfrac{16-\sqrt6}{36}&\tfrac{16+\sqrt6}{36}&\tfrac19\end{pmatrix},\quad \mathbf b=\text{letzte Zeile von }\mathsf A .
\]
\textbf{Tests:} Ordnungsbedingungen ($\sum b_i=1$, $\sum b_ic_i=\tfrac12$, $\sum b_ic_i^2=\tfrac13$, \dots\ bis zur
Ordnung), $\mathsf A\mathbf 1=\mathbf c$; für konstantes $M$ ist $P_h$ bei Gauß--Legendre der
$(s,s)$-Padé-Approximant von $e^{hM}$ $\Rightarrow$ Vergleich mit \verb|expm|, Fehler $O(h^{2s+1})$ je Schritt.

\subsection{Stabilität und Schrittweite}
Gauß--Legendre ist A-stabil, $|R(z)|=1$ auf der imaginären Achse (propagierende Moden ohne
künstliche Dämpfung). Evaneszente Moden ($\lambda$ reell, groß) werden nur für $h|\lambda|\lesssim1$ genau
wiedergegeben; zu große Schritte verfälschen sie (bei Gauß nicht gedämpft, $|R(\pm\infty)|=1$).
Daher:
\begin{equation}
  h\le c_s/\rho_{\max},\qquad c_2\approx0.5,\ c_3\approx1.0,
  \label{eq:stepcrit}
\end{equation}
plus Fehlersteuerung per Schrittverdopplung: $\epsilon=\|P_hP_h-P_{2h}\|_F/\|P_{2h}\|_F$ mit
$P_h$ über $[y,y+h]$, $[y+h,y+2h]$; Ziel $\epsilon\le\texttt{rk\_tol}$ (Standard $10^{-8}$),
Schrittweitenanpassung $h\leftarrow h\,(\texttt{rk\_tol}/\epsilon)^{1/(2s+1)}\cdot0.9$.
Für $y$-unabhängiges $M$ (Lamellen, homogene Bereiche) genügt eine Padé-Auswertung mit Skalieren
und Quadrieren: $P=(P_{h/2^k})^{2^k}$.

\subsection{Scheiben-S-Matrix aus dem RK-Ergebnis}
\begin{equation}
  T^{(\text{Scheibe})}=\Psi_{\text{ref}}^{-1}\,T_F\,\Psi_{\text{ref}},\qquad S^{(\text{Scheibe})}\ \text{aus \eqref{eq:T2S}}
  \label{eq:rkslice}
\end{equation}
(Buch Anh.~IV (A.IV.14'); dort Blockgröße für 2D falsch angegeben: richtig $n_h=2n_f$).
Die Scheibendicke folgt \eqref{eq:slicecrit}, damit $T_{11}$ gut konditioniert ist; Diagnose:
Konditionszahl von $T_{11}$ schätzen und bei $>10^{10}$ die Scheibe halbieren.

\subsection{Dichte Ausgabe für das Nahfeld}
Gauß--Kollokation liefert ein Polynom vom Grad $s$ je Schritt:
\begin{equation}
  F(y+\vartheta h)=F(y)+h\sum_{j=1}^{s}\Big(\int_0^{\vartheta}\ell_j(\tau)\,d\tau\Big)K_j,\qquad 0\le\vartheta\le1,
  \label{eq:dense}
\end{equation}
mit den Lagrange-Polynomen $\ell_j$ zu den Knoten $\mathbf c$. Für das Nahfeld wird nur \emph{ein}
Zustandsvektor (die tatsächliche Lösung) von der Unterkante jeder Scheibe aus neu integriert
(dieselbe Schrittfolge, gespeichert in \verb|RKSliceRecord|).

\subsection{Rechenaufwand (Richtwerte)}
Pro Schritt $\approx\frac83(sn)^3$ reelle Flops (LU) $+8(sn)^2n$ (Lösen). 2D, $N_x=N_z=10$:
$n=1764$, $s=2$: $\approx4.7\cdot10^{11}+2\cdot10^{11}$ Flops. Der RK-Löser ist daher für 2D teuer und
hauptsächlich für Profile $g(x,z)$ und Referenzrechnungen gedacht; für $y$-invariante 2D-Schichten
ist RCWA vorzuziehen.

%=====================================================================
\section{Löser B: RCWA (Eigenmoden je Scheibe)}\label{sec:rcwa}
%=====================================================================

\subsection{Eigenwertproblem}
In einer Scheibe mit $y$-unabhängigem $M$ (Treppenapproximation: $\varepsilon$ in Scheibenmitte):
$M=W\Lambda W^{-1}$, $F(y)=\sum_kc_kw_ke^{\lambda_k(y-y_{\text{ref},k})}$.
\paragraph{Reduziertes Problem (isotrop/lamellar, $M_{11}=M_{12}=M_{21}=M_{22}=0=M_{33}=M_{34}=M_{43}=M_{44}$):}
Mit $\mathcal P=\begin{pmatrix}M_{13}&M_{14}\\M_{23}&M_{24}\end{pmatrix}$, $\mathcal R=\begin{pmatrix}M_{31}&M_{32}\\M_{41}&M_{42}\end{pmatrix}$:
\begin{equation}
  M=\begin{pmatrix}0&\mathcal P\\\mathcal R&0\end{pmatrix},\quad
  \mathcal P\mathcal R\,u=\mu\,u,\quad \kappa=\sqrt{-\mu}\ (\text{Zweig \eqref{eq:beta_branch}}),\quad
  w_\pm=\begin{pmatrix}u\\\mathcal Ru/\lambda_\pm\end{pmatrix},\ \lambda_\pm=\pm i\kappa .
  \label{eq:rcwa_reduced}
\end{equation}
(Prüfung: $Mw_\pm=(\mathcal P\mathcal Ru/\lambda_\pm;\mathcal Ru)=\lambda_\pm w_\pm$ wegen $\mu=\lambda_\pm^2$.)
\enquote{Auf}: $\lambda_+=i\kappa$ ($\operatorname{Re}\lambda_+\le0$), \enquote{ab}: $\lambda_-=-i\kappa$.
Für $\kappa\approx0$ (Grenzfall) $w=(u;\mathcal Ru/\lambda)$ schlecht definiert $\Rightarrow$ kleine
Regularisierung wie bei \eqref{eq:beta_branch}.
\paragraph{Allgemein (anisotrop, Profile mit $Q_{yx}\neq0$):} volle Eigenzerlegung der $4n_f$-Matrix.
Klassifikation (robuste Form von \eqref{eq:classify}): Bewertung $s_k=\operatorname{Re}\lambda_k$, für
$|\operatorname{Re}\lambda_k|\le\tau$ stattdessen $\mp\tau/2$ nach dem Vorzeichen von $S_y$; sortieren und die untere
Hälfte als aufwärts nehmen. (Die reine Schwellenwert-Variante scheiterte in metallischen Scheiben.)
\paragraph{TE/TM-Entkopplung:} Im 1D-in-plane-Fall verschwinden alle Blöcke zwischen $(E_x,H_z)$ und
$(E_z,H_x)$ exakt. RK und RCWA erkennen das automatisch und rechnen beide Hälften getrennt
(Laufzeitgewinn RK $\approx6$, RCWA $\approx2{,}5$).

\subsection{Stabile Scheiben-S-Matrix}
Moden auf ($W_u,\Lambda_u$, Bezug Unterkante $y_u$) und ab ($W_d,\Lambda_d$, Bezug Oberkante $y_o$):
\begin{equation}
  F(y)=W_ue^{\Lambda_u(y-y_u)}\mathbf c_u+W_de^{\Lambda_d(y-y_o)}\mathbf c_d,\qquad
  X_u=e^{\Lambda_uh},\ X_d=e^{-\Lambda_dh}\quad(|X|\le1).
  \label{eq:rcwa_field}
\end{equation}
Mit $\Psi_{\text{ref}}^{-1}W_u=\begin{pmatrix}A_u\\B_u\end{pmatrix}$, $\Psi_{\text{ref}}^{-1}W_d=\begin{pmatrix}A_d\\B_d\end{pmatrix}$
(obere Hälfte: ab-Amplitude, untere: auf-Amplitude) gilt an Unter- und Oberkante
\begin{equation}
  \begin{aligned}
  \mathbf a_u&=A_u\mathbf c_u+A_dX_d\mathbf c_d, & \mathbf b_u&=B_u\mathbf c_u+B_dX_d\mathbf c_d,\\
  \mathbf a_o&=A_uX_u\mathbf c_u+A_d\mathbf c_d, & \mathbf b_o&=B_uX_u\mathbf c_u+B_d\mathbf c_d .
  \end{aligned}
  \label{eq:rcwa_bc}
\end{equation}
Die Eingänge $(\mathbf b_u,\mathbf a_o)$ bestimmen
\begin{equation}
  \underbrace{\begin{pmatrix}B_u&B_dX_d\\A_uX_u&A_d\end{pmatrix}}_{=:G}\begin{pmatrix}\mathbf c_u\\\mathbf c_d\end{pmatrix}=
  \begin{pmatrix}\mathbf b_u\\\mathbf a_o\end{pmatrix},\qquad
  \begin{pmatrix}\mathbf b_o\\\mathbf a_u\end{pmatrix}=
  \underbrace{\begin{pmatrix}B_uX_u&B_d\\A_u&A_dX_d\end{pmatrix}}_{=:H}\begin{pmatrix}\mathbf c_u\\\mathbf c_d\end{pmatrix},
  \qquad S=HG^{-1}.
  \label{eq:rcwa_S}
\end{equation}
Alle Einträge sind beschränkt; keine wachsenden Exponentiale. $G$ als LU speichern
(\verb|RCWASliceRecord|) für die Feldrekonstruktion. Homogene Zwischenschichten: dieselben Formeln mit
analytischen Moden ($W=\Psi_X$, $\Lambda=\operatorname{diag}(\mp i\beta)$).

\subsection{Querprüfung RK $\leftrightarrow$ RCWA}
Für $y$-unabhängiges $M$ (Lamellen) müssen beide Löser bis auf die RK-Toleranz identische Effizienzen
liefern; für glatte Profile müssen RCWA mit $K\to\infty$ Scheiben und RK gegen denselben Wert
konvergieren -- \emph{außer} im TM-Fall bei Metallen mit Li-Lamellen-Faktorisierung (Buch VI.5:
Treppenapproximation konvergiert dort extrem langsam). Deshalb verwendet RCWA auf Profilschichten
\textbf{dieselbe FFF-Faktorisierung mit dem Profil-Normalenfeld} ($\varepsilon$ in Scheibenmitte, $\mathbf N$ des
glatten Profils), nicht die Lamellenformel. Kapitelprüfung: Al-Sinusgitter, FFF mit 60/150/300 Scheiben
$\to0{,}853/0{,}865/0{,}869$, Referenz $0{,}8696$ (C-Methode).

%=====================================================================
\section{Nahfeld}\label{sec:field}
%=====================================================================

\subsection{Ablauf}
\begin{enumerate}[nosep]
  \item Gesamtlösung \eqref{eq:solution}, interne Amplituden $(\mathbf a_q,\mathbf b_q)$ an allen Grenzflächen \eqref{eq:internal}.
  \item Fourierkoeffizienten $F(y)=([E_x],[E_z],[H_x],[H_z])$ an den gewünschten Höhen:
    \begin{itemize}[nosep]
      \item homogene Schicht/Superstrat/Substrat: Moden mit Phasen $e^{\lambda(y-y_{\text{ref}})}$;
      \item RCWA-Scheibe: $(\mathbf c_u;\mathbf c_d)=G^{-1}(\mathbf b_u;\mathbf a_o)$, dann \eqref{eq:rcwa_field};
      \item RK-Scheibe: $F(y_u)=\Psi_{\text{ref}}(\mathbf a_u;\mathbf b_u)$, erneute Integration mit dichter Ausgabe \eqref{eq:dense}.
    \end{itemize}
  \item $[E_y]$, $[H_y]$ aus \eqref{eq:Ey}, \eqref{eq:Hy} mit den $Q$-Matrizen der Höhe $y$.
  \item Synthese im Ortsraum: $u(x,y,z)=\sum_ru_r(y)e^{i(\alpha_nx+\gamma_mz)}$, effizient per inverser FFT
        auf einem Gitter (Zero-Padding), danach Multiplikation mit $e^{i(\alpha_0x+\gamma_0z)}$.
\end{enumerate}

\subsection{Umsetzung (M6)}
\begin{itemize}[nosep]
  \item Ein Lauf mit Aufzeichnung speichert jedes Element des Stapels (Grenzfläche zum Substrat, modale Scheibe
        mit $W$, $\Lambda$, $G$, RK-Scheibe mit den Transfermatrizen $T(y_0\to y_j)$ an genau den angeforderten
        Höhen $y_j$). Die RK-Integration hält dafür an den Abtasthöhen an; dichte Ausgabe \eqref{eq:dense}
        ist damit nicht nötig.
  \item Interne Amplituden aus \eqref{eq:internal} mit kumulierten $S_L$, $S_U$ je Grenzfläche.
  \item Schnittstelle \verb|Solver::fieldCoefficients(stack, wl, y)| liefert $F$, $[E_y]$, $[H_y]$ je Höhe;
        \verb|Solver::nearField| erzeugt die Schnitte. Gemeinsame Schichtmodelle (\verb|detail::LayerModel|:
        $Q(y)$, $\overline{\overline\varepsilon}(x,y)$, $\mathbf N(x)$, $\tp{\mathbf N\mathbf N^T}$) sorgen dafür, dass
        Lösung und Rekonstruktion dieselben Daten verwenden (RCWA-Treppe: $Q$ und $\varepsilon$ der Scheibenmitte).
\end{itemize}

\subsection{FFF-konsistente Rekonstruktion (empfohlen in modulierten Schichten)}
Die direkte Synthese von $E_x,E_y$ aus abgeschnittenen Reihen zeigt Gibbs-Oszillationen an
Materialgrenzen. Besser: stetige Größen synthetisieren und punktweise zurückrechnen:
\begin{equation}
  [\mathbf F_\varepsilon]=\tp{C_\varepsilon}^{-1}[\mathbf E]\ \ \text{(Variante A)}\ \Rightarrow\ \mathbf F_\varepsilon(x,z)\ \text{(stetig)}
  \ \Rightarrow\ \mathbf E(x,z)=C_\varepsilon(x,z)\,\mathbf F_\varepsilon(x,z),
  \label{eq:recon}
\end{equation}
umgesetzt (für isotrope und anisotrope Schichten gleichermaßen) über die Projektorform \eqref{eq:projector}:
$[\mathbf D]=Q[\mathbf E]$, $[\mathbf d_n]=\tp{P_N}[\mathbf D]$, $[\mathbf e_t]=[\mathbf E]-\tp{P_N}[\mathbf E]$ synthetisieren
(stetig) und punktweise $\mathbf E(x)=H_1(x)\,\mathbf e_t(x)+\xi(x)^{-1}\mathbf d_n(x)$; isotrop ist das
$\mathbf E=\mathbf e_t+\mathbf d_n/\varepsilon$. Befund (V12): an einer Lamellenwand springt $D_x=\varepsilon E_x$ bei
direkter Synthese um 270\,\%, FFF-konsistent um 3\,\% (Gitterabstand 2\,nm). $\mathbf H$ ist überall stetig ($\mu=1$) und wird direkt synthetisiert.
Beide Varianten (direkt / FFF-konsistent) sind auswählbar (\verb|reconstruction: direct \| fff|).

\subsection{Schnitte und Ausgabe}
\begin{center}
\small
\begin{tabular}{lll}
\toprule
& horizontaler Schnitt ($y=y_0$) & vertikaler Schnitt \\
\midrule
1D & Linie $u(x,y_0)$, $x\in[x_a,x_b]$ & Karte $u(x,y)$, $z=0$ \\
2D & Karte $u(x,y_0,z)$ & Karte $u(x,y,z_0)$ (optional zusätzlich $u(x_0,y,z)$) \\
\bottomrule
\end{tabular}
\end{center}
Größen: alle sechs komplexen Komponenten, $|\mathbf E|^2$, $|\mathbf H|^2$, $\operatorname{Re}(\mathbf Ee^{-i\omega t})$ für
gegebene Phase, Poynting-Vektor $\tfrac12\operatorname{Re}(\mathbf E\times\mathbf H^*)$, Materialindex-Karte (für Konturen).
Ausgabeformat: je Schnitt ein Verzeichnis (\verb|nearfield_horizontal|, \verb|nearfield_vertical|) mit
\verb|Ex|\dots\verb|Hz.npy| (complex128), \verb|absE2|, \verb|absH2|, \verb|Sy|, \verb|material.npy| (float64),
C-Ordnung, Form $(n_y,n_x)$ bzw.\ $(n_x)$, plus \verb|meta.json|; $\mathbf H$ als $Z_0\mathbf H$, $|\mathbf E_{\text{inc}}|=1$ (Achsen in µm, Wellenlänge,
Einfall, Komponentenliste). Optional HDF5 (\verb|FG_ENABLE_HDF5|). Das Skript
\verb|scripts/plot_nearfield.py| erzeugt PNG/PDF.

%=====================================================================
\section{Numerische Details und Robustheit}\label{sec:numerics}
%=====================================================================
\begin{itemize}
  \item Durchgehend \verb|std::complex<double>| bzw.\ \verb|cuDoubleComplex|; keine einfache Genauigkeit
        im Kern. (Gemischte Genauigkeit nur optional für LU mit iterativer Verfeinerung, M10.)
  \item Matrizen nie explizit invertieren, wenn ein Gleichungssystem genügt (LU + Lösen).
        Ausnahme: $\Li$, $\tp{C_\varepsilon}^{-1}$, $Q_{\varepsilon,yy}^{-1}$, die mehrfach gebraucht werden
        (dann LU speichern und Produkte per Lösen bilden).
  \item Eigenwertklassifikation mit Toleranzen relativ zu $\max|\lambda|$; entartete Eigenwerte
        (Symmetrie, Normaleinfall) sind erlaubt -- die Eigenvektoren spannen den Raum auf, die
        S-Matrix ist davon unabhängig.
  \item Rayleigh-Anomalien: \eqref{eq:beta_branch}; zusätzlich Option, die Wellenlänge um
        $10^{-12}$ relativ zu verschieben.
  \item Diagnosen je Rechnung: Energiebilanz $\sum e^{(r)}+\sum e^{(t)}$ (verlustfrei $=1$), Konditionszahlen
        ($T_{11}$, $G$, $Q_{\varepsilon,yy}$), Anzahl RK-Schritte, Klassifikationsreserve.
  \item Reproduzierbarkeit: deterministische Reihenfolgen, keine Zufallszahlen im Kern.
  \item \textbf{Eigen-Falle:} \verb|dot()| konjugiert das erste Argument, \verb|cross()| liefert für komplexe
        Vektoren das konjugierte Kreuzprodukt. Bilineare Formen ($\xi=\mathbf N^T\overline{\overline\varepsilon}\mathbf N$,
        $\mathbf v\times\mathbf T$, $\mathbf E\cdot\hat{\mathbf s}$) explizit ausschreiben (Befund M5).
  \item Symmetrien (optional, später): spiegelsymmetrische Strukturen bei Normaleinfall erlauben
        Halbierung der Ordnungen -- \emph{nicht} in der ersten Version.
\end{itemize}

%=====================================================================
\section{CUDA-Backend (RTX 6000 Ada)}\label{sec:cuda}
%=====================================================================

\subsection{Hardwareeigenschaften und Konsequenzen}
RTX 6000 Ada: Compute Capability 8.9 (\verb|sm_89|), 48\,GB GDDR6. Die FP64-Rate beträgt bei dieser
Architektur nur einen kleinen Bruchteil der FP32-Rate (Größenordnung 1/64; vor M10 mit einem
ZGEMM-Benchmark messen). Konsequenzen:
\begin{itemize}[nosep]
  \item Die GPU lohnt sich vor allem bei großen dichten Matrizen (2D, $n\gtrsim1000$) und bei vielen
        unabhängigen Rechnungen (Wellenlängen-/Winkel-Scans, gebatcht).
  \item Für 1D-Probleme ($n\le$ einige Hundert) ist der CPU-Pfad meist gleich schnell oder schneller;
        die Backend-Auswahl erfolgt automatisch nach Problemgröße (Schwelle konfigurierbar).
  \item Der CPU-Pfad bleibt Referenz für alle Tests (GPU-Ergebnisse müssen auf $10^{-10}$ relativ übereinstimmen).
\end{itemize}

\subsection{Bibliotheken}
\begin{center}
\small
\begin{tabular}{p{2.6cm}p{6.3cm}p{5.3cm}}
\toprule
Operation & GPU & CPU \\
\midrule
GEMM & cuBLAS \verb|cublasZgemm| (+ \verb|StridedBatched|) & Eigen / BLAS (OpenBLAS, FlexiBLAS) \\
LU / Lösen & cuSOLVER \verb|cusolverDnXgetrf|, \verb|Xgetrs| (64-Bit-API) & Eigen \verb|PartialPivLU| / LAPACK \verb|zgetrf| \\
Eigenwerte (allgemein) & cuSOLVER \verb|cusolverDnXgeev| (ab CUDA 12.6 Update 2; hybrid CPU/GPU, gepinnter Host-Speicher) & Eigen \verb|ComplexEigenSolver| / LAPACK \verb|zgeev| \\
FFT & cuFFT (Z2Z) & FFTW3 oder PocketFFT (header-only) \\
\bottomrule
\end{tabular}
\end{center}
Eigene Kernel: Aufbau (block-)Toeplitz-Matrizen aus Koeffizientengittern, Diagonalskalierungen
($\alpha X$, $X\gamma$), punktweise Produkte $C_\varepsilon(x)\mathbf F(x)$, Feldsynthese-Nachbearbeitung.
Zielarchitektur: \verb|CMAKE_CUDA_ARCHITECTURES=89|.
Speicherverwaltung über RAII-Klassen (\verb|DeviceBuffer<T>|), Streams je Wellenlänge,
Fehlerprüfung per Makro nach jedem Aufruf.

\subsection{Schnittstelle}
Der Physikcode ruft nur die abstrakte Schnittstelle \verb|fg::linalg::Backend| (GEMM, LU, Solve, Eig,
FFT, Speichertransfers) auf. Implementierungen: \verb|CpuBackend| (Pflicht, zuerst), \verb|CudaBackend|
(M10). Kein CUDA-Code außerhalb von \verb|src/cuda/| und \verb|src/linalg/cuda_backend.*|.

\subsection{Umsetzung (M10, v1.8)}
\begin{itemize}[nosep]
  \item \verb|CudaBackend| (\verb|src/cuda/cuda_backend.cpp|) nutzt nur die Host-API: \verb|cublasZgemm|,
        \verb|cusolverDnXgetrf|/\verb|Xgetrs| (64 Bit), \verb|cusolverDnXgeev| (rechte Eigenvektoren, auf
        Einheitsnorm normiert wie Eigen; Host-Workspace gepinnt). RAII für Device-/Pinned-Speicher, je Backend
        ein Stream, Fehlerprüfung per Makro. Ohne eigene Kernel ist kein \verb|nvcc| nötig
        (\verb|find_package(CUDAToolkit 12.6)|, \verb|CUDA::cudart|, \verb|cublas|, \verb|cusolver|); eigene
        Kernel (Toeplitz-Aufbau, Synthese) und cuFFT folgen bei Bedarf (M11).
  \item Dispatch: Die $O(n^3)$-Stellen (RCWA-Eigenprobleme, stabile Schicht-S-Matrix, Sternprodukt, $T\to S$,
        RK-Schrittpropagator) fragen \verb|linalg::offload(n)|; ist es falsch, läuft der unveränderte
        Eigen-Code (CPU-Referenz bitgleich), sonst \verb|linalg::device()|. Einstellung thread-lokal über
        \verb|SolverOptions::backend| (\verb|auto|, \verb|cpu|, \verb|cuda|) und \verb|backendThreshold| (1000).
  \item Batching über Wellenlängen: \verb|solveScan| verteilt die Scanpunkte auf OpenMP-Threads, jeder mit
        eigenem Löser und (GPU) eigenem Backend/Stream; Ergebnisse identisch mit sequentiellen Rechnungen.
  \item Ohne GPU (Sandbox) werden die GPU-Tests übersprungen; Bau und Linken gegen CUDA 12.8 sind geprüft,
        die Abnahme GPU $=$ CPU ($10^{-10}$) und der Benchmark-Bericht (\verb|docs/benchmarks.md|) erfolgen auf
        dem Zielsystem.
\end{itemize}

%=====================================================================
\section{Software-Architektur}\label{sec:arch}
%=====================================================================

\subsection{Verzeichnisse}
{\footnotesize\begin{verbatim}
include/fg/core      Typen, Konstanten, Fehler, IndexMap
include/fg/fourier   Koeffizienten (Intervall/Rechteck/Polygon/FFT), Toeplitz
include/fg/geometry  Material, Schichttypen, Stapel, Normalenfelder, Scheibenzerlegung
include/fg/physics   Einfall, Q-Matrizen (FFF), Propagationsmatrix M, Modenbasen, Effizienzen
include/fg/smatrix   S-Matrix, Sternprodukt, T->S, Gesamtlösung, interne Amplituden
include/fg/solvers   Butcher-Tableaus, RK-Löser, RCWA-Löser, gemeinsame Schnittstelle
include/fg/field     Nahfeld: Schnitte, Rekonstruktion, Synthese
include/fg/linalg    Backend-Schnittstelle, CPU- und CUDA-Backend
include/fg/io        JSON-Konfiguration, Ergebnis- und Feldausgabe
src/...              Implementierungen (gleiche Gliederung), src/cuda/ Kernel
apps/fgsolve         Kommandozeilenprogramm
tests/unit           GoogleTest, je Modul
tests/validation     Referenzfälle (Abschnitt Validierung), CTest-Label "validation"
examples/            JSON-Konfigurationen der Referenzfälle
scripts/             Setup AlmaLinux 10, Plot-Skripte, Doku-Build
docs/                dieser Plan, Architektur, Roadmap, Validierung, Konventionen
\end{verbatim}}

\subsection{Zentrale Typen (Skizze)}
{\footnotesize\begin{verbatim}
using cplx = std::complex<double>;  using MatC = Eigen::MatrixXcd;  using VecC = Eigen::VectorXcd;
struct Material { bool isotropic; cplx eps; Eigen::Matrix3cd epsTensor; };   // mu = 1
struct Incidence { double lambda_um, phi_deg, psi_deg, delta_deg; };
enum class Dim { D1, D2 };   enum class Factorization { Laurent, LiLamellar, FFF_iso, FFF_A, FFF_B };
struct Truncation { int Nx, Nz; };                      // Nz = 0 im 1D-Fall
class IndexMap;                                         // (n,m) <-> r
struct LayerSpec { /* variant: Homogeneous, Lamellar1D, Profile1D, Polygon2D, Profile2D */ };
struct Stack { Material substrate, superstrate; std::vector<LayerSpec> layers; double dx, dz; };
struct QBlocks { MatC b[3][3]; };                       // Q_eps,ij
MatC buildPropagationMatrix(const QBlocks&, const VecC& alpha, const VecC& gamma); // eq:Mgeneral
struct ModeBasis { MatC Psi, PsiInv; VecC lambda; };    // homogene Medien
struct SMatrix { MatC S11, S12, S21, S22; };            // Konvention eq:Sdef
SMatrix star(const SMatrix& lower, const SMatrix& upper);
SMatrix fromTransfer(const MatC& T);
class Solver { virtual SMatrix sliceSMatrix(...) = 0; };  // RKSolver, RCWASolver
struct Result { std::vector<OrderEfficiency> R, T; double absorption, energyBalance; Diagnostics diag; };
struct NearFieldRequest { double y_horizontal; double z_vertical; GridSpec grid; Quantities q; };
\end{verbatim}}

\subsection{Konfiguration (JSON, Kurzfassung; Details in \texttt{docs/config\_format.md})}
{\footnotesize\begin{verbatim}
{ "dimension": "1D", "solver": "rk" | "rcwa",
  "truncation": {"Nx": 15, "Nz": 0}, "factorization": "auto",
  "period_um": {"x": 0.5, "z": 0.0},
  "incidence": {"lambda_um": 0.6328, "phi_deg": 30, "psi_deg": 0, "delta_deg": 90},
  "superstrate": {"n": 1.0}, "substrate": {"n": [1.38, 7.62]},
  "layers": [{"type": "profile1d", "shape": "sinus", "depth_um": 0.2,
              "below": "substrate", "above": "superstrate", "slices": 40}],
  "rk": {"tableau": "gauss2", "tol": 1e-8}, "rcwa": {"slices": 60},
  "nearfield": {"horizontal_y_um": 0.1, "vertical_z_um": 0.0, "nx": 400, "ny": 300,
                "quantities": ["E","H","absE2"], "reconstruction": "fff"},
  "output": {"dir": "out/al_sinus"} }
\end{verbatim}}

\subsection{Code-Regeln}
C++20, \verb|clang-format| (Datei im Repo), keine rohen \verb|new|/\verb|delete|, Fehler über
\verb|fg::Error| (Ausnahmen) mit Kontext, keine globalen Zustände, Einheiten nur an der IO-Grenze
umrechnen, jede öffentliche Funktion mit Kommentar inkl.\ Formelreferenz auf diesen Plan
(z.\,B.\ \verb|// Plan (7.3)|).

%=====================================================================
\section{Validierung und Tests}\label{sec:validation}
%=====================================================================

\subsection{Unit-Tests (Pflicht, je Modul)}
\begin{enumerate}[nosep]
  \item IndexMap: Bijektion $(n,m)\leftrightarrow r$ für alle Paare; 1D-Spezialfall.
  \item Fourierkoeffizienten: Intervall vs.\ numerische Quadratur; Polygon \eqref{eq:polygon} vs.\ Rechteck;
        Summe der Koeffizienten zweier komplementärer Gebiete $=\delta_{p0}$.
  \item Toeplitz: Laurent-Regel für glatte Funktionen (Inneres exakt); Beispiel Buch IV.1
        (inverse Regel bei $N=0$ exakt $=1$, Laurent $1{,}5625$).
  \item Butcher-Tableaus: Ordnungsbedingungen; $P_h$ vs.\ \verb|expm| (Padé-Ordnung).
  \item Propagationsmatrix: homogenes isotropes Medium $\Rightarrow$ $M\Psi=\Psi\Lambda$ mit \eqref{eq:Psi_iso};
        Vergleich der 16 Blöcke mit direkter Elimination aus \eqref{eq:mx1}--\eqref{eq:mx3} (Zufalls-$Q$).
  \item S-Matrix: Sternprodukt vs.\ T-Produkt; neutrales Element; Assoziativität.
  \item $Q_\varepsilon$: Grenzfälle aus Abschnitt~\ref{sec:fff}; $A_\varepsilon C_\varepsilon=\mathbb 1$ für Zufallstensoren;
        Variante B isotrop $\approx$ \eqref{eq:Qiso}.
  \item Modenbasis anisotrop: isotroper Tensor reproduziert \eqref{eq:Psi_iso}; uniaxial mit optischer Achse
        $\parallel y$: $\lambda=\pm i\sqrt{\varepsilon_o-\kappa^2}$ (TE-artig) und $\pm i\sqrt{\varepsilon_o(1-\kappa^2/\varepsilon_e)}$.
\end{enumerate}

\subsection{Referenzfälle (Validierung)}
\begin{center}
\small
\begin{tabular}{p{0.6cm}p{6.6cm}p{6.6cm}}
\toprule
ID & Fall & Referenz / Kriterium \\
\midrule
V1 & ebene Grenzfläche, alle Winkel, TE/TM, beliebiger Azimut & Fresnel analytisch, $10^{-12}$ \\
V2 & homogene Schicht auf Substrat (Airy), auch absorbierend & analytisch, $10^{-12}$ \\
V3 & Lamellengitter dielektrisch ($\nu=1{,}5$ bzw.\ $2{,}5$), TE/TM/konisch & Energie $=1\pm10^{-10}$;
     RK = RCWA auf RK-Toleranz; Reziprozität nur bis auf Abschneidung (die Ordnungsmengen beider
     Einfallsrichtungen sind um eine Ordnung verschoben): Abweichung muss mit $N$ fallen \\
V4 & Al-Sinusgitter: $d=0{,}5\,\mu$m, $g=\tfrac a2(1+\cos Kx)$, $a=0{,}2\,\mu$m, $\lambda=0{,}6328\,\mu$m,
     $\Phi=30^\circ$, $\nu=1{,}38+7{,}62i$, Ordnung $-1$ & TE $0{,}3235$, TM $0{,}8696$ (intern, C-Methode;
     FFF-Treppe 300 Scheiben $0{,}8687$); Toleranz $2\cdot10^{-3}$ \\
V5 & tiefes Glas-Sinusgitter, $d=0{,}5$, $a=1{,}0\,\mu$m, $\nu=1{,}5$, $30^\circ$, Transmission $-1$ & TE $0{,}820$
     (intern, 120 Scheiben); Buch Abb.\ III.4 $\approx0{,}82$ \\
V6 & Dielektrisches Sinusgitter $\nu=2{,}5$, $a/d=1$, $d=1\,\mu$m, $30^\circ$ (Buch Abb.\ IV.3), TM &
     FFF-Fehler $10^{-4}$ bei $N\approx8$ (qualitativ: FFF $\gg$ klassisch) \\
V7 & Kobalt-Sinusgitter anisotrop (Buch Kap.\ VIII): $d=0{,}6$, $a=0{,}5\,\mu$m, $30^\circ$, TM-Einfall,
     $\varepsilon_{ii}=-8{,}19+16{,}38i$, $\varepsilon_{xz}=-\varepsilon_{zx}=-0{,}495-0{,}106i$ &
     intern ($N=20$, 100 Scheiben): TM$_0$ $0{,}0567$, TM$_{-1}$ $0{,}1435$, TE$_0$ $1{,}3\cdot10^{-7}$,
     TE$_{-1}$ $2{,}37\cdot10^{-6}$; Buch $\approx0{,}055/0{,}148/\approx10^{-7}/2{,}6\cdot10^{-6}$; Toleranz 3\,\%.
     M5: $0{,}05673/0{,}14355/1{,}275\cdot10^{-7}/2{,}372\cdot10^{-6}$; RK ($N=8$) = RCWA (400 Scheiben) auf $10^{-4}$ \\
V7b & einachsiges Substrat, Achse $\parallel y$, $\varepsilon_o=2{,}4$, $\varepsilon_e=3{,}1$ & analytisch ($r_s$, $r_p$ s.\ Text), $10^{-12}$ \\
V7c & verlustfreie biaxiale gedrehte Schicht, konisch & Energie $10^{-12}$; homogener Pfad = Lamellenpfad (RCWA $10^{-10}$, RK $10^{-8}$) \\
V8 & Anisotroper Code mit isotropem Tensor & identisch mit isotropem Code ($10^{-10}$ bei gleicher Faktorisierung).
     2D (M8): Laurent mit Tensor = isotrop (Polygon, Eierkarton); \verb|FFF_B| auf $z$-invariantem Streifen = Li;
     anisotrope $z$-invariante Strukturen = 1D-anisotrop ($10^{-10}$/$10^{-12}$); Spiegeltest $x\leftrightarrow z$
     mit gespiegeltem Tensor ($10^{-9}$). Konvergenz: Co-Säulen \verb|FFF_B| $T=0{,}548\ldots0{,}551$ ($N=2\ldots7$),
     Laurent $0{,}42\ldots0{,}49$; anisotroper Eierkarton: A, B, Laurent konsistent \\
V9 & 2D-Code mit $z$-invarianter Struktur ($g(x,z)=g(x)$, $N_z=0$) & identisch mit 1D-konisch.
     M7: Lamellen, \texttt{Polygon2D}-Streifen (NF-Tensor, \verb|FFF_sym|/\verb|FFF_iso| = Li),
     \verb|sinus_x| und 1D-Profil im 2D-Stapel, RK und RCWA: $10^{-12}$ ($N_z=0$); RCWA mit $N_z>0$: Ordnungen
     $(n,0)$ $10^{-10}$, übrige Ordnungen leistungsfrei; RK mit $N_z>0$ nur auf RK-Toleranz (Schrittweitensteuerung
     sieht die entkoppelten Blöcke) \\
V10 & Lochgitter (Buch IX.1.7): Ag-Film $0{,}2\,\mu$m auf Glas, quadratische Löcher $0{,}25\,\mu$m,
      Periode $0{,}9\,\mu$m, $\nu_{\text{Ag}}=0{,}1+8{,}94i$, Normaleinfall & $T(\lambda=1{,}388\,\mu$m$)\approx7{,}5\,\%$
      (Buch); Plasmonpol der ebenen Grenzfläche bei $1{,}369406\,\mu$m. M7 (RCWA, \verb|FFF_sym|, NF-Tensor):
      Peak stets oberhalb des Pols; $\lambda_{\text{Peak}}=1{,}383/1{,}386/1{,}391/1{,}404/1{,}3876/1{,}3889\,\mu$m und
      $T_{\max}=4{,}1/3{,}4/10{,}3/5{,}3/4{,}2/6{,}2\,\%$ für $N=5\ldots10$: ab $N=9$ Peak bei der Buch-Wellenlänge,
      Höhe nähert sich 7,5\,\%; Tensor- und Vektorvariante von Alg.~NF gleichwertig. Offen: $N=11\ldots15$
      (Zielsystem). Test: qualitativ, Polarisationsunabhängigkeit $10^{-9}$ \\
V11 & Fig.\ VI.10 (Al, $40^\circ$, $\nu=1{,}3+7{,}6i$): Li-Lamellentreppe vs.\ FFF & qualitativ: Treppe mit
      Li-Formel konvergiert langsam, FFF schnell \\
V13 & Gold-Sinus TM: $d=1$, $a=0{,}25\,\mu$m, $\nu=0{,}2+3{,}4i$, $20^\circ$ &
      C-Methode: $e_{-1}=0{,}26614$, $e_0=0{,}02362$; \verb|FFF_rig| $10^{-3}$, \verb|FFF_sym| $3\cdot10^{-3}$;
      Buchvariante versagt (dokumentiert) \\
V14 & Gold-Dreieck TM (Spitze bei $0{,}7\,\mu$m, sonst wie V13) & offen: nur physikalische Grenzen
      ($\sum e<1$); Laurent bei $N\le64$ noch nicht konvergiert \\
V12 & Nahfeld: Airy-Schicht analytisch; Stetigkeit an allen Element-/Scheibengrenzen (RK, RCWA, TE/TM);
      Flusskonstanz (verlustfreie Lamellen); RK-Feld = RCWA-Feld; $D_x$-Stetigkeit direkt vs.\ FFF &
      $10^{-12}$; $10^{-6}$; $10^{-12}$/$10^{-8}$; $10^{-7}$; FFF $<0{,}1$, direkt $>1$.
      2D (M9): $z$-invariante Strukturen = 1D-Nahfeld ($10^{-9}$, Karten = Linie $\cdot e^{i\gamma_0(z-z_0)}$),
      Stetigkeit (RK, RCWA) $10^{-6}$, Flusskonstanz $10^{-12}$, RK = RCWA $10^{-7}$, $D_x$-Sprung an der
      Säulenwand direkt 293\,\%, FFF 5\,\% / 3\,\% (anisotrop) \\
\bottomrule
\end{tabular}
\end{center}
\enquote{Intern} = in der Prüfung der Buchkapitel mit unabhängigen Python-Implementierungen berechnet
(nicht aus der Literatur); bei Abweichungen zuerst die eigene Implementierung, dann diese Werte prüfen.
Zusätzlich für jeden Referenzfall eine Konvergenzstudie in $N$ und (RCWA) in der Scheibenzahl als
Testausgabe (CSV) ablegen.

%=====================================================================
\section{Meilensteine}\label{sec:milestones}
%=====================================================================
\begin{longtable}{p{1.0cm}p{6.6cm}p{6.6cm}}
\toprule
M & Inhalt & Abnahme \\
\midrule
\endhead
M0 & Build-System, Abhängigkeiten, CI (CPU), Grundgerüst, Doku & \verb|cmake --build| + \verb|ctest| grün auf AlmaLinux 10 und in CI \\
M1 & IndexMap, Fourierkoeffizienten, Toeplitz, Butcher, S-Matrix-Kern, isotrope Modenbasis, $M$-Aufbau & Unit-Tests 1--6, 8 (isotrop) \\
M2 & 1D in-plane isotrop: homogene Schichten, Lamellen (RCWA, Li und FFF\_iso), Gesamtlösung, Effizienzen & V1, V2, V3 (TE/TM) \\
M3 & 1D Profile: Nullstellen-Schnitte, Normalenfeld, FFF\_iso; RK-Löser (Gauß 2/3, Schrittsteuerung); RCWA-Treppe mit FFF & V4, V5, V6; RK = RCWA auf Lamellen \\
M4 & 1D konisch ($4n_f$): volle $M$, Polarisation, s/p-Ausgabe & V1/V2 konisch, V3 konisch, Reziprozität \\
M5 & 1D anisotrop: Tensoren, FFF\_A, FFF\_B, anisotrope Substrate (4×4-Moden, Klassifikation) & V7, V8, Unit-Test 7 \\
M6 & Nahfeld 1D: interne Amplituden, RCWA/RK-Rekonstruktion, direkt/FFF, Schnitte, Plot-Skript & V12 (1D) \\
M7 & 2D isotrop: IndexMap 2D, Polygon-Koeffizienten, Normalenfeld (Alg.\ NF), RCWA; RK für Profile2D & V9, V10
  (erledigt v0.6; V10 nur qualitativ, quantitativ offen) \\
M8 & 2D anisotrop & V8 (2D), Konvergenzstudie (erledigt v0.7) \\
M9 & Nahfeld 2D (horizontal $x$--$z$, vertikal $x$--$y$) & V12 (2D) (erledigt v0.8) \\
M10 & CUDA-Backend (cuBLAS, cuSOLVER getrf/geev, cuFFT), Batching über Wellenlängen, Benchmarks & GPU = CPU auf $10^{-10}$; Benchmark-Bericht
  (v0.9: implementiert, gebaut und gelinkt; GPU-Abnahme auf dem Zielsystem offen) \\
M11 & Parameter-Scans, Performance, Doku für Anwender & Beispielskripte, Laufzeittabelle
  (erledigt v0.10: Scans $\lambda\times\varphi\times\psi\times\delta$, Sweeps per JSON-Pointer, paralleler Scan,
  \verb|plot_efficiencies.py|, \verb|docs/user_guide.md|) \\
\bottomrule
\end{longtable}

%=====================================================================
\section{Buchfehler mit Relevanz für die Implementierung}\label{sec:errata}
%=====================================================================
\begin{center}
\small
\begin{tabular}{p{2.3cm}p{5.4cm}p{6.6cm}}
\toprule
Stelle & Buch & korrekt (in diesem Plan verwendet) \\
\midrule
(II.19)/(II.20) & $\beta_0^{(1)}$ im einfallenden Term & $\beta_0^{(2)}$ \\
TM-Matrix S.\,46 & Indizes $(m,n)$ vertauscht & $M_{12,nm}=(k^2)_{n-m}$ usw. \\
(III.24)/(III.26) & $C$ und $C'$ vertauscht & Grenzflächen über $\Psi_{\text{ref}}^{-1}\Psi_X$ \\
S.\,56 & $A'^{(1)}=A^{(1)}$ & nur für $\nu_1=\nu_M$; Grenzfläche Substrat$\to$ref nötig \\
(III.29) & $T_{12}$, $T_{21}$ vertauscht & -- (Sternprodukt verwenden) \\
(III.45) & $S_{21}$ statt $S_{22}$ & \eqref{eq:star} \\
(IV.39)/(IV.40) & $\delta_{n,p}$ fehlt & Startwerte als Einheitsmatrix (Abschn.\ RK) \\
konisch & $\beta$ ohne $\gamma_0$ & $\beta=\sqrt{\varepsilon-\alpha^2-\gamma^2}$ \\
S.\,102 (Text) & Spaltenreihenfolge von $C_\varepsilon$ & \eqref{eq:Ceps} \\
S.\,108 & $C_\varepsilon$ $y$-unabhängig & hängt über $\varepsilon(y)$ von $y$ ab \\
(VI.13) & $L$ in allen Blöcken & nur Diagonalblöcke (schräge Lamellen) \\
(VII.72) & Vorfaktor $\sqrt{\mu_0/\varepsilon_0}/(ik_j)$ & $1/(i\omega\varepsilon_j)$ (nicht benötigt) \\
(VIII.2) & Vorzeichen $C_{23}$ & \eqref{eq:C1d} \\
(VIII.3) & $M_{32}$ ohne $-i\alpha^2/(\omega\mu_0)$ & \eqref{eq:Mmu1b} \\
(IX.1) & $\beta_0=k_0\cos\Psi$ & $\cos\Phi$, \eqref{eq:incidence} \\
(IX.14) & Division durch $N+1$ & \eqref{eq:indexmap} \\
(IX.16)--(IX.19) & Indizes $m(s)$, $\alpha_{m(r)}$ & alle zur Komponente $r$, $\alpha_{n(r)}$ \\
(IX.20) & Transmission ohne $k_2\to k_1$ & \eqref{eq:eff_closed} \\
Anh.\ IV & 2D-Blockgröße $(4N+2)^2$ & $2(2N+1)^2$ \\
(IV.32) & Vertauschung $\tp{N_i}\leftrightarrow\Li$ (Näherung) & versagt bei schwach absorbierenden Metallen;
  \eqref{eq:Qrig}/\eqref{eq:Qsym} \\
\bottomrule
\end{tabular}
\end{center}

%=====================================================================
\section*{Literatur}
\addcontentsline{toc}{section}{Literatur}
\begin{enumerate}[label={[\arabic*]},nosep,leftmargin=2em]
  \item M.~Nevière, E.~Popov, \emph{Light Propagation in Periodic Media: Differential Theory and Design}, Marcel Dekker, 2003.
  \item L.~Li, \enquote{Use of Fourier series in the analysis of discontinuous periodic structures}, JOSA A 13, 1870 (1996).
  \item L.~Li, \enquote{Formulation and comparison of two recursive matrix algorithms for modeling layered diffraction gratings}, JOSA A 13, 1024 (1996).
  \item L.~Li, \enquote{New formulation of the Fourier modal method for crossed surface-relief gratings}, JOSA A 14, 2758 (1997).
  \item E.~Popov, M.~Nevière, \enquote{Maxwell equations in Fourier space: fast-converging formulation for diffraction by arbitrary shaped, periodic, anisotropic media}, JOSA A 18, 2886 (2001).
  \item K.~Watanabe, R.~Petit, M.~Nevière, \enquote{Differential theory of gratings made of anisotropic materials}, JOSA A 19, 325 (2002).
  \item T.~Schuster, J.~Ruoff, N.~Kerwien, S.~Rafler, W.~Osten, \enquote{Normal vector method for convergence improvement using the RCWA for crossed gratings}, JOSA A 24, 2880 (2007).
  \item P.~Götz, T.~Schuster, K.~Frenner, S.~Rafler, W.~Osten, \enquote{Normal vector method for the RCWA with automated vector field generation}, Opt.\ Express 16, 17295 (2008).
  \item E.~Hairer, G.~Wanner, \emph{Solving Ordinary Differential Equations II}, Springer (Butcher-Tableaus Gauß/Radau).
  \item NVIDIA cuSOLVER Documentation, \verb|cusolverDnXgeev| (seit CUDA 12.6 Update 2).
\end{enumerate}
(Literaturangaben 7, 8 vor Verwendung bibliographisch prüfen.)

{\sloppy
%=====================================================================
\clearpage
\section{Erweiterungen ab Version 2.0: Überblick}\label{sec:ext}
%=====================================================================
Die Abschnitte \ref{sec:ext}--\ref{sec:extcross} spezifizieren die zweite Ausbaustufe (Meilensteine M12--M19).
Grundlage ist die Verbesserungsliste des Projekts:
CAD-Slicer für NURBS-Geometrie aus STEP-Dateien (über die eigene Bibliothek Nurbforge),
Python-Anbindung als Wheel, Differenzierbarkeit für Optimierung, mehr Parallelität, wählbare einfache
Genauigkeit (performant auch auf einer RTX~3090), PML für isolierte Strukturen und eine nicht gleichmäßige
Fourierdarstellung für metallische Strukturen. Es gelten alle Regeln der Abschnitte
\ref{sec:scope}--\ref{sec:errata} weiter, soweit hier nichts anderes steht.

\subsection{Annahmen und Festlegungen}
\begin{enumerate}[nosep]
  \item \enquote{PML für iso Strukturen} wird als \emph{isolierte}, also aperiodische Strukturen gelesen
        (Abschnitt~\ref{sec:pml}). Sollte \enquote{isotrop} gemeint sein, entfällt M18 in dieser Form.
  \item \enquote{Nicht gleichmäßige Fouriertransformation für Metalle} wird als adaptive räumliche Auflösung
        (ASR, Koordinatentransformation $x=f(X)$) umgesetzt (Abschnitt~\ref{sec:ctrans}). Das ist die in der
        FMM-Literatur etablierte Form; eine NUFFT im numerischen Sinn löst das Konvergenzproblem an
        Metallkanten nicht.
  \item PML und ASR verwenden \emph{dieselbe} Maschinerie (Koordinatentransformation im Fourierraum);
        M17 ist deshalb Voraussetzung für M18.
  \item Nurbforge (\url{https://github.com/PhotonPower/Nurbforge}, MIT, C++20) wird optional eingebunden
        (\verb|FG_ENABLE_NURBFORGE|). Der Kern von FourierGrating bleibt ohne Nurbforge baubar.
  \item Standard bleibt \verb|std::complex<double>|. Einfache Genauigkeit wird nur innerhalb des
        Linear-Algebra-Backends verwendet (M16).
  \item Wenn eine neue Option ausgeschaltet ist, rechnen alle bestehenden Pfade \textbf{bitgleich} wie in
        Version~1.9. Das wird mit Golden Files geprüft (Abschnitt~\ref{sec:extcross}).
\end{enumerate}

\subsection{Meilensteine und Abhängigkeiten}
\begin{center}
\small
\begin{tabular}{p{1.0cm}p{6.4cm}p{3.0cm}p{4.0cm}}
\toprule
M & Inhalt & setzt voraus & prüfbar in der Sandbox \\
\midrule
M12 & Python-Anbindung (pybind11, Wheel), \verb|Simulation|-Objekt & M11 & ja \\
M13 & Gekrümmte Querschnitte: Konturkoeffizienten, \verb|Contour2D|, Normalenfeld mit 3D-Normalen,
      höhenabhängige Schichtmodelle & M9 & ja \\
M14 & CAD-Slicer: STEP $\to$ Schichten (Nurbforge, Bandzerlegung, Materialzuordnung) & M13 (+ Nurbforge 0.6) & ja \\
M15 & Profiling und Parallelisierung (Scheiben, Baumreduktion, Koeffizienten, GPU-residente Ketten) & M10, M14 & Korrektheit ja, Leistung nein \\
M16 & Wählbare Genauigkeit FP64 / gemischt / FP32 (RTX 3090, RTX 6000 Ada) & M15 & CPU-Pfad ja, GPU nein \\
M17 & Koordinatentransformationen, adaptive räumliche Auflösung (Metalle) & M13 & ja \\
M18 & PML, isolierte Strukturen, strukturierte Ein-/Ausgangsmedien & M17 & ja \\
M19 & Differenzierbarkeit (diskrete Adjungierte), PyTorch/JAX-Adapter, Optimierungsbeispiele & M12, M13 (besser M14) & ja \\
\bottomrule
\end{tabular}
\end{center}
Die Meilensteine werden in dieser Reihenfolge abgearbeitet. Die offene GPU-Abnahme von M10 blockiert
M12--M14 und M17--M19 nicht; M15 (GPU-Teil) und M16 (GPU-Teil) werden erst auf einem Rechner mit GPU
abgenommen. Die Sandbox hat nur einen Kern und keine GPU: Testgrößen klein halten, Leistungsmessungen
gehören auf das Zielsystem (AlmaLinux 10, RTX 6000 Ada) bzw.\ den Windows-Rechner mit RTX~3090.

%=====================================================================
\section{M12: Python-Anbindung}\label{sec:python}
%=====================================================================

\subsection{Ziel}
\verb|pip install fouriergrating| (bzw.\ \verb|pip install .| im Repository) liefert ein Modul, das alles kann,
was \verb|fgsolve| kann, ohne Umweg über Dateien, und das für Optimierungsschleifen (M19) geeignet ist.

\subsection{Build}
\begin{itemize}[nosep]
  \item pybind11 ($\ge2.12$) und scikit-build-core ($\ge0.10$), wie in Nurbforge (gleiche Struktur,
        damit beide Projekte gleich gepflegt werden). CMake-Option \verb|FG_BUILD_PYTHON| (Standard \verb|OFF|;
        scikit-build setzt sie auf \verb|ON|).
  \item Verzeichnisse: \verb|python/bindings/| (C++-Modul \verb|fouriergrating._core|),
        \verb|python/fouriergrating/| (Python-Paket), \verb|tests/python/| (pytest), \verb|pyproject.toml| im Wurzelverzeichnis.
  \item Wheels per cibuildwheel: \verb|manylinux_2_28| x86\_64 (installiert auf AlmaLinux 8--10), CPython 3.10--3.13,
        CPU-Variante mit OpenMP (\verb|libgomp| bündelt auditwheel). Workflow \verb|.github/workflows/wheels.yml| für
        \verb|v*|-Tags, Artefakte wie bei Nurbforge. Windows-Wheel: optional, später (der Windows-Build nutzt bisher MinGW;
        MSVC-Kompatibilität ist vorher zu prüfen).
  \item CUDA: kein CUDA im CI-Wheel. Lokaler Build auf dem Zielsystem (im Benutzerhandbuch dokumentieren):
{\footnotesize\begin{verbatim}
pip install . -C cmake.define.FG_ENABLE_CUDA=ON -C cmake.define.CMAKE_CUDA_ARCHITECTURES="86;89"
\end{verbatim}}
\end{itemize}

\subsection{API}
Die JSON-Konfiguration bleibt die kanonische Beschreibung: Ein Python-\verb|dict| mit demselben Schema wie
\verb|docs/config_format.md| wird über nlohmann::json an \verb|io::loadConfig| übergeben. Damit gibt es genau
eine Semantik für CLI und Python.
{\footnotesize\begin{verbatim}
import fouriergrating as fg
cfg = fg.load_config("examples/v3_lamellar_dielectric.json")   # -> dict (Pfade relativ zur Datei aufgeloest)
res = fg.solve(cfg)                  # ein Punkt (erster Scanpunkt); -> fg.Result
res.reflected                        # numpy structured array: n, m, eff, eff_s, eff_p, amp_s, amp_p
res.transmitted, res.absorption, res.diagnostics   # dict: energy_balance, rk_steps, warnings, timings
scan = fg.solve_scan(cfg)            # Scan/Sweep aus cfg -> fg.ScanResult (Spalten als numpy-Arrays, to_pandas())
nf = fg.near_field(cfg)              # dict: Gitter + Felder als numpy-Arrays + meta (wie meta.json)
sim = fg.Simulation(cfg)             # Aufbau einmal (Schichtmodelle, Normalenfelder, IndexMap) ...
sim.solve(lambda_um=0.6328, phi_deg=30.0, psi_deg=0.0, delta_deg=90.0)   # ... viele Rechnungen
sim.set_parameter("/layers/0/thickness_um", 0.25)   # JSON-Pointer wie bei Sweeps; baut nur Betroffenes neu
fg.set_backend("auto", threshold=1000); fg.cuda_available(); fg.__version__
\end{verbatim}}
\begin{itemize}[nosep]
  \item Einheiten wie in JSON (µm, Grad); die Normierung bleibt intern.
  \item Während jeder Rechnung wird der GIL freigegeben (\verb|py::gil_scoped_release|); mehrere Python-Threads dürfen
        parallel rechnen (jeder mit eigener \verb|Simulation|).
  \item Ausnahmen: \verb|fg::Error| $\to$ \verb|fouriergrating.FourierGratingError| (Unterklasse von \verb|RuntimeError|),
        \verb|NotImplemented| $\to$ \verb|NotImplementedError|.
  \item \verb|Simulation| speichert die Schichtmodelle und baut beim Ändern eines Parameters nur die betroffenen
        Schichten neu (Grundlage für M19).
  \item Typannotationen (\verb|py.typed|, \verb|_core.pyi|).
\end{itemize}

\subsection{Tests und Abnahme}
\begin{itemize}[nosep]
  \item pytest: Für alle Dateien in \verb|examples/| (Nahfeld nur für die kleinen) stimmen Python-Ergebnisse mit
        \verb|fgsolve| (CSV) auf $10^{-12}$ relativ überein.
  \item GIL-Test: zwei Threads mit je einer \verb|Simulation| laufen gleichzeitig und liefern dieselben Werte wie sequentiell.
  \item \verb|set_parameter|: Ergebnis gleich einem Neuaufbau aus der geänderten Konfiguration (bitgleich).
  \item Wheel aus CI installiert in einer frischen venv; \verb|import fouriergrating; fg.solve(...)| läuft.
  \item Dokumentation \verb|docs/python.md|, Beispiel \verb|examples/python/scan_lamellar.py|.
\end{itemize}

%=====================================================================
\section{M13: Gekrümmte Querschnitte und höhenabhängige Schichtmodelle}\label{sec:contour}
%=====================================================================

\subsection{Ziel}
Querschnitte $y=\text{const}$ mit gekrümmten Rändern (Kreise, Ellipsen, NURBS) bekommen \emph{exakte}
Fourierkoeffizienten (bis auf Quadraturgenauigkeit $\approx10^{-14}$), ohne Polygonisierung und ohne Abtastung.
Außerdem werden Schichten möglich, deren Querschnitt und Normalenfeld sich mit $y$ ändern und die nicht als
Profil $y=g(x,z)$ darstellbar sind (Grundlage für den CAD-Slicer, M14). Das Modul ist von Nurbforge unabhängig.

\subsection{Konturdarstellung}
\begin{itemize}[nosep]
  \item \verb|fg::geometry::BezierSegment|: rationale Bézierkurve in der $(x,z)$-Ebene, Grad $p\le9$,
        Kontrollpunkte $\mathbf P_j$ und Gewichte $w_j>0$.
  \item \verb|fg::geometry::Contour2D|: geschlossene Kette von Segmenten (Ende = Anfang des nächsten, Toleranz
        $10^{-12}\,d$). Orientierung gegen den Uhrzeigersinn in $(x,z)$ heißt: Material links (wie \eqref{eq:polygon}).
  \item \verb|fg::geometry::Region|: eine äußere Kontur, beliebig viele Löcher, Materialname.
        Konstruktoren: Polygon, Kreis und Ellipse (exakt als rationale Quadratik, vier Segmente),
        NURBS (Grad, Knoten, Kontrollpunkte, Gewichte; Zerlegung in Bézierstücke per Knoteneinfügung, Piegl--Tiller A5.6).
\end{itemize}

\subsection{Fourierkoeffizienten aus Konturintegralen}
Für ein Gebiet $\Omega$ mit Rand $\partial\Omega$ (äußere Kontur gegen, Löcher im Uhrzeigersinn) und
$\mathbf k=(pK_x,qK_z)\neq0$ folgt wie bei \eqref{eq:polygon} aus dem Gaußschen Satz
\begin{equation}
  f_{pq}=\frac{1}{d_xd_z}\iint_\Omega e^{-i\mathbf k\cdot\mathbf r}\,d^2r
  =\frac{i}{|\mathbf k|^2d_xd_z}\sum_{\text{Segmente}}\int_0^1\big(k_x\,z'(t)-k_z\,x'(t)\big)\,
  e^{-i(k_xx(t)+k_zz(t))}\,dt,
  \label{eq:contour}
\end{equation}
denn für eine gegen den Uhrzeigersinn laufende Kurve $\mathbf r(t)=(x(t),z(t))$ ist $\mathbf n\,ds=(z'(t),-x'(t))\,dt$.
Für $\mathbf k=0$:
\begin{equation}
  f_{00}=\frac{1}{d_xd_z}\cdot\frac12\sum_{\text{Segmente}}\int_0^1\big(x(t)z'(t)-z(t)x'(t)\big)\,dt .
  \label{eq:contour0}
\end{equation}
\textbf{Prüfung:} Für Grad~1 ergibt \eqref{eq:contour} exakt \eqref{eq:polygon} (die Integration über eine Strecke
liefert $\ell_je^{-i\mathbf k\cdot\mathbf c_j}\operatorname{sinc}(\mathbf k\cdot\mathbf t_j/2)$).
Für einen Kreis mit Radius $R$ und Mittelpunkt $\mathbf c$ gilt analytisch
\begin{equation}
  f_{pq}=\frac{2\pi R\,J_1(|\mathbf k|R)}{|\mathbf k|\,d_xd_z}\,e^{-i\mathbf k\cdot\mathbf c},\qquad f_{00}=\frac{\pi R^2}{d_xd_z}.
  \label{eq:circle}
\end{equation}
\textbf{Periodische Bilder:} Da $e^{-i\mathbf k\cdot\mathbf r}$ für alle $\mathbf k$ des reziproken Gitters periodisch
ist, darf ein Gebiet über den Zellrand ragen; \eqref{eq:contour} liefert dann den Koeffizienten der periodisierten
Indikatorfunktion, solange sich die periodischen Bilder nicht überlappen (wie bei \verb|Polygon2D|).

\paragraph{Quadratur.} Je Bézierstück Gauß--Legendre. Der Integrand ist glatt, oszilliert aber mit der Phase
$\Phi=|\mathbf k|_{\max}L$ ($L$ = Bogenlänge des Stücks, $|\mathbf k|_{\max}$ für $|p|\le2N_x$, $|q|\le2N_z$).
Adaptiv: Stück mit $n$ und $2n$ Punkten auswerten ($n_0=\max(16,\lceil\Phi/2\rceil+8)$); weichen die Beiträge um
mehr als $10^{-14}\max_{pq}|f_{pq}|\,d_xd_z$ ab, Stück halbieren (de Casteljau, rational im homogenen Raum).
Die Polynomordnung ändert sich bei rationalen Stücken nicht; die Gewichte gehen nur über $x(t),z(t)$ ein.

\paragraph{Effiziente Auswertung aller $(p,q)$.} Mit den Quadraturknoten $t_\ell$, Gewichten $\omega_\ell$ und
$E^x_{\ell p}=e^{-ipK_xx(t_\ell)}$, $E^z_{\ell q}=e^{-iqK_zz(t_\ell)}$ (Rekursion über Potenzen) gilt
\[
  G^{(1)}=\sum_\ell\omega_\ell z'(t_\ell)\,E^x_{\ell\cdot}{}^{\!T}E^z_{\ell\cdot},\qquad
  G^{(2)}=\sum_\ell\omega_\ell x'(t_\ell)\,E^x_{\ell\cdot}{}^{\!T}E^z_{\ell\cdot},\qquad
  f_{pq}=\frac{i\,(k_xG^{(1)}_{pq}-k_zG^{(2)}_{pq})}{|\mathbf k|^2d_xd_z},
\]
also zwei Matrixprodukte $(E^x)^T\operatorname{diag}(\cdot)E^z$. Ergebnis in \verb|fourier::CoefficientGrid|.
Implementierung: \verb|fourier::regionCoefficients| (Argumente: Region, Perioden, $P_x$, $P_z$) in
\verb|src/fourier/contour.cpp|; Formelverweise \verb|// Plan eq:contour|.

\subsection{Neuer Schichttyp \texttt{Contour2D}}
\verb|Contour2DLayer| ist die Verallgemeinerung von \verb|Polygon2DLayer|: $y$-invariant, Hintergrundmaterial,
Liste von \verb|Region|s (paarweise disjunkt; geprüft über Kurvenschnitte der Ränder und Enthaltensein eines Punktes).
Koeffizienten: $\varepsilon=\varepsilon_{\text{bg}}+\sum_r(\varepsilon_r-\varepsilon_{\text{bg}})\chi_r$ (Linearität;
für $1/\varepsilon$ und Tensoren analog wie bei \verb|Polygon2D|). Faktorisierungen wie \verb|Polygon2D|
(Laurent, FFF\_iso/FFF\_sym mit Tensorfeld, FFF\_B; FFF\_rig nicht zulässig).
JSON:
{\footnotesize\begin{verbatim}
{"type": "contour2d", "thickness_um": 0.2, "background": "air",
 "regions": [ {"material": "Au", "circle": {"center_um": [0.25, 0.25], "radius_um": 0.1}},
              {"material": "Au", "ellipse": {"center_um": [0.6, 0.6], "radii_um": [0.12, 0.06], "angle_deg": 30}},
              {"material": "SiO2", "polygon_um": [[0,0],[0.1,0],[0.1,0.1]]},
              {"material": "Si", "nurbs": {"degree": 2, "knots": [...], "points_um": [[x,z],...],
                                           "weights": [...]}, "holes": [ ... ]} ]}
\end{verbatim}}

\subsection{Normalenfeld mit Randstichproben und 3D-Normalen}
Algorithmus~NF (Abschnitt~\ref{sec:fff}) wird so verallgemeinert, dass er statt Polygonkanten eine Liste von
\emph{Randstichproben} $(\mathbf p_s,\mathbf n_s)$ verarbeitet, mit $\mathbf n_s$ als 3D-Einheitsvektor
(bei $y$-invarianten Schichten $n_{s,y}=0$, bei geneigten Wänden aus M14 $n_{s,y}\neq0$):
\begin{enumerate}[nosep]
  \item Kurven adaptiv abtasten, bis die Sehnenabweichung $\le10^{-3}$ Rasterabstände beträgt; $\mathbf n_s$ exakt aus
        der Kurventangente (2D) bzw.\ von außen geliefert (3D, M14). Die Prüfung \enquote{Materialgrenze ja/nein}
        bleibt wie in Algorithmus~NF (Material beidseits verschieden, periodische Punktsuche).
  \item Nächste Stichprobe je Rasterpunkt mit Bucket-Gitter (Zellgröße $\approx$ Stichprobenabstand), periodische Bilder.
        Anforderung: $512^2$ Raster mit $10^4$ Stichproben in unter 1\,s auf einem Kern.
        Zwischen zwei benachbarten Stichproben wird der Fußpunkt auf der Sehne bestimmt und die Normale linear
        interpoliert und normiert.
  \item Tensorvariante: geglättet wird das $3\times3$-Projektorfeld $\mathbf n\mathbf n^T$; Koeffizienten aller sechs
        Komponenten $N_iN_j$. Für $n_{s,y}=0$ ist das Ergebnis bitgleich mit der bisherigen $2\times2$-Variante
        (Regressionstest gegen \verb|Polygon2D|).
\end{enumerate}

\subsection{Höhenabhängige Schichtmodelle}
Die interne Schnittstelle \verb|detail::LayerModel| (Datei \verb|src/solvers/layer_model.hpp|) wird erweitert:
{\footnotesize\begin{verbatim}
virtual Vec3D normalAt(double x_um, double y_um, double z_um) const;          // bisher (x, z)
virtual Eigen::Matrix3d projectorAt(double x_um, double y_um, double z_um) const;
virtual const NormalFieldCoefficients& NN(double y_um) const;                // bisher ohne y
\end{verbatim}}
$y$-invariante Schichten und Profile ignorieren $y$ (Profile: Normale überträgt sich, Abschnitt~\ref{sec:fff});
alle Aufrufer (Q-Matrizen, FFF-konsistente Rekonstruktion, Nahfeld) übergeben die Höhe. Neues Modell
\verb|detail::SectionedLayer| für Schichten, deren Querschnitt sich mit $y$ ändert:
{\footnotesize\begin{verbatim}
struct CrossSection { std::vector<Region> regions; std::vector<BoundarySample> samples; };
class SectionProvider {                       // M14: Nurbforge; Tests: analytische Provider
 public: virtual CrossSection section(double y_um) const = 0;   // y relativ zur Schichtunterkante
         virtual double thickness_um() const = 0; };
\end{verbatim}}
\verb|SectionedLayer| berechnet bei Bedarf Querschnitt, Koeffizienten von $\varepsilon$, $1/\varepsilon$, $N_iN_j$
und $Q(y)$ und speichert sie in einem Cache (Schlüssel: $y$ bitgenau; dieselben Knoten werden bei der Nahfeldrekonstruktion
wieder gebraucht, Abschnitt~\ref{sec:layers}). RK wertet an den Stufenknoten aus; RCWA verwendet die Scheibenmitten
(Treppe, \verb|slices|). Der Kontaminationsschutz \eqref{eq:slicecrit} gilt unverändert.

\paragraph{Analytischer Testprovider.} Für die Abnahme ohne Nurbforge: Kegelstumpf (Kreis mit Radius
$R(y)=R_0+(R_1-R_0)y/h$, 3D-Normale $\propto(\cos\theta,\,-(R_1-R_0)/h,\,\sin\theta)$ mit dem Umfangswinkel $\theta$)
und Kugelkappe. Beide sind auch als \verb|Profile2D| mit Funktion $g(x,z)$ darstellbar (Referenz).

\subsection{Tests und Abnahme (M13)}
\begin{itemize}[nosep]
  \item Unit: Grad-1-Konturen $=$ \eqref{eq:polygon} auf $10^{-14}$; Kreis $=$ \eqref{eq:circle} auf $10^{-13}$ (bis
        $|p|,|q|\le60$); Ellipse $=$ gestreckter Kreis; Verschiebung um eine Periode ändert nichts ($10^{-13}$);
        Loch $=$ negative Fläche; Adaptivität (Ergebnis unabhängig von $n_0$ auf $10^{-14}$).
  \item Unit: NF mit Randstichproben auf Polygonen bitgleich zur bisherigen Implementierung; 3D-Normalen mit $n_y=0$
        bitgleich; Projektor symmetrisch, Spur 1.
  \item \textbf{V15} (Säulen mit Kreisquerschnitt): \verb|Contour2D| mit exaktem Kreis gegen \verb|Polygon2D| mit $L$ Ecken
        ($L=16,32,\dots,512$): Abweichung der Effizienzen fällt wie $O(L^{-2})$, Richardson-Extrapolation stimmt mit
        \verb|Contour2D| auf $10^{-8}$ überein. RK $=$ RCWA ($y$-invariant).
  \item \textbf{V16a} (geneigte Wand ohne CAD): \verb|SectionedLayer| mit Kegelstumpf-Provider gegen \verb|Profile2D|
        (gleiche Geometrie, RK, gleiche Faktorisierung): $10^{-6}$ bei \verb|rk.tol| $=10^{-9}$; Energieerhaltung bei
        verlustfreiem Material $10^{-10}$.
  \item Alle bisherigen Tests bitgleich (Golden Files).
\end{itemize}

%=====================================================================
\section{M14: CAD-Slicer (STEP $\to$ Schichten) mit Nurbforge}\label{sec:cad}
%=====================================================================

\subsection{Ziel und Datenfluss}
Eine STEP-Datei (B-rep, NURBS und analytische Flächen) beschreibt den Inhalt einer Einheitszelle. FourierGrating
zerlegt sie entlang $y$ in Bänder, ordnet Materialien zu und erzeugt daraus homogene, $y$-invariante
(\verb|Contour2D|) und $y$-variable (\verb|SectionedLayer|) Schichten.
{\footnotesize\begin{verbatim}
STEP --nf::read_step--> Shells (+ Namen) --Achsen/Einheiten--> FG-Koordinaten (um)
  --nf::critical_heights--> Baender [y_k, y_k+1]
  --je Band: homogen | y-invariant (1 Schnitt) | y-variabel (Schnitte bei Bedarf)
  --nf::section--> Schleifen + Flaechennormalen --> Region/Contour2D + Randstichproben --> M13
\end{verbatim}}

\subsection{Vorleistung in Nurbforge (Version 0.6.0, eigener Pull Request)}
Diese Funktionen gehören fachlich in den CAD-Kern und werden \emph{dort} implementiert und getestet
(Repository \verb|PhotonPower/Nurbforge|, Regeln dort in \verb|CONTRIBUTING.md|). FourierGrating bindet danach das
Tag \verb|v0.6.0| ein.
\begin{enumerate}[nosep]
  \item \verb|nf::PlanarSection nf::section(const Shell&, const Plane&, const SectionOptions&)|:
        geschlossene Schleifen in den Ebenenkoordinaten $(s,t)$, je Schleife eine Kette von \verb|NurbsCurve|
        (z = 0). Orientierung: Festkörper links (äußere Schleifen gegen den Uhrzeigersinn, Hohlräume im Uhrzeigersinn).
        Je Kurve: Index der Quellfläche und die p-Kurve, damit die nach außen zeigende Flächennormale an jedem
        Schnittpunkt ausgewertet werden kann. Umsetzung: \verb|planar_section| je Fläche, Kurven auf den
        Trimmbereich beschneiden (\verb|TrimmedSurface::contains|, Schnitt mit den Trimmkurven), Stücke an den Enden
        verketten (Toleranz der Shell). Offene Ketten $\Rightarrow$ \verb|GeometryError| mit Ort.
        Ebenen, die eine Fläche tangential schneiden, sind nicht erlaubt (der Aufrufer wählt keine kritischen Höhen).
  \item \verb|std::vector<double> nf::critical_heights(const Shell&, const Vec3& dir, const Tolerance&)|:
        sortierte Höhen $h=\mathbf r\cdot\mathbf d$ von Ecken, ebenen Flächen mit $\mathbf n\parallel\mathbf d$, inneren
        Extrema der Höhe auf Flächen ($S_u\cdot\mathbf d=S_v\cdot\mathbf d=0$, Newton aus Bézier-Unterteilung) und auf
        Kanten ($C'\cdot\mathbf d=0$). Zwischen zwei kritischen Höhen ändert sich die Topologie der Schnitte nicht.
  \item \verb|bool nf::is_extrusion_band(const Shell&, const Vec3& dir, double h0, double h1, const Tolerance&)|:
        alle Flächen, die das Band schneiden, haben dort $\mathbf n\cdot\mathbf d=0$ (Prüfung an Stichproben und an
        den Bézier-Kontrollnetzen).
  \item STEP-Namen: \verb|read_step_named| liefert je Shell den Namen (PRODUCT bzw.\ MANIFOLD\_SOLID\_BREP-Name).
  \item Thread-Sicherheit aller \verb|const|-Abfragen dokumentieren und mit einem TSAN-Test prüfen (für M15).
  \item Tests in Nurbforge: Quader, Zylinder, Kugel, Kegel, Torus, Quader mit Hohlraum; Schleifenzahl, Orientierung,
        Fläche gegen analytische Werte ($10^{-9}$); kritische Höhen (Kugel: Pole; Torus: vier Höhen bei
        $\mathbf d\parallel$ Achse \dots).
\end{enumerate}

\subsection{Einbindung in FourierGrating}
\begin{itemize}[nosep]
  \item CMake-Option \verb|FG_ENABLE_NURBFORGE| (Standard \verb|OFF| im Kern-CI, eigener CI-Job mit \verb|ON|);
        \verb|FetchContent| auf das Tag \verb|v0.6.0| (\verb|NF_BUILD_PYTHON=OFF|, \verb|NF_BUILD_TESTS=OFF|), statisch gelinkt.
        Die Python-Wheels (M12) werden mit \verb|ON| gebaut. Lizenzhinweis (MIT) in \verb|README.md|.
  \item Neuer Namensraum \verb|fg::cad| (\verb|include/fg/cad|, \verb|src/cad|). Kein anderes Modul inkludiert Nurbforge-Header.
  \item Python-Nutzer von Nurbforge übergeben Geometrie als STEP-Text (\verb|nf.to_step(shells)|); C++-Objekte werden
        nicht zwischen den beiden Python-Modulen ausgetauscht (ABI-Risiko).
\end{itemize}

\subsection{Koordinaten, Einheiten, Einheitszelle}
\begin{itemize}[nosep]
  \item STEP wird in mm gelesen (Nurbforge rechnet aus der Dateieinheit um), dann mit 1000 auf µm skaliert;
        die Option \verb|unit| (\verb|auto|, \verb|mm|, \verb|um|) korrigiert falsche Dateieinheiten.
  \item Achsen: Standard CAD-$+Z\to$ FG-$+y$ (Stapelrichtung), CAD-$X\to x$; $z$ folgt aus Rechtshändigkeit
        $\hat{\mathbf z}=\hat{\mathbf x}\times\hat{\mathbf y}$, also FG-$z=-$CAD-$Y$. Konfigurierbar über
        \verb|"axes": {"y": "+Z", "x": "+X"}|; eine Spiegelung (Linkssystem) ist verboten, weil sie chirale Strukturen
        verfälscht. Test V16b prüft das.
  \item Zelle: Ursprung \verb|cell_origin_um| $=(x_0,z_0)$ (Standard: Minimum der Hüllquader-Projektion) und Perioden aus
        \verb|period_um|. Geometrie darf über die Zelle ragen (periodische Bilder, M13). Überlappen sich Bilder in
        einem Schnitt (Kurvenschnitt der Schleife mit ihren um $\pm d_x,\pm d_z$ verschobenen Kopien), bricht der Slicer
        mit einer Meldung ab; \verb|"periodic_union": true| vereinigt vorher in 3D mit den verschobenen Kopien
        (\verb|nf::boolean|, Union).
  \item $y$-Bereich: \verb|y_range_um| (Standard: Hüllquader). Der CAD-Block ist ein Eintrag in \verb|layers| und
        wird beim Laden in mehrere Schichten (Bänder) expandiert.
\end{itemize}

\subsection{Materialzuordnung}
\verb|"bodies"| ordnet Shells (Index oder STEP-Name) Materialien zu; nicht genannte Shells sind ein Fehler.
\verb|"background"| ist das Material außerhalb aller Körper (Name oder \verb|"superstrate"|/\verb|"substrate"|).
Körper müssen disjunkt sein. Überlappen sie, gilt die Reihenfolge der Liste als Priorität:
$B_i\leftarrow B_i\setminus\bigcup_{j<i}B_j$ (\verb|nf::boolean|, Difference) bei \verb|"overlap": "priority"|,
sonst Abbruch. Hohlräume (BREP\_WITH\_VOIDS) werden Löcher. Permittivität im Schnitt:
$\varepsilon=\varepsilon_{\text{bg}}+\sum_b(\varepsilon_b-\varepsilon_{\text{bg}})\chi_b$.

\subsection{Bandzerlegung und Klassifikation}
\begin{enumerate}[nosep]
  \item Kritische Höhen aller Shells im $y$-Bereich bestimmen und zusammenführen; Höhen mit Abstand
        $<$ \verb|critical_tol_um| (Standard $10^{-6}$) werden zusammengelegt. Bänder dünner als
        \verb|min_band_um| (Standard $10^{-5}$) werden dem Nachbarband zugeschlagen (Warnung).
  \item Je Band: kein Rand im Band $\Rightarrow$ \verb|Homogeneous| (Material per \verb|point_in_shell| in Bandmitte,
        Zellmitte); \verb|is_extrusion_band| $\Rightarrow$ \verb|Contour2D| aus \emph{einem} Schnitt in Bandmitte
        (exakt); sonst \verb|SectionedLayer| mit \verb|NurbforgeSectionProvider|.
  \item Es wird nie genau auf einer kritischen Höhe geschnitten. RK-Stufenknoten liegen im Inneren; RCWA-Scheibenmitten
        ebenfalls. Bandgrenzen sind Schichtgrenzen des Stapels.
\end{enumerate}

\subsection{Schnitt $\to$ Regionen und Randstichproben}
Schleifen aus \verb|nf::section| werden in Bézierstücke zerlegt (\verb|nf| Bézier-Zerlegung), in FG-Koordinaten
abgebildet und zu \verb|Region|s geordnet (äußere Schleife = positive Fläche; Löcher der kleinsten umgebenden
äußeren Schleife zugeordnet). Randstichproben: Punkte auf den Schleifen (Abstand wie M13) mit der 3D-Flächennormale
der Quellfläche, umgerechnet in FG-Achsen. Für das Normalenfeld genügt die Projektorform, das Vorzeichen ist egal.

\paragraph{Kosten und Interpolation in $y$.} Ein Schnitt kostet Flächen-Ebenen-Schnitte (Marching); RK fordert viele
Höhen an. Deshalb: (a) Cache je Höhe; (b) Option \verb|"y_interpolation": "chebyshev"|: je Band die
Koeffizientengitter von $\varepsilon$, $1/\varepsilon$ (bzw.\ Tensorgrößen) und $N_iN_j$ an $K$ Tschebyschow-Knoten
(Standard $K=16$) berechnen und baryzentrisch interpolieren. Innerhalb eines Bandes sind die Koeffizienten glatt in
$y$; an Bandenden können Wurzelsingularitäten auftreten (Kuppenspitze). Fehlerkontrolle: Vergleich $K$ gegen $2K$,
Ziel $10^{-10}\max|f_{pq}|$, sonst $K$ verdoppeln (bis 256, dann exakte Schnitte mit Warnung).

\subsection{Konfiguration}
{\footnotesize\begin{verbatim}
{"type": "cad", "file": "pillars.step", "unit": "auto", "axes": {"y": "+Z", "x": "+X"},
 "y_range_um": [0.0, 0.35], "cell_origin_um": [0.0, 0.0], "background": "superstrate",
 "bodies": [ {"name": "Pillar", "material": "Au"}, {"index": 1, "material": "SiO2"} ],
 "overlap": "error", "periodic_union": false,
 "slicing": {"critical_tol_um": 1e-6, "min_band_um": 1e-5, "rcwa_slices": 40,
             "y_interpolation": "none", "chebyshev_nodes": 16},
 "factorization": "auto", "normal_field": {"method": "tensor", "grid": 512, "sigma": 3}}
\end{verbatim}}
\verb|fgsolve --cad-report| schreibt die Bandzerlegung (Höhen, Typ je Band, Flächenanteile je Material, Warnungen)
als JSON; \verb|scripts/plot_cad_sections.py| zeichnet ausgewählte Schnitte.

\subsection{Tests und Abnahme (M14)}
STEP-Testdateien werden zur Testlaufzeit mit Nurbforge erzeugt (\verb|nf::write_step| in ein temporäres
Verzeichnis); zusätzlich zwei kleine Dateien unter \verb|tests/data/cad/| für Regressionstests.
\begin{itemize}[nosep]
  \item \textbf{V16b} (Prisma, Chiralität): extrudiertes L-Profil $\to$ ein $y$-invariantes Band; Koeffizienten
        $=$ \verb|Polygon2D| mit von Hand abgebildetem Polygon ($10^{-13}$); Effizienzen RK und RCWA $10^{-12}$.
        Gespiegelte Achsenwahl muss abgelehnt werden.
  \item \textbf{V17} (Zylindersäulen): STEP-Zylinder $=$ \verb|Contour2D| mit Kreis ($10^{-12}$), Koeffizienten gegen
        \eqref{eq:circle}.
  \item \textbf{V18} (Kugelkappe, Kegelstumpf, geneigte Wände): \verb|SectionedLayer| aus CAD gegen \verb|Profile2D| mit
        Funktion $g(x,z)$ (RK, \verb|rk.tol| $=10^{-9}$): $10^{-6}$; gegen den analytischen Provider aus M13: $10^{-9}$;
        RCWA-Treppe konvergiert mit der Scheibenzahl gegen denselben Wert; Energie $10^{-10}$ (verlustfrei).
  \item Zwei Materialien (Au-Säule mit SiO$_2$-Deckel) und Hohlraum: RK $=$ RCWA im Grenzwert, Energie, Reziprozität.
  \item Tschebyschow-Interpolation: Ergebnis gleich exakten Schnitten auf $10^{-9}$.
  \item Doku \verb|docs/cad.md|, Beispiel \verb|examples/cad_pillar_array.json|, Skript
        \verb|scripts/make_cad_examples.py| (erzeugt die STEP-Beispiele mit dem Nurbforge-Python-Paket).
\end{itemize}

%=====================================================================
\section{M15: Profiling und Parallelisierung}\label{sec:parallel}
%=====================================================================

\subsection{Ziel}
Eine \emph{einzelne} Rechnung (nicht nur ein Scan) nutzt alle Kerne bzw.\ die GPU gut aus. Ergebnisse sind
deterministisch und unabhängig von der Threadzahl.

\subsection{Schritt 1: Messen}
\begin{itemize}[nosep]
  \item \verb|fg::prof|: Scope-Timer (\verb|FG_PROF_SCOPE("rcwa.eig")|), threadsichere Aggregation (Summe, Anzahl, Max
        je Name), Ausgabe als JSON (\verb|fgsolve --profile|, Python \verb|res.diagnostics["timings"]|). Kosten
        abgeschaltet $<1\,\%$.
  \item Profile für V3, V9, V10 ($N=10$), V13 und das CAD-Beispiel auf dem Zielsystem (CPU und GPU) in
        \verb|docs/benchmarks.md| (\enquote{Profil vor M15}). Die folgenden Schritte werden nach diesem Profil priorisiert;
        was nicht messbar hilft, entfällt.
\end{itemize}

\subsection{Schritt 2: Threadmodell}
Drei Ebenen: Scanpunkte (bestehend, \verb|solveScan|), Schichten/Scheiben einer Rechnung (neu), BLAS/Eigen
innerhalb einer Matrixoperation. Konfiguration \verb|"threads"| mit den Schlüsseln \verb|scan|, \verb|slices| (je Zahl oder \verb|auto|) und \verb|blas| (Standard 1),
Umgebungsvariable \verb|FG_NUM_THREADS|. Regel gegen Überbelegung: Produkt der Ebenen $\le$ Kernzahl; \verb|auto|
verteilt zuerst auf die Scanebene (wenn Scanpunkte $\ge$ Kerne), sonst auf die Scheibenebene. OpenMP-Tasks;
verschachtelte Parallelität explizit (\verb|omp_set_max_active_levels(2)|), Eigen über \verb|Eigen::setNbThreads|.

\subsection{Schritt 3: Parallele Scheiben-S-Matrizen und Baumreduktion}
\begin{itemize}[nosep]
  \item Die S-Matrizen aller Schichten und Scheiben (RCWA-Scheiben, RK-Stabilisierungsscheiben) sind unabhängig und werden
        parallel berechnet; Ablage nach Index.
  \item Das Sternprodukt ist assoziativ. Verkettung als balancierter Binärbaum mit \emph{fester} Form (Paare $(0,1),(2,3),\dots$,
        dann rekursiv), unabhängig von der Threadzahl $\Rightarrow$ bitgleiche Ergebnisse für jede Threadzahl.
        Abweichung zur sequentiellen Verkettung $\le10^{-12}$ relativ (Test). Option \verb!"reduction": "tree" | "sequential"!.
  \item Nahfeld (\verb|storeInternal|) braucht $S_L$ und $S_U$ an jeder inneren Grenzfläche: zunächst sequentiell
        (Präfix/Suffix); parallele Präfixsumme nur, wenn das Profil es verlangt.
  \item RK: Stufenauswertungen $Q(y+c_ih)$ einer Scheibe parallel; Scheiben untereinander parallel.
\end{itemize}

\subsection{Schritt 4: Geometrie und Koeffizienten}
Raster von Algorithmus~NF, Konturquadratur (je Thread eigenes Teilgitter, Summation in fester Reihenfolge),
Profil-Schnittpunkte und CAD-Schnitte je Höhe parallel. Für Nurbforge nur die als threadsicher dokumentierten
\verb|const|-Abfragen verwenden (M14).

\subsection{Schritt 5: GPU-residente Ketten}
Das CUDA-Backend arbeitet bisher über die Host-API (jede Operation kopiert Ein- und Ausgaben). Bei FP32 (M16) wird der
PCIe-Transfer sonst zum Engpass (Beispiel $n=3362$: eine Matrix 90--180\,MB, CGEMM auf der RTX~3090 nur wenige ms).
Daher:
\begin{itemize}[nosep]
  \item \verb|linalg::DeviceMatrix| (RAII-Handle) und Operationen \verb|gemm|, \verb|solve|, \verb|eig|, \verb|add|,
        \verb|scaleRows| auf Handles; \verb|upload|/\verb|download| explizit.
  \item RCWA-Scheiben-S und Sternproduktkette laufen vollständig auf dem Gerät; nur das Gesamt-$S$ (bzw.\ die für das
        Nahfeld nötigen Matrizen) wird zurückkopiert. Mehrere Streams für unabhängige Scheiben.
  \item Kontrolle mit Nsight Systems: Anteil der Transferzeit $<10\,\%$ bei 2D, $N\ge10$.
\end{itemize}

\paragraph{Umsetzung (v2.7).} Die Baumform wird ohne vollständige Ablage aller Scheiben erreicht: ausgerichtete
Blöcke von $2^m\ge$ Threads Elementen werden parallel paarweise reduziert, vollständige Teilbäume wie ein
Binärzähler zusammengefasst und der Rest von rechts gefaltet; das ergibt dieselbe Form wie die stufenweise
Paarbildung für jede Blockgröße (Test bitgleich), Speicher $O(\text{Threads}+\log K)$. Standard bleibt
\verb!"reduction": "sequential"! (bitgleich zu v0.10); die parallele Berechnung der Scheiben selbst ändert keine Bits und
ist immer aktiv. \verb|auto| gibt der Scanebene $\min(\text{Scanpunkte},\text{Kerne})$ Threads und der Scheibenebene den
Rest. Innere Schleifen (Raster von Algorithmus~NF zeilenweise, Gauß-Glättung, Konturstücke je Segment,
Tschebyschow-Knoten) laufen nur parallel, wenn der Aufrufer selbst sequentiell ist; Summationsreihenfolgen bleiben
erhalten (bitgleich). Ohne OpenMP verteilt \verb|parallelFor| auf \verb|std::thread|; der TSAN-Job baut so, weil
libgomp nicht instrumentiert ist. MSVC: OpenMP~2.0 genügt (\verb|omp_set_nested| statt \verb|omp_set_max_active_levels|).
GPU: \verb|DeviceMatrix| mit Speicherpool je Backend (kein \verb|cudaFree| im Lauf); ein Stream je Thread, eine Matrix
wird erst nach \verb|synchronize()| des erzeugenden Threads weitergegeben, beim Freigeben durch einen anderen Thread
wartet dieser auf seinen Stream. Die residenten Operationen haben eine Host-Referenz (Eigen), gegen die der
Gerätepfad auch ohne GPU geprüft wird. Höchstens \verb|FG_GPU_STREAMS| (Standard 4) Scheiben gleichzeitig auf dem Gerät.

\subsection{Tests und Abnahme (M15)}
\begin{itemize}[nosep]
  \item Determinismus: Threadzahl 1, 2, 4 liefern bitgleiche Ergebnisse (alle Validierungsfälle, kleine $N$).
  \item Baum = sequentiell auf $10^{-12}$; GPU-residente Kette = Host-Kette auf $10^{-10}$ (GPU-Rechner).
  \item CI-Job mit ThreadSanitizer (Unit-Tests, kleine Größen) ohne Befund.
  \item Leistungsbericht in \verb|docs/benchmarks.md| (Zielsystem und RTX~3090): Laufzeit einer 2D-Rechnung mit
        $\ge32$ Scheiben vor/nach, parallele Effizienz (Zielwert $\ge60\,\%$ auf 16 Kernen), Transferanteil GPU.
\end{itemize}

%=====================================================================
\section{M16: Wählbare Genauigkeit (FP64 / gemischt / FP32)}\label{sec:precision}
%=====================================================================

\subsection{Motivation und Grundsatz}
Auf der RTX~3090 (GA102, \verb|sm_86|) und der RTX 6000 Ada (\verb|sm_89|) beträgt die FP64-Rate etwa 1/64 der
FP32-Rate. Die $O(n^3)$-Operationen (GEMM, LU, Eigenwerte) dominieren die Laufzeit. Die Genauigkeit wird deshalb
\textbf{am Backend} gewählt: Matrizen bleiben im Physikcode \verb|MatC| (FP64); das Backend wandelt an seiner Grenze
($O(n^2)$) um und rechnet die $O(n^3)$-Operation in der gewählten Genauigkeit. Der Physikcode wird nicht templatisiert.

\subsection{Modi}
\begin{center}
\small
\begin{tabular}{p{4.6cm}p{2.6cm}p{3.6cm}p{3.0cm}}
\toprule
Operation & \verb|fp64| (Standard) & \verb|mixed| & \verb|fp32| \\
\midrule
Toeplitz, $Q$, $M$, Koeffizienten, Normalenfeld & FP64 & FP64 & FP64 \\
GEMM & Z & C & C \\
LU / Lösen & Z & C-Faktorisierung + FP64-Nachiteration (IRS) & C \\
Eigenwerte (RCWA) & Z & C & C \\
Exponentiale, stabile $S$-Formeln, Effizienzen & FP64 & FP64 & FP64 \\
\bottomrule
\end{tabular}
\end{center}
GPU: \verb|cublasCgemm| (Math-Mode ausdrücklich \verb|CUBLAS_DEFAULT_MATH|, \textbf{kein TF32}; zusätzlich
\verb|NVIDIA_TF32_OVERRIDE=0| in der Doku), \verb|cusolverDnXgetrf|/\verb|Xgetrs| und \verb|cusolverDnXgeev| mit
\verb|CUDA_C_32F|, gemischtes Lösen mit der IRS-API (\verb|cusolverDnZCgesv| bzw.\ \verb|cusolverDnIRSXgesv|).
CPU: Eigen \verb|MatrixXcf| (bzw.\ LAPACK \verb|cgetrf|/\verb|cgeev|), Nachiteration selbst implementiert:
$x\leftarrow x+A^{-1}_{32}(b-Ax)$ mit Residuum in FP64, höchstens 10 Schritte, Abbruch bei
$\|b-Ax\|\le\epsilon_{64}\|A\|\|x\|\cdot10$, sonst Rückfall auf FP64-LU (Diagnose).

\subsection{Angepasste Parameter}
\begin{itemize}[nosep]
  \item Scheibenkriterium \eqref{eq:slicecrit}: Wachstumsgrenze $G_{\max}=\epsilon^{-1/2}$ der Arbeitsgenauigkeit der
        Propagation. FP64: $6.7\cdot10^{7}\approx10^8$ (wie bisher), FP32/mixed: $\approx4\cdot10^{3}$, also
        $\ln G_{\max}\approx8.3$ statt 18.4 $\Rightarrow$ mehr, dünnere Scheiben. \verb|growthLimit| bekommt den Standard
        \verb|auto| (abhängig vom Modus); ein expliziter Wert gilt weiter.
  \item RK: \verb|rk.tol| wird in FP32 auf $\ge10^{-5}$ angehoben (Warnung).
  \item Eigenwertklassifikation: Toleranzen relativ zu $\max|\lambda|$ mit $\epsilon_{32}$ skalieren.
\end{itemize}

\subsection{Schutzmechanismen}
Diagnosen je Rechnung: Energiebilanz (verlustfrei), Konditionszahlen ($T_{11}$, $G$), Zahl der Nachiterationen.
Option \verb|"precision_fallback": true| (Standard in \verb|fp32|): Rechnung automatisch in FP64 wiederholen, wenn bei
verlustfreien Strukturen $|\sum e-1|>10^{-3}$ oder $\operatorname{cond}(T_{11})>10^{5}$; Warnung im Ergebnis.

\subsection{Standardwahl (Entscheidung des Projektinhabers)}
\verb|"precision": "auto"| ist der Standard und bedeutet: CPU-Pfad immer \verb|fp64|; auf der GPU \verb|fp32| mit
\verb|precision_fallback| auf einer GeForce RTX~3090 (Anwenderinstallation Windows~11), sonst \verb|fp64|
(Server, RTX 6000 Ada). Erkennung über \verb|cudaGetDeviceProperties| (Name und Compute Capability 8.6); die Regel steht
in einer Tabelle \verb|src/cuda/device_policy.cpp|, damit weitere Karten ergänzt werden können. Überschreibbar mit
\verb|FG_PRECISION| und in der Konfiguration. Das Ergebnis meldet die tatsächlich verwendete Genauigkeit
(\verb|diagnostics.precision|) und ob ein Rückfall auf FP64 stattfand.

\subsection{Build}
\verb|CMAKE_CUDA_ARCHITECTURES| \verb|"86;89"| in den CUDA-Presets (RTX 3090 und RTX 6000 Ada); Preset
\verb|win-cuda| ebenfalls \verb|86|. Speichertabelle (Abschnitt Architektur) um FP32 und 24\,GB (3090) ergänzen.

\paragraph{Umsetzung (v2.8).} \verb|linalg::Precision| wird je Thread gesetzt; in FP32/Mixed läuft jede
$O(n^3)$-Operation über ein Backend (GPU über der Schwelle, sonst ein CPU-Backend in einfacher Genauigkeit: Eigen
\verb|MatrixXcf|, eigene Nachiteration). Damit ist die Studie auch ohne GPU möglich. Auf der GPU tragen die
Gerätematrizen den Elementtyp (Z oder C); die Host-API (\verb|gemm|/\verb|solve|/\verb|eig| auf \verb|MatC|) ist auf
den residenten Operationen aufgebaut. FP32 läuft resident (Umwandlung beim Hoch-/Herunterladen); Mixed nutzt die
Host-API, weil die Typumwandlung auf dem Gerät Kernel bräuchte, und löst mit \verb|cusolverDnZCgesv|. Die
Klassifikationstoleranz ist $10^{-9}$ (FP64) bzw. $10^{-4}$ (C-Eigenwerte) relativ zu $\max|\lambda|$; die
Konsistenzprüfung lässt $2\tau$ Überlappung zu (nach dem Fluss getrennte Moden mit $|\operatorname{Re}\lambda|\le\tau$).
Der Rückfall prüft die Energiebilanz; $\operatorname{cond}(T_{11})$ wird nicht berechnet. Die Produkte
$\Psi_{\text{ref}}^{-1}\Psi$ (Substrat, RK-Scheiben, Host-Pfad der Scheiben-$S$) laufen jetzt ebenfalls über das Backend
(vorher dichte FP64-GEMM auf der CPU, bei $N=20$ die Hälfte der Laufzeit).

\subsection{Tests und Abnahme (M16)}
\begin{itemize}[nosep]
  \item Unit: Backend-Operationen in \verb|fp32| gegen \verb|fp64| (gut konditionierte Zufallsmatrizen, $n=50\dots400$):
        relativer Fehler $\le10^{-5}$; IRS-Lösen erreicht FP64-Genauigkeit ($10^{-12}$) für $\operatorname{cond}\le10^6$.
  \item Genauigkeitsstudie: alle Validierungsfälle (V1--V18) in allen drei Modi, Tabelle in \verb|docs/validation.md|
        (maximaler relativer Fehler der Effizienzen gegen \verb|fp64|, Energiebilanz, Fallback ausgelöst ja/nein).
        Zielwerte vor der Studie: \verb|mixed| $\le10^{-6}$, \verb|fp32| $\le10^{-4}$ für V1--V9. \textbf{Nach der Studie
        (v2.8):} \verb|mixed| und \verb|fp32| $\le10^{-4}$ für V1--V9 (gemessen $\le2\cdot10^{-5}$ auf CPU und RTX~3090);
        resonante 2D-Fälle (V10, V15) bis $\approx10^{-4}$. \verb|mixed| ist kaum genauer als \verb|fp32|, weil die
        Eigenprobleme (in beiden Modi C) den Fehler bestimmen; $10^{-6}$ ist mit C-Eigenwerten nicht erreichbar.
  \item Laufzeittabelle RTX~3090 und RTX 6000 Ada: 2D, $N=10,15,20$, alle Modi.
  \item \verb|AGENTS.md| Abschnitt 4: \enquote{Kern ausschließlich \texttt{complex<double>}; einfache Genauigkeit nur im
        Backend (M16).}
\end{itemize}

%=====================================================================
\section{M17: Koordinatentransformationen, adaptive räumliche Auflösung}\label{sec:ctrans}
%=====================================================================

\subsection{Formalismus}
Neue Koordinaten $X,Z$ mit
\begin{equation}
  x=f(X),\qquad z=g(Z),\qquad f(X+d_x)=f(X)+d_x,\quad g(Z+d_z)=g(Z)+d_z
  \label{eq:ct_map}
\end{equation}
(separabel, $f',g'\neq0$; ASR: reell und monoton; PML, M18: komplex in Randzonen). Die Maxwellgleichungen bleiben
unverändert, nur $\partial_x=(1/f'(X))\,\partial_X$, $\partial_z=(1/g'(Z))\,\partial_Z$. Da $f(X)-X$ periodisch ist, sind
die Felder in $X$ pseudoperiodisch mit denselben $\alpha_n$:
$e^{i\alpha_0f(X)}=e^{i\alpha_0X}\cdot(\text{periodisch})$. In der Entwicklung \eqref{eq:pseudo} nach $X,Z$ gilt
\begin{equation}
  \partial_x\to i\,A,\quad A=\tp{1/f'}\,\alpha,\qquad \partial_z\to i\,G,\quad G=\tp{1/g'}\,\gamma .
  \label{eq:ct_deriv}
\end{equation}
$\tp{1/f'}$ ist im 2D-Einzelindexraum Toeplitz in $n$ und Einheit in $m$ (umgekehrt für $\tp{1/g'}$). Deshalb
kommutieren $\tp{1/f'}$ und $\tp{1/g'}$ mit $\gamma$ bzw.\ $\alpha$ und miteinander, also $AG=GA$. $1/f'$ ist stetig;
die Laurent-Regel ist zulässig.

\begin{tcolorbox}[title=Ersetzungsregel]
In \eqref{eq:Ey}, \eqref{eq:Hy}, \eqref{eq:Mgeneral}, \eqref{eq:Mmu1a}, \eqref{eq:Mmu1b} wird jedes $\alpha$ durch $A$ und
jedes $\gamma$ durch $G$ ersetzt, \textbf{unter Beibehaltung der Reihenfolge}. Die Kurzformen in \eqref{eq:Mmu1b}
stammen aus $\alpha\,\mathcal U\,\alpha$ usw.\ mit $\mathcal U=\mathbb 1$; daraus werden
$\alpha^2\to AA$, $\alpha\gamma\to AG$, $\gamma^2\to GG$, $\gamma\alpha\to GA$. Beispiel:
$M_{14}=iA\,\mathcal E\,A-i\mathbb 1$, $M_{32}=i\big(Q_{\varepsilon,zz}-Q_{\varepsilon,zy}\mathcal EQ_{\varepsilon,yz}-AA\big)$.
Unit-Test: Für $f'\equiv g'\equiv1$ ist der Pfad mit Matrizen $A,G$ gleich dem Diagonalpfad ($10^{-15}$);
der Diagonalpfad bleibt Standard (bitgleich).
\end{tcolorbox}

\subsection{Permittivität und Normalenfeld in transformierten Koordinaten}
$\varepsilon$ wird als Funktion von $(X,Z)$ entwickelt: $\varepsilon(f(X),y,g(Z))$.
\begin{itemize}[nosep]
  \item Lamellen: Bruchstellen $X_j=f^{-1}(x_j)$, dann \eqref{eq:interval} unverändert.
  \item Polygone, Konturen, CAD-Schnitte: Rand abbilden $X(t)=f^{-1}(x(t))$, $Z(t)=g^{-1}(z(t))$ (Newton; $f$ monoton im
        physikalischen Bereich), $X'=x'/f'(X)$, dann \eqref{eq:contour} mit derselben adaptiven Quadratur.
  \item Profile: Nullstellen von $g(f(X))-y$ in $X$.
  \item Normalenfeld: \emph{physikalische} Normalen an den Bildpunkten. Die Feldkomponenten $E_x,E_z,\dots$ bleiben
        kartesische Komponenten im physikalischen Raum, daher werden die Normalen nicht transformiert.
\end{itemize}

\subsection{Homogene Halbräume, Einfall und Effizienzen}
In transformierten Koordinaten sind die Ordnungen keine Eigenvektoren mehr. Daher:
\begin{itemize}[nosep]
  \item Modenbasis von Substrat und Superstrat numerisch: Eigenzerlegung von $M$ (mit $A$, $G$, $Q_{\varepsilon}=\varepsilon\mathbb 1$),
        Klassifikation auf/ab wie bei anisotropen Medien (Abschnitt~\ref{sec:homog}).
  \item Einfallende Welle: Koeffizienten von $e^{i\alpha_0f(X)+i\gamma_0g(Z)}$ in der Basis $e^{i(\alpha_nX+\gamma_mZ)}$,
        \begin{equation}
          c_{nm}=\frac{1}{d_xd_z}\int_0^{d_x}\!\!\int_0^{d_z}e^{i\alpha_0(f(X)-X)+i\gamma_0(g(Z)-Z)}\,e^{-i(nK_xX+mK_zZ)}\,dX\,dZ
          \label{eq:ct_inc}
        \end{equation}
        (separabel, DFT), multipliziert mit der Polarisation \eqref{eq:polarization}; danach in die Modenbasis projizieren.
  \item Physikalische Rayleigh-Amplituden aus dem reflektierten bzw.\ transmittierten Tangentialfeld $U(X)=\sum_mu_me^{i\alpha_mX}$
        (1D; 2D separabel analog):
        \begin{equation}
          R_n=\frac{1}{d_x}\int_0^{d_x}U(X)\,e^{-i\alpha_nf(X)}\,f'(X)\,dX .
          \label{eq:ct_rayleigh}
        \end{equation}
        Der Integrand ist periodisch und glatt; Trapezregel mit $\ge8(2N_x+1)$ Punkten. Effizienzen dann mit
        \eqref{eq:eff_closed} aus den physikalischen Amplituden.
  \item Nahfeld: Synthese in $X$, Ausgabe auf einem Gitter in $x$ über $X=f^{-1}(x)$.
\end{itemize}

\subsection{ASR-Abbildung (Granet)}
Für Bruchstellen $0\le x_0<x_1<\dots<x_J=x_0+d_x$ (Standard: alle Materialsprünge aller Schichten in $x$, oder Liste aus
der Konfiguration) und Stärke $0\le\eta<1$, auf $[x_{j-1},x_j]$ mit $\Delta_j=x_j-x_{j-1}$:
\begin{equation}
  f(X)=X-\frac{\eta\,\Delta_j}{2\pi}\sin\frac{2\pi(X-x_{j-1})}{\Delta_j},\qquad
  f'(X)=1-\eta\cos\frac{2\pi(X-x_{j-1})}{\Delta_j}.
  \label{eq:asr}
\end{equation}
$f(x_j)=x_j$, $f'$ ist stetig (an jeder Bruchstelle $1-\eta$), monoton für $\eta<1$; ein gleichmäßiges Raster in $X$ wird an
den Bruchstellen in $x$ um den Faktor $1-\eta$ verdichtet. Lamellenbruchstellen bleiben unverändert ($X_j=x_j$).
Koeffizienten von $1/f'$ per DFT mit $\ge64(2N_x+1)$ Punkten ($1/f'$ ist stetig mit stetiger Ableitung).
Konfiguration:
{\footnotesize\begin{verbatim}
"coordinate_transform": {"x": {"type": "asr", "eta": 0.8, "breakpoints_um": "auto"}, "z": {...}}
\end{verbatim}}
Gekrümmte Ränder in 2D (\enquote{matched coordinates}, nicht separabel) sind \emph{nicht} Teil von M17.

\paragraph{Umsetzung (v2.9).} \verb|CoordinateMap| (\verb|src/geometry/coordinate_transform.cpp|) mit $f$, $f'$, $f^{-1}$
(Newton mit Bisektion). Alle Integrale über eine Periode in $X$ (Koeffizienten von $1/f'$, $C$ aus
\eqref{eq:ct_rayleigh}, $D$ aus \eqref{eq:ct_inc}, Normalenfelder der Profile) mit zusammengesetzter
Gauß-Legendre-Quadratur zwischen den Bruchstellen (16 Punkte, Panelbreite $\le d/(2\pi p_{\max})$ und
$\le$ halber Abstand der Polstelle von $1/f'$); das ersetzt DFT ($\ge64(2N+1)$ Punkte) und Trapezregel, die an den
Bruchstellen ($f'''$ springt) nur algebraisch konvergieren. Halbräume: Eigenzerlegung von $M$ (reduzierter Pfad), dann
Klassifikation aller $4n_f$ Moden mit \eqref{eq:classify} (Fluss für $|\operatorname{Re}\lambda|\le\tau$): im reduzierten
Pfad entscheidet sonst das Vorzeichen des Rundungsrauschens von $\operatorname{Im}\kappa$ propagierender Moden über
auf/ab -- in inneren Schichten unschädlich, in der Referenzbasis nicht. Homogene isotrope Medien (Halbräume und
Schichten): $PR=(A^2+G^2-\varepsilon)\otimes\mathbb 1_2$, die TE/TM-Paare sind exakt entartet und die Eigenvektoren der
vollen reduzierten Matrix (Rückwärtssubstitution der Schur-Form, nicht normal) für kleines $\eta$ schlecht konditioniert
(Energiefehler bis $3\cdot10^{-3}$ bei $\eta=10^{-6}$); daher wird nur der $n_f$-Block zerlegt und seine Eigenvektoren
für beide Komponenten verwendet. Einfall: $D$ angewandt auf das Tangentialfeld
der einfallenden Welle, Projektion auf die abwärts laufenden Moden, Anteile auf evaneszenten abwärts laufenden Moden
(Abschneiderauschen, über dem Stapel mit $e^{\lambda y}$ verstärkt) werden verworfen; Effizienzen: $C$ auf die Tangentialfelder, dann die
analytische Basis der physikalischen Ordnungen. Polygone: Rand als Kontur, abgebildet und mit \eqref{eq:contour}
integriert; Raster von Algorithmus NF bilinear auf das $X,Z$-Raster umgetastet. Profile 2D: Linien bei $z=g(Z)$, die
Bruchstellen von $g$ als zusätzliche kritische Linien. \verb|split_polarizations| wird mit Transformation nicht
genutzt. Standard-Bruchstellen \enquote{auto}: Materialsprünge der Lamellen, Ecken der Polygone.

\subsection{Tests und Abnahme (M17)}
\begin{itemize}[nosep]
  \item $\eta=0$: bitgleich mit dem Pfad ohne Transformation.
  \item V1, V2, V3, V9 mit $\eta=0.5$: Ergebnisse auf $10^{-8}$ (bei moderatem $N$) wie ohne Transformation;
        Energiebilanz und Reziprozität wie bisher. \textbf{Nach der Umsetzung (v2.9):} $10^{-8}$ gilt für die ebenen Fälle
        V1, V2 (eine Bruchstelle: $10^{-12}$; zwei: $10^{-8}$ ab $N=16$). Bei V3, V9 unterscheiden sich die Ergebnisse bei
        gleichem $N$ um den Abschneidefehler ($10^{-6}$ bei $N=40$); geprüft werden der gemeinsame Grenzwert (Abweichung
        von $N=120$ ohne Transformation kleiner als ohne ASR), die Energiebilanz ($10^{-7}$ bei $N=40$, die Umrechnung der
        Ordnungen konvergiert wie $N^{-4}$), RK $=$ RCWA und z-invariantes 2D $=$ 1D konisch ($10^{-10}$). Zusätzlich:
        transformierter Pfad mit $f'\equiv g'\equiv1$ $=$ Diagonalpfad ($10^{-12}$).
  \item \textbf{V19} (Metall-Lamellengitter TM, z.\,B.\ Gold oder Aluminium wie V13 als Lamelle): Referenz ohne ASR bei sehr
        großem $N$; mit ASR ($\eta=0.8$) gleicher Grenzwert ($10^{-4}$); Konvergenztabelle (benötigtes $N$ für $10^{-3}$
        mit/ohne ASR) in \verb|docs/validation.md|. Literaturvergleich mit Granet (1999), sofern die dortigen Daten
        bibliographisch geprüft sind.
\end{itemize}

%=====================================================================
\section{M18: PML und isolierte Strukturen}\label{sec:pml}
%=====================================================================

\subsection{Ziel}
Aperiodische (isolierte) Strukturen: Ein einzelnes Objekt wird in einer Rechenzelle mit absorbierenden Randzonen (PML)
gerechnet. Zielanwendungen (Projektinhaber): \textbf{Steg} und \textbf{Graben} (1D) sowie \textbf{Kontaktloch} (2D, Loch
in einem Film), dazu Säulen und Wellenleiterenden. Weil diese Anwendungen schrägen Einfall brauchen, gehört die
Streufeldformulierung (Abschnitt~\ref{sec:scatfield}) zum Pflichtumfang von M18. Die PML ist eine komplexe Koordinatenstreckung nach
\eqref{eq:ct_map}--\eqref{eq:ct_deriv} (M17).

\subsection{PML-Abbildung}
Rechenzelle $[0,d_x)$, physikalischer Bereich $[w,d_x-w]$ mit $f(X)=X$, Randzonen der Breite $w$ an beiden Enden.
Mit der Eindringtiefe $\xi\in[0,1]$ ($\xi=(w-X)/w$ links, $\xi=(X-d_x+w)/w$ rechts) und $u(\xi)=\sin^2(\pi\xi/2)$:
\begin{equation}
  \frac{1}{f'(X)}=\frac{1-u(\xi)}{1+i\beta\,u(\xi)} .
  \label{eq:pml}
\end{equation}
Eigenschaften: $1/f'=1$ und stetig differenzierbar am Übergang ($u'(0)=0$); $1/f'\to0$ am Zellrand (Abbildung
nach $|x|\to\infty$, vgl.\ Hugonin/Lalanne); $\operatorname{Im}f'>0$ dämpft bei der Zeitkonvention $e^{-i\omega t}$
auslaufende Wellen. Standard $w=0.15\,d_x$, $\beta=1$. \eqref{eq:pml} ist ein Vorschlag; der Agent vergleicht ihn in der
Konvergenzstudie mit der Formulierung von Hugonin/Lalanne (JOSA A 22, 1844, 2005; vor Verwendung prüfen) und
übernimmt die bessere als Standard. Analog in $z$ (2D).

\subsection{Voraussetzungen und Einfall}
\begin{itemize}[nosep]
  \item In den PML-Zonen ist jede Schicht in $x$ (bzw.\ $z$) homogen; Strukturen reichen nicht hinein. Geprüft beim
        Aufbau (Schnittkurven/Bruchstellen in den Zonen $\Rightarrow$ Fehler). Geschichtete Hintergründe (Substrat,
        Filme) laufen durch die Zonen.
  \item Einfall (a): ebene Welle bei Normaleinfall ($\alpha_0=\gamma_0=0$). Das konstante Feld ist auch in gestreckten
        Koordinaten exakt ($c_{nm}=\delta_{n0}\delta_{m0}$). Schräger Einfall ist mit seitlichen PML in der
        Totalfeldrechnung nicht darstellbar (die einfallende Welle würde in den Zonen gedämpft); dafür gibt es die
        Streufeldformulierung (c).
  \item Einfall (b): Mode eines \emph{strukturierten} Ein- oder Ausgangsmediums. Substrat und Superstrat dürfen
        $y$-invariant strukturiert sein (Lamellen, Polygon, Kontur), z.\,B.\ ein Wellenleiter, der sich in $y$ ausbreitet.
        Modenbasis aus der RCWA-Eigenzerlegung dieses Mediums; Auswahl über effektiven Index oder Modennummer.
        Normierung der Moden unkonjugiert, $\int(\mathbf E\times\mathbf H)\cdot\hat{\mathbf y}\,dx=1$ (übliche Wahl mit PML);
        für geführte Moden ist $|r|^2$ dann die Leistungsreflexion.
  \item Einfall (c): ebene Welle unter beliebigem Winkel über die Streufeldformulierung (Abschnitt~\ref{sec:scatfield}).
  \item Graben und Kontaktloch sind Löcher in Filmen, die durch die PML-Zonen laufen; das ist zulässig, weil der Film in den
        Zonen $x$-invariant ist. Hintergrund einer Schicht ist dann das Filmmaterial.
\end{itemize}

\subsection{Streufeldformulierung (schräger Einfall)}\label{sec:scatfield}
Hintergrundstapel: jede strukturierte Schicht durch ihr Hintergrundmaterial ersetzt. Seine Lösung für die einfallende
ebene Welle ist analytisch (Airy/Fresnel, bestehender Code für V2):
$\mathbf E_{\text{bg}}=e^{i(\alpha_0x+\gamma_0z)}\mathbf e(y)$, $\mathbf H_{\text{bg}}$ analog, in jeder Schicht eine Summe aus
$e^{\pm i\beta y}$-Termen. Mit $\mathbf E=\mathbf E_{\text{bg}}+\mathbf E_s$ erfüllt das Streufeld die Maxwellgleichungen
mit der Polarisationsquelle
\begin{equation}
  \mathbf D_s=\overline{\overline\varepsilon}\,\mathbf E_s+\mathbf P,\qquad
  \mathbf P=(\overline{\overline\varepsilon}-\overline{\overline\varepsilon}_{\text{bg}})\,\mathbf E_{\text{bg}},
  \label{eq:scat_source}
\end{equation}
die nur in der Struktur, also im physikalischen Bereich ($f(X)=X$), von null verschieden ist. Das Streufeld wird ohne
Blochphase entwickelt ($\alpha_n=nK_x$, $\gamma_m=mK_z$); es klingt in den PML-Zonen ab.
\begin{itemize}[nosep]
  \item \textbf{Koeffizienten der Quelle:} $\mathbf E_{\text{bg}}$ ist innerhalb einer Schicht in $x,z$ stetig, also gilt für das
        Produkt die Laurent-Regel exakt. Für jedes Material $b$ der Schicht:
        $[\chi_be^{i(\alpha_0x+\gamma_0z)}]_{pq}$ ist \eqref{eq:contour} bzw.\ \eqref{eq:polygon}/\eqref{eq:interval} mit dem
        verschobenen Wellenvektor $\mathbf k=(pK_x-\alpha_0,\ qK_z-\gamma_0)$ (der Fall $\mathbf k=0$ tritt nur für
        $\alpha_0\in K_x\mathbb Z$ auf). Damit $[\mathbf P](y)=\sum_b[\chi_be^{i(\cdot)}](\varepsilon_b-\varepsilon_{\text{bg}})\mathbf e(y)$.
  \item \textbf{Inhomogene Propagationsgleichung:} Aus \eqref{eq:Ey} mit $[D_y]=\dots+[P_y]$ und den Zeilen von \eqref{eq:ode}
        ($Q_\mu=\mathbb 1$, mit $A,G$ aus M17):
        \begin{equation}
          \frac{dF_s}{dy}=M\,F_s+\sigma(y),\qquad
          \sigma=\begin{pmatrix}-iA\,\mathcal E[P_y]\\ -iG\,\mathcal E[P_y]\\ i\big([P_z]-Q_{\varepsilon,zy}\mathcal E[P_y]\big)\\
          -i\big([P_x]-Q_{\varepsilon,xy}\mathcal E[P_y]\big)\end{pmatrix}.
          \label{eq:scat_ode}
        \end{equation}
        Vorzeichen und Faktorisierung werden mit dem Test \enquote{Streufeld + Hintergrund $=$ Totalfeld} (V20d) geprüft.
  \item \textbf{RK:} Das implizite RK-Verfahren behandelt \eqref{eq:scat_ode} direkt (Quellterm in den Stufengleichungen).
        Jede Scheibe liefert neben $S$ zwei Quellvektoren: die von der Quelle erzeugten aus- und abwärts laufenden
        Amplituden bei verschwindendem Einfall.
  \item \textbf{RCWA:} Je Scheibe partikuläre Lösung je Exponentialterm von $\sigma$: für $\sigma=\sigma_\kappa e^{i\kappa y}$ ist
        $F_p=(i\kappa\mathbb 1-M)^{-1}\sigma_\kappa e^{i\kappa y}$; liegt $i\kappa$ nahe an einem Eigenwert von $M$
        (Abstand $<10^{-8}\max|\lambda|$), wird die Scheibe mit RK gerechnet (Warnung).
  \item \textbf{Verkettung:} erweitertes Sternprodukt, das die Quellvektoren der Teilstapel mitführt (Quellen wirken wie
        zusätzliche einlaufende Amplituden an den Scheibengrenzen). Herleitung im Plan nachtragen, sobald implementiert,
        mit Test gegen die sequentielle Rechnung.
  \item \textbf{Ausgaben:} wie unten; Leistungsflüsse und Nahfeld des Streufeldes direkt, das Totalfeld als Summe.
\end{itemize}

\subsection{Ausgaben}
\begin{itemize}[nosep]
  \item Gestreutes Feld $\mathbf E_s=\mathbf E-\mathbf E_{\text{bg}}$ mit der Hintergrundlösung: derselbe Stapel, jede
        strukturierte Schicht durch ihr Hintergrundmaterial ersetzt, gleiche Transformation und gleiches $N$
        (Diskretisierungsfehler heben sich auf).
  \item Leistungsflüsse des gestreuten Feldes durch Ebenen $y=\text{const}$ über den physikalischen Bereich,
        $P=\tfrac12\operatorname{Re}\int(\mathbf E_s\times\overline{\mathbf H_s})\cdot\hat{\mathbf y}\,dx$, normiert auf die
        einfallende Intensität: Streubreite (1D, in µm) bzw.\ Streuquerschnitt (2D). Absorption aus der Flussbilanz.
  \item Bei modalem Einfall: Modenkoeffizienten $r_{jk}$, $t_{jk}$ direkt aus der S-Matrix in der Modenbasis,
        Leistungsanteile $|r_{jk}|^2$, $|t_{jk}|^2$ für geführte Moden, Abstrahlverlust $1-\sum$.
  \item Nahfeld (gesamt und gestreut) auf dem physikalischen Bereich.
\end{itemize}
Konfiguration:
{\footnotesize\begin{verbatim}
"aperiodic": {"pml": {"x": {"width_um": 0.3, "beta": 1.0}}, "excitation": "plane_wave"}   // oder {"mode": 0}
\end{verbatim}}

\paragraph{Umsetzung (v2.10).} \verb|PmlMap| (\verb|src/geometry/aperiodic.cpp|) liefert $\tp{1/f'}$ nach
\eqref{eq:pml} (Gauß-Legendre zwischen $0,w,d-w,d$); die Geometrie bleibt unverändert, weil $f(X)=X$ überall dort gilt,
wo Strukturen liegen (\verb|checkPmlZones|, auch für CAD-Schnitte und strukturierte Halbräume). Normaleinfall:
Totalfeld und Hintergrundstapel (\verb|backgroundStack|: Material am Zellrand bzw. Hintergrund, ein 1D-Profil wird
in zwei Schichten bei seiner Höhe in den Zonen geteilt) mit demselben $N$. Streufeld: statt \eqref{eq:scat_source} mit
exakten Laurent-Koeffizienten wird die Quelle als $\sigma=(M-M_{\text{bg}})F_{\text{bg}}$ gebildet ($M_{\text{bg}}$:
Hintergrundmaterial mit denselben $A,G$). Das ist zeilenweise identisch mit \eqref{eq:scat_ode}, benutzt aber für
$P=(Q_\varepsilon-\varepsilon_{\text{bg}})[\mathbf E_{\text{bg}}]$ dieselbe Faktorisierung wie das Totalfeld: Bei Li- bzw.
FFF-Regeln ist $\varepsilon\mathbf E_s$ nicht stetig, die exakten Koeffizienten von $\Delta\varepsilon\mathbf E_{\text{bg}}$
wären mit der Faktorisierung von $\mathbf E_s$ inkonsistent. Damit gilt V20d bis auf Rundung. Hintergrundfeld:
Hintergrundstapel bei $N=0$ (Airy), je Schicht vier Exponentialmoden. Schräger Einfall mit PML: das Streufeld wird
ohne Blochphase entwickelt, $\sigma$ in der Blochbasis ($\alpha_0+nK$, ohne PML-Faktor, weil $\sigma$ im physikalischen
Bereich liegt) gebildet und mit $T_{mn}=\frac1d\int_w^{d-w}e^{i(\alpha_0+(n-m)K)x}dx$ entwickelt (stetig bei
$\alpha_0\to0$). RCWA: partikuläre Lösung je Exponentialterm, aus den Werten an den Scheibengrenzen
$u=g^1_b-S_{11}g^0_b-S_{12}g^1_a$, $v=g^0_a-S_{21}g^0_b-S_{22}g^1_a$ ($g=\Psi^{-1}F_p$); Resonanz ($<10^{-8}\max|\lambda|$):
RK mit Warnung. RK: erweitertes System $(F;1)$. Verkettung: $u=u_U+U_{11}(1-L_{12}U_{21})^{-1}(u_L+L_{12}v_U)$,
$v=v_L+L_{22}(1-U_{21}L_{12})^{-1}(v_U+U_{21}u_L)$ (sequentiell und Baum), innere Amplituden analog. Ausgaben: Flüsse über das
Fenster mit dem Hintergrund als zusätzlicher \enquote{Ordnung} $(\alpha_0,\gamma_0)$; Absorption aus dem Fluss des
Totalfeldes durch die geschlossene Box (Fenster $\times$ Schichten, Seiten mit Gauß-Knoten in $y$) minus Hintergrund;
zusätzlich der seitliche Streufluss, denn nur oben $+$ unten $+$ seitlich ist von der Fenstergröße unabhängig.
Intensität: $|\mathbf S|$ der einfallenden Welle. Strukturierte Halbräume: Modenbasis aus der Eigenzerlegung von $M$;
geführte Moden: $n_{\text{eff}}=|\operatorname{Im}\lambda|$ über dem Mantelindex, $|\operatorname{Re}\lambda|<10^{-3}n_{\text{eff}}$
und $\ge90\,\%$ von $\int|E|^2$ im Fenster; Leistung über den konjugierten Fluss im Fenster (statt der unkonjugierten
Normierung, deren Integral über die gestreckte Zelle nicht konvergiert). Klassifikation der Halbraummoden: propagierende
Moden ($|\operatorname{Re}\lambda|\le10^{-6}|\lambda|$) nach dem Fluss (mit PML ist das Rauschen in
$\operatorname{Re}\lambda$ bis $10^{-7}|\lambda|$).

\subsection{Tests und Abnahme (M18)}
\begin{itemize}[nosep]
  \item \textbf{Nach der Umsetzung (v2.10):} V20a, V20d, V20e wie gefordert. V20b: Konvergenz und Unabhängigkeit von
        $w,\beta$ deutlich unter $10^{-3}$; die Absorption verlustfreier Stege ist $\le2\cdot10^{-4}$ statt $10^{-6}$ (das
        abgeschnittene PML-System erhält die Energie nur bis auf den Diskretisierungsfehler). V20c: abruptes Ende
        konvergiert auf $10^{-4}$, der Literaturvergleich ist offen (Wert nicht geprüft). V20f: Gesamtstreuung konvergiert
        auf $10^{-4}$, oben/unten nur $\sim1/L$; Reziprozität über das Winkelspektrum des Streufeldes auf einer Linie über der Struktur
        (asymmetrisch, TE/TM): $5\cdot10^{-2}$ bei $L=8\,$µm, fallend mit $L$ (Toleranz $0{,}1$; begrenzt durch die endliche Linie). V20g: Symmetrie, CAD $=$ Contour2D und
        z-invariantes 2D $=$ 1D auf $10^{-10}$; die Konvergenzstudie gehört auf das Zielsystem. Vergleich mit
        Hugonin/Lalanne nicht durchgeführt. Messwerte in \verb|docs/validation.md|.
  \item \textbf{V20a}: homogener Schichtstapel mit PML bei Normaleinfall $=$ Fresnel/Airy auf $10^{-10}$.
  \item \textbf{V20b}: isolierter dielektrischer Steg auf Substrat (1D, TE und TM): Streubreite konvergiert bei
        Zellbreite $4,6,8\times$ Stegbreite auf $10^{-3}$ und hängt nicht von $w$, $\beta$ ab ($10^{-3}$);
        Absorption $=0$ bei verlustfreiem Material ($10^{-6}$).
  \item \textbf{V20c}: Reflexion der Grundmode an einem abrupt endenden Schichtwellenleiter gegen Silberstein et~al.\
        (JOSA A 18, 2865, 2001; Wert vor Verwendung bibliographisch prüfen), Toleranz $10^{-3}$.
  \item \textbf{V20d} (Streufeldformulierung, ohne PML): periodisches Lamellengitter und \verb|Contour2D|-Säulen unter
        schrägem Einfall (auch konisch). Hintergrund $+$ Streufeld (Entwicklung mit $\alpha_n=\alpha_0+nK_x$) $=$ Totalfeldrechnung
        auf $10^{-10}$ (RCWA) bzw.\ RK-Toleranz. Damit sind \eqref{eq:scat_source}--\eqref{eq:scat_ode} unabhängig von der PML geprüft.
  \item \textbf{V20e}: Normaleinfall mit PML: Streufeldformulierung $=$ Totalfeld minus Hintergrund (Einfall a) auf $10^{-8}$.
  \item \textbf{V20f} (Graben und Steg schräg, 1D, TE/TM, $\Phi=0^\circ,30^\circ,60^\circ$): Streubreite konvergiert in der
        Zellbreite auf $10^{-3}$, unabhängig von $w,\beta$; Reziprozität der Streuamplituden bei Vertauschung von Einfalls- und
        Beobachtungsrichtung (qualitativ, Toleranz nach Studie festlegen).
  \item \textbf{V20g} (Kontaktloch, 2D, PML in $x$ und $z$): rundes Loch in einem Film auf Substrat, Normal- und Schrägeinfall;
        Konvergenz in Zellgröße und $N$ (Zielsystem); Kontaktloch als \verb|Contour2D| und als CAD-Körper identisch ($10^{-10}$).
\end{itemize}

%=====================================================================
\section{M19: Differenzierbarkeit und Optimierung}\label{sec:adjoint}
%=====================================================================

\subsection{Ziel und Umfang}
Gradienten reeller Zielfunktionen $L$ von Effizienzen oder komplexen Amplituden nach Strukturparametern, mit einem
Aufwand von höchstens 3--5 Vorwärtsrechnungen \emph{unabhängig von der Zahl der Parameter}. Parameter (Phase~1):
Dicke $y$-invarianter Schichten; Materialpermittivität (komplex); Lamellen-Bruchstellen; Polygonecken;
Kontrollpunkte und Gewichte von \verb|Contour2D|; Profilparameter (Tiefe, Verschiebung); CAD-Kontrollpunkte für
seitliche Bewegungen (Abschnitt~\ref{sec:adj_shape}). Phase~2 (optional): Einfallswinkel und Wellenlänge.

\subsection{Methode: diskrete Adjungierte}
Abgeleitet wird die \emph{tatsächlich berechnete} diskrete Abbildung (erst diskretisieren, dann differenzieren), im
Rückwärtsmodus über die aufgezeichneten Zwischenergebnisse (\verb|storeInternal| plus neue Felder). Zwei bewusste
Festlegungen, im Ergebnis als Diagnose ausgewiesen:
\begin{enumerate}[nosep]
  \item RK-Schrittfolge, Scheibenzerlegung und Bandtopologie (M14) gelten als fest.
  \item Das Normalenfeld ($\tp{N_iN_j}$) gilt als parameterunabhängig (\enquote{eingefroren}). Es beeinflusst nur die
        Konvergenzgeschwindigkeit, nicht den Grenzwert; Algorithmus~NF ist wegen der Nächster-Punkt-Suche nicht glatt.
        Der Unterschied zum vollen Differenzenquotienten wird in den Tests gemessen und fällt mit $N$.
\end{enumerate}

\subsection{Konvention (komplex, PyTorch-kompatibel)}
Für reelles $L$ und komplexe Größe $A$ heißt $\bar A:=\partial L/\partial\operatorname{Re}A+i\,\partial L/\partial\operatorname{Im}A$;
dann $dL=\operatorname{Re}\operatorname{tr}(\bar A^HdA)$. Rückwärtsregeln (Kotangenten werden \emph{addiert}):
\begin{center}
\small
\begin{tabular}{p{4.2cm}p{10.2cm}}
\toprule
Vorwärts & Rückwärts \\
\midrule
$C=AB$ & $\bar A\mathrel{+}=\bar CB^H$, $\bar B\mathrel{+}=A^H\bar C$ \\
$X=A^{-1}B$ (LU) & $\hat B=A^{-H}\bar X$ (Lösen mit $A^H$), $\bar B\mathrel{+}=\hat B$, $\bar A\mathrel{+}=-\hat BX^H$ \\
$C=A^{-1}$ & $\bar A\mathrel{+}=-C^H\bar CC^H$ \\
$z=\phi(w)$ holomorph & $\bar w\mathrel{+}=\overline{\phi'(w)}\,\bar z$ (z.\,B.\ $e^{\lambda h}$, $\sqrt{\cdot}$ für $\beta$) \\
$h$ reell, $X=\operatorname{diag}(e^{\lambda_ih})$ & $\bar h\mathrel{+}=\operatorname{Re}\sum_i\overline{\bar X_{ii}}\,\lambda_iX_{ii}$ (normiert: Faktor $k_0$ für µm) \\
$A=W\Lambda W^{-1}$ (Eig) & $\bar A=W^{-H}\big(\operatorname{diag}(\bar\lambda)+\overline{P}\circ(W^H\bar W)\big)W^H$,
  $P_{ij}=1/(\lambda_j-\lambda_i)$ ($i\neq j$), $P_{ii}=0$ \\
Toeplitz $T=\tp f$, $T_{rr'}=f_{n-n',m-m'}$ & $\bar f_{pq}\mathrel{+}=\sum_{n-n'=p,\ m-m'=q}\bar T_{rr'}$ (Summe über die Diagonalen) \\
$e=\frac{a|B_x|^2+b|B_z|^2+2c\operatorname{Re}(B_x\overline{B_z})}{D}$ \eqref{eq:eff_closed} &
  $\bar B_x\mathrel{+}=2\bar e(aB_x+cB_z)/D$, $\bar B_z\mathrel{+}=2\bar e(bB_z+cB_x)/D$ ($a,b,c,D$ reell) \\
\bottomrule
\end{tabular}
\end{center}
\textbf{Eigenzerlegung:} Die Formel setzt voraus, dass $L$ von der Normierung und Phase der Eigenvektoren nicht abhängt
(gilt für die S-Matrix). Entartete Cluster ($|\lambda_i-\lambda_j|<\tau\max|\lambda|$, $\tau=10^{-10}$) liegen in derselben
Klasse (auf/ab) und gehen nur über ihren Unterraum ein: dort $P_{ij}:=0$. Fast entartete Paare werden mit
Lorentz-Verbreiterung $P_{ij}=\overline{(\lambda_j-\lambda_i)}/(|\lambda_j-\lambda_i|^2+\delta^2)$, $\delta=10^{-12}\max|\lambda|$
behandelt. Test mit Normaleinfall auf symmetrische Gitter (exakte Entartung).

\textbf{RK-Propagator} eines Schrittes (lineare ODE, $s$ Stufen, Butcher $\mathcal A,b,c$):
$\mathbb M=\operatorname{blockdiag}(M(y+c_ih))$, $Z=\mathbb 1-h\,\mathbb M(\mathcal A\otimes\mathbb 1)$,
$K=Z^{-1}\mathbb M(\mathbf 1\otimes\mathbb 1)$, $P=\mathbb 1+h(b^T\otimes\mathbb 1)K$. Rückwärts über die Regeln oben:
$\bar K=h(b\otimes\mathbb 1)\bar P$, $\bar R=Z^{-H}\bar K$, $\bar Z=-\bar RK^H$,
$\bar{\mathbb M}\mathrel{+}=-h\bar Z(\mathcal A\otimes\mathbb 1)^H+\bar R(\mathbf 1\otimes\mathbb 1)^H$, davon die Diagonalblöcke.
(Ist der RK-Schritt im Code anders organisiert, gilt dieselbe Herleitung für die dortige Form.)

Sternprodukt \eqref{eq:star}, $T\to S$, Grenzflächen, $Q$-Matrizen (Laurent, Li-Inverse, FFF-Varianten) und der Aufbau
von $M$ \eqref{eq:Mgeneral} werden aus diesen Primitiven zusammengesetzt; jede Funktion bekommt eine eigene
Rückwärtsfunktion in \verb|src/adjoint/| und einen eigenen Gradiententest.

\subsection{Formableitungen der Koeffizienten}\label{sec:adj_shape}
Bewegt sich der Rand eines Gebiets mit der Geschwindigkeit $\mathbf V=\partial\mathbf r/\partial\theta$, dann gilt
(Reynolds/Hadamard, mit $\mathbf n\,ds=(z',-x')\,dt$), für \emph{alle} $\mathbf k$ einschließlich $\mathbf k=0$:
\begin{equation}
  \frac{\partial f_{pq}}{\partial\theta}=\frac{1}{d_xd_z}\oint\big(V_x\,z'(t)-V_z\,x'(t)\big)\,e^{-i\mathbf k\cdot\mathbf r(t)}\,dt .
  \label{eq:shape2d}
\end{equation}
NURBS-Kontur $\mathbf r(t)=\sum_jR_j(t)\mathbf P_j$: $\partial\mathbf r/\partial P_{j,x}=R_j(t)\hat{\mathbf x}$,
$\partial\mathbf r/\partial P_{j,z}=R_j(t)\hat{\mathbf z}$, für Gewichte
$\partial\mathbf r/\partial w_j=\frac{B_j(t)}{\sum_iw_iB_i(t)}(\mathbf P_j-\mathbf r(t))$. Polygonecken sind der Fall Grad~1
(Hutfunktionen). Lamellenbruchstellen ($\varepsilon=\sum_j\varepsilon_j\chi_{[x_j,x_{j+1}]}$): $\partial f_p/\partial x_j=-\frac{\varepsilon_j-\varepsilon_{j-1}}{d}\,e^{-ipKx_j}$
(Wert rechts minus links der Bruchstelle). Gleiche Quadratur wie \eqref{eq:contour}.

\textbf{CAD (3D):} Ein Flächenparameter $\theta$ (z.\,B.\ Kontrollpunkt $\mathbf P_{ij}$ in Richtung $\mathbf e$:
$\mathbf V=R_{ij}(u,v)\,\mathbf e$) verschiebt die Schnittkurve in der Ebene mit der Normalgeschwindigkeit
$v_n=(\mathbf V\cdot\mathbf N)/|\mathbf N_\parallel|$ ($\mathbf N$ Flächennormale, $\mathbf N_\parallel$ ihr Anteil in der
$(x,z)$-Ebene):
\begin{equation}
  \frac{\partial f_{pq}(y)}{\partial\theta}=\frac{1}{d_xd_z}\oint\frac{\mathbf V\cdot\mathbf N}{|\mathbf N_\parallel|}\,
  e^{-i\mathbf k\cdot\mathbf r}\,ds .
  \label{eq:shape3d}
\end{equation}
$(u,v)$ entlang der Schnittkurve aus den p-Kurven (M14). Phase~1 erlaubt nur Parameter, die keine kritische Höhe
verschieben (seitliche Bewegungen); bewegt ein Parameter eine waagrechte ebene Fläche, ist zusätzlich
$\partial y_k/\partial\theta=\mathbf V\cdot\hat{\mathbf y}$ über die Dickengradienten der beiden angrenzenden Bänder
einzurechnen; andere Fälle $\Rightarrow$ \verb|NotImplemented("M19 Phase 2")|.

\subsection{Schnittstellen}
{\footnotesize\begin{verbatim}
// C++ (include/fg/adjoint/gradient.hpp)
struct ParameterRef { std::string pointer; };      // JSON-Pointer, z. B. "/layers/2/regions/0/nurbs/points_um/3/0"
struct Objective { std::vector<std::pair<OrderRef, double>> weights; };   // L = sum w * e  (oder Kotangenten)
struct GradientResult { double value; std::vector<double> grad; Diagnostics diag; };
GradientResult gradient(const SolverOptions&, const Stack&, const Wavelength&, const Objective&,
                        const std::vector<ParameterRef>&);
// allgemeine VJP: Kotangenten fuer alle Effizienzen und/oder komplexen Amplituden -> Parameterkotangenten
std::vector<double> vjp(const Simulation&, const Cotangents&, const std::vector<ParameterRef>&);
\end{verbatim}}
Python: \verb|sim.value_and_grad(objective, params)|, \verb|sim.vjp(cotangents, params)|;
\verb|fouriergrating.torch| (\verb|torch.autograd.Function|: vorwärts Effizienzen als Tensor, rückwärts \verb|vjp|) und
\verb|fouriergrating.jax| (\verb|jax.custom_vjp| mit \verb|jax.pure_callback|); beide optional importierbar.
Komplexe Parameter (Permittivität) folgen der Konvention oben und sind damit direkt mit PyTorch kompatibel.
Speicher: Zwischenergebnisse je Scheibe; Option \verb|"checkpoint_every": k| (nur jede $k$-te Scheibe speichern, Rest neu
rechnen). GPU: Rückwärtsfunktionen über das Backend (M15/M16).

\paragraph{Umsetzung (v2.11).} Rückwärtsregeln in \verb|src/adjoint/primitives.cpp| (Tabelle oben; zusätzlich
$T\to S$, Scheiben-S-Matrix $S=HG^{-1}$ mit $\bar H=\bar SG^{-H}$, $\bar G=-S^H\bar H$, Propagationsmatrix mit Diagonal-
und dichten $A,G$, Q-Varianten Laurent/Li/FFF\_iso/FFF\_sym/FFF\_rig einschließlich der Kotangenten von $\tp{N_iN_j}$ und
$\tp{N_i}$, RK-Schritt mit Ableitung nach $h$). Der Treiber (\verb|src/solvers/gradient.cpp|) zeichnet die Kette auf
(Scheibenmatrizen, Moden, S-Matrizen, Präfixprodukte; RK: die akzeptierten Schritte, \verb|RKSteps|) und läuft
rückwärts bis zu den Koeffizientengittern. Die Schichtmodelle liefern dazu \verb|adjointSlice|: Faktorisierung,
Gitter und die Zerlegung $\varepsilon=\sum_k\varepsilon_k\chi_k$ ($1/\varepsilon$ ebenso). Formableitungen:
\eqref{eq:shape2d} mit $\mathbf V(t)=(\mathbf r_+(t)-\mathbf r_-(t))/2\delta$ aus den bei $\theta\pm\delta$ neu
erzeugten Rändern gleicher Segmentstruktur (für Kontrollpunkte exakt, $\mathbf r$ ist linear in ihnen), bei
Schnittschichten und CAD-Körpern je Scheibenhöhe aus den Schnitten (\eqref{eq:shape3d} in dieser Form);
Lamellenkanten und Profilkreuzungen analytisch. Abweichung von Festlegung 2: Das Normalenfeld der Profile ist
analytisch und glatt und wird mitdifferenziert (zentraler Differenzenquotient der Koeffizienten, ohne
$O(n^3)$-Aufwand); eingefroren wäre der Gradient bei FFF\_rig um eine Größenordnung falsch. Das Rasterfeld von
Algorithmus~NF (2D) bleibt eingefroren. Schnittstellen: \verb|fg::gradient| (JSON-Pointer), \verb|StackParameter|
(C++, beliebige Geometrie über den Stapel bei $\theta+\delta$), \verb|Simulation::valueAndGrad|/\verb|vjp|, Python
\verb|value_and_grad|/\verb|vjp|, \verb|fouriergrating.torch|, \verb|fouriergrating.jax|. Phase~2 (v2.12): Wellenlänge,
Einfallswinkel $\Phi,\Psi$, Polarisation $\delta$ und Permittivität von Superstrat (reell) und Substrat. Sie wirken über
$\alpha_n,\gamma_m$ (Kotangenten aus allen $M$ und den analytischen Basen der Halbräume, Rückwärtsregel von
\verb|isotropicModes|), über $k_0$ (normierte Dicken, $K_x=\lambda/d_x$) und über Einfallsvektor und
Effizienznormierung; diese hängen von wenigen Skalaren ab und werden mit $O(n_f)$ Aufwand per zentralem
Differenzenquotienten abgeleitet. \verb|checkpoint_every|: gespeichert werden nur jedes $k$-te Präfixprodukt, die
Scheibendaten werden rückwärts je Abschnitt neu gerechnet (Gradient bitgleich). Anisotrope Materialien: Rückwärtsregel
von \eqref{eq:QaniA}, \eqref{eq:QaniB} und Laurent bis zu den Gittern der Koeffizientenfunktionen; deren Ableitung nach
den Regionstensoren per zentralem Differenzenquotienten der Koeffizienten (holomorph, $O(n_f)$), Kanten analytisch,
Profiltiefe über das mitbewegte Normalenfeld; Lamellen, Profile und 2D-Querschnitte mit Laurent. Nicht umgesetzt:
Rückwärtsfunktionen auf dem GPU-Backend, anisotrope 2D-Querschnitte mit FFF\_B (Rasterfeld von Algorithmus~NF),
\verb|profile2d| mit anisotropen Medien, Koordinatentransformationen und aperiodische Strukturen.

\subsection{Tests und Abnahme (M19)}
\begin{itemize}[nosep]
  \item \textbf{Nach der Umsetzung (v2.11):} Rückwärtsfunktionen $10^{-6}$; Ende-zu-Ende V3, V9, V13, V15, V18
        $\le7\cdot10^{-7}$ (Laurent/Li $\le10^{-7}$), Entartung $10^{-9}$, RK $=$ RCWA $2\cdot10^{-9}$, Aufwand
        $1{,}3$--$2{,}3\times$; Messwerte in \verb|docs/validation.md|. Phase~2 (v2.12): $\lambda,\Phi,\Psi,\delta$ und
        Superstrat/Substrat gegen Differenzenquotienten $\le2\cdot10^{-7}$ (RK $2\cdot10^{-6}$, adaptive Schritte);
        anisotrope Lamellen, Profile (Laurent, FFF\_A, FFF\_B) und 2D-Säulen (Laurent) $\le3\cdot10^{-7}$.
  \item Jede Rückwärtsfunktion: zentraler Differenzenquotient mit Zufallsmatrizen ($h=10^{-7}$, Real- und Imaginärteil),
        relative Übereinstimmung $10^{-6}$.
  \item Ende-zu-Ende je Parameterart auf V3, V9, V13, V15, V18: Adjungierte gegen Differenzenquotient mit eingefrorenem
        Normalenfeld und fester Scheibenfolge $10^{-6}$; gegen den vollen Differenzenquotienten: Abweichung tabellieren.
  \item Entartung: Normaleinfall auf symmetrisches Lamellengitter, Gradient endlich und $=$ Differenzenquotient.
  \item RK- und RCWA-Gradient stimmen im Grenzwert überein (Lamellen: $10^{-8}$).
  \item Aufwand: Gradient $\le5\times$ Vorwärtsrechnung (gemessen, Bericht).
  \item Beispiel 1 (\verb|examples/python/optimize_littrow.py|): binäres Gitter, Littrow $-1$. Ordnung, Dicke und Tastgrad
        mit \verb|scipy.optimize.minimize| (L-BFGS-B); Optimum $=$ Rasterscan-Optimum auf $10^{-3}$.
  \item Beispiel 2 (\verb|examples/python/optimize_nurbs_pillar.py|): Kontrollpunkte einer \verb|Contour2D|-Säule mit
        PyTorch-Adam; Zielfunktion fällt monoton (Demonstration, kein Referenzwert).
\end{itemize}

%=====================================================================
\section{Querschnittsthemen der Version 2}\label{sec:extcross}
%=====================================================================
\begin{itemize}[nosep]
  \item \textbf{Golden Files vor dem ersten Eingriff (Beginn M12):} Für alle \verb|examples/| Effizienzen (17 signifikante
        Stellen) mit dem Stand v0.10 nach \verb|tests/regression/golden/| schreiben. Test \verb|Regression.*|: mit
        Standardoptionen auf derselben Maschine bitgleich, in CI (anderer Compiler bzw.\ andere C-Bibliothek) $10^{-11}$ bezogen auf die
        einfallende Leistung, also absolut auf den Effizienzen. Korrektur v2.3: rein relativ scheitern Ordnungen
        mit Effizienz $\approx10^{-14}$. Korrektur v2.4: $10^{-13}$ hält nur zwischen MSVC und MinGW-GCC
        ($\le1{,}2\cdot10^{-14}$); Linux-GCC~13 gegen MinGW-GCC~16 weicht bei adaptivem RK (Gold-Echelette) und
        beim 2D-Lochgitter (V10) bis $1{,}3\cdot10^{-12}$ ab (Rundung der Mathematikbibliotheken, verstärkt durch
        Schrittweitensteuerung und Kondition). Python $=$ \verb|fgsolve| gilt im selben Build bitgleich.
  \item \textbf{CI-Matrix:} \verb|cpu-release| (wie bisher), \verb|cpu-release| + \verb|FG_ENABLE_NURBFORGE=ON| (ab M14),
        Python-Wheel-Build und pytest (ab M12), ThreadSanitizer (ab M15), Doku.
  \item \textbf{Versionen:} je Meilenstein eine Minor-Version (M12 $\to$ 0.11.0, \dots, M19 $\to$ 0.18.0); Plan-Version 2.x;
        \verb|CHANGELOG.md| je Meilenstein.
  \item \textbf{Dokumentation:} \verb|docs/config_format.md| (neue Schichttypen und Blöcke), \verb|docs/user_guide.md|,
        \verb|docs/architecture.md| (Module \verb|cad|, \verb|adjoint|, Python), \verb|docs/validation.md| (V15--V20),
        \verb|docs/benchmarks.md|.
  \item \textbf{Fehlerfunde im Plan:} wie bisher (Test, Plan und Code im selben Commit, CHANGELOG).
\end{itemize}

\subsection{Entscheidungen des Projektinhabers (5.~November 2026)}
\begin{enumerate}[nosep]
  \item M18: \enquote{iso} heißt isolierte Strukturen (Steg, Graben, Kontaktloch usw.) $\Rightarrow$ Streufeldformulierung für
        schrägen Einfall ist Pflicht (Abschnitt~\ref{sec:scatfield}).
  \item Installationen: Server AlmaLinux~10 mit RTX 6000 Ada; Anwender Windows~11 mit RTX~3090 $\Rightarrow$ MSVC-Port,
        Windows-Wheels (CPU und \verb|fouriergrating-cuda12|), Server aus den Quellen (M12).
  \item CAD-Daten erzeugt ein Sprachmodell mit Nurbforge $\Rightarrow$ einzige STEP-Quelle ist der native Nurbforge-Schreiber;
        Geometrieprüfung mit Fehlercodes für die Prüfschleife des Sprachmodells (Abschnitt~\ref{sec:cadcheck}); Einheit µm.
  \item Auf der RTX~3090 ist FP32 der Standard (mit Rückfall auf FP64), auf dem Server FP64 (M16).
\end{enumerate}
Weitere Rückfragen nur, wenn ein Meilenstein auf eine im Plan nicht entschiedene Frage stößt.

\subsection*{Literatur zu Version 2 (vor Verwendung bibliographisch prüfen)}
\begin{enumerate}[label={[\arabic*]},nosep,leftmargin=2em,start=11]
  \item L.~Piegl, W.~Tiller, \emph{The NURBS Book}, 2.\ Aufl., Springer 1997.
  \item G.~Granet, \enquote{Reformulation of the lamellar grating problem through the concept of adaptive spatial resolution}, JOSA A 16, 2510 (1999).
  \item T.~Vallius, M.~Honkanen, \enquote{Reformulation of the Fourier modal method with adaptive spatial resolution: application to multilevel profiles}, Opt.\ Express 10, 24 (2002).
  \item T.~Weiss, G.~Granet, N.~A.~Gippius, S.~G.~Tikhodeev, H.~Giessen, \enquote{Matched coordinates and adaptive spatial resolution in the Fourier modal method}, Opt.\ Express 17, 8051 (2009).
  \item J.~P.~Hugonin, P.~Lalanne, \enquote{Perfectly matched layers as nonlinear coordinate transforms: a generalized formalization}, JOSA A 22, 1844 (2005).
  \item E.~Silberstein, P.~Lalanne, J.~P.~Hugonin, Q.~Cao, \enquote{Use of grating theories in integrated optics}, JOSA A 18, 2865 (2001).
  \item M.~B.~Giles, \enquote{Collected matrix derivative results for forward and reverse mode algorithmic differentiation}, in: Advances in Automatic Differentiation, Springer (2008).
  \item C.~Kim, B.~Lee, \enquote{TORCWA: GPU-accelerated Fourier modal method and gradient-based optimization for metasurface design}, Comput.\ Phys.\ Commun.\ 282, 108552 (2023).
  \item NVIDIA cuSOLVER Documentation, Iterative-Refinement-Löser (\verb|cusolverDnIRSXgesv|, \verb|cusolverDn<t1><t2>gesv|).
\end{enumerate}

%=====================================================================
\clearpage
\section{Erweiterungen ab Version 2.14: Zylinderkoordinaten (M20, M21)}\label{sec:cylext}
%=====================================================================
\providecommand{\hk}[1]{\langle\!\langle #1\rangle\!\rangle}
Die Abschnitte \ref{sec:cylext}--\ref{sec:cylP} spezifizieren die dritte Ausbaustufe: Strukturen mit
\textbf{Rotationssymmetrie} um die Stapelachse $y$ (M20) und mit \textbf{$P$-zähliger Drehsymmetrie} (M21).
Beispiele: Lochblenden und Bullseye-Strukturen in Metallfilmen, Zonenplatten, Mikrolinsen, Säulen, Kugeln und
Kegelstümpfe auf Substraten, kreisförmige Bragg-Gitter, Ring- und Scheibenresonatoren (M20) sowie quadratische und
hexagonale Säulen, Sternstrukturen, Mikropropeller und Einzelelemente von Metaoberflächen (M21).

Die Strukturen sind isoliert (aperiodisch in $x$ und $z$), liegen aber auf einem in $x,z$ unbegrenzten
Schichtstapel. Kartesisch sind sie mit M18 (PML in $x$ und $z$) rechenbar, aber teuer: $n_f=(2N+1)^2$ mit $N$
proportional zur Zellgröße. In Zylinderkoordinaten entkoppeln die azimutalen Harmonischen $e^{im\varphi}$
exakt (M20) bzw.\ in Klassen $m_0+P\mathbb Z$ (M21), und je Harmonischer bleibt ein Problem in $(r,y)$ mit
$n_f=N_k$ Stützstellen eines Hankel-Spektrums. Die Ableitungen sind in dieser Basis diagonal, die
Propagationsgleichung hat dieselbe Form $dF/dy=M(y)F$ wie bisher; RK, RCWA, S-Matrix, Sternprodukt,
Streufeldformulierung, Adjungierte und CUDA-Backend werden unverändert wiederverwendet.

Das Buch behandelt Zylinderkoordinaten in Kap.~X--XV; Abschnitt~\ref{sec:scope} hat sie bisher ausgeschlossen. Die
Formulierung hier folgt der Idee von Bonod, Popov und Nevière (Fourier--Bessel-Basis, Faktorisierungsregeln
[20--22]), ist aber eine \textbf{Eigenherleitung dieses Projekts}
(Galerkin-Matrizen mit analytischem Hintergrund, stetige Hilfsvariablen für die inverse Regel); vor der Umsetzung
mit den Arbeiten vergleichen und Abweichungen hier dokumentieren.

\subsection{Annahmen und Festlegungen}
\begin{enumerate}[nosep]
  \item Die Zylinderachse ist die Stapelachse $y$, Achsenposition $x=z=0$. Alle Schichtgrenzen bleiben Ebenen
        $y=\text{const}$, Substrat und Superstrat bleiben homogen und isotrop.
  \item Materialien in M20 isotrop. Anisotrope Tensoren sind nur zulässig, wenn ihre Komponenten in der Basis
        $(\mathbf e_r,\mathbf e_y,\mathbf e_\varphi)$ nicht von $\varphi$ abhängen (z.\,B.\ radial orientiert einachsig);
        das ist \emph{nicht} Teil von M20/M21.
  \item Der Hintergrund jeder Schicht (Material für $r>R_s$, $R_s$ = Außenradius der Struktur) ist homogen. Alle
        strukturbedingten Anteile haben kompakten Träger $[0,R_s]$.
  \item Anregung: Strahlen mit Hankel-Spektrum (Gauß, benutzerdefiniert) im Totalfeld; ebene Wellen beliebiger
        Richtung ausschließlich über die Streufeldformulierung (Abschnitt~\ref{sec:scatfield}), denn ihr Spektrum ist
        eine Deltafunktion.
  \item Mit ausgeschalteter Option (kein Block \verb|"cylindrical"|) bleiben alle bisherigen Pfade bitgleich
        (Golden Files).
\end{enumerate}

\subsection{Meilensteine und Abhängigkeiten}
\begin{center}
\small
\begin{tabular}{p{1.0cm}p{6.6cm}p{3.0cm}p{3.8cm}}
\toprule
M & Inhalt & setzt voraus & prüfbar in der Sandbox \\
\midrule
M20 & Rotationssymmetrische Strukturen: Fourier--Bessel-Spektrum, Galerkin-Matrizen, Laurent und Li in $r$,
      Strahlen und Streufeld, Fernfeld, Nahfeld, Ringschichten, CAD-Rotationskörper & M18 (Streufeld), M14 (CAD) & ja \\
M21 & $P$-zählige Drehsymmetrie: gekoppelte Harmonische, Winkelkoeffizienten aus Querschnitten (Polygon, Kontur, CAD),
      Symmetrieklassen & M20, M13 & ja \\
\bottomrule
\end{tabular}
\end{center}
Gradienten (M19) werden für beide Meilensteine mitgeführt, sobald die Vorwärtsrechnung abgenommen ist.

%=====================================================================
\section{M20: Rotationssymmetrische Strukturen}\label{sec:cyl}
%=====================================================================

\subsection{Ziel}
Beugung an Strukturen mit $\varepsilon=\varepsilon(r,y)$ auf beliebigen Schichtstapeln, angeregt durch Strahlen auf der
Achse (Gauß, linear, zirkular, radial oder azimutal polarisiert) oder ebene Wellen unter beliebigem Winkel. Ausgaben:
reflektierte, transmittierte und absorbierte Leistung, Streu- und Extinktionsquerschnitt, Winkelspektrum und
Fernfeld, Nahfeld in horizontalen und vertikalen Schnitten.

\subsection{Koordinaten und Maxwellgleichungen}
$x=r\cos\varphi$, $z=r\sin\varphi$ (der Azimut läuft von $x$ nach $z$, wie $\Psi$ in \eqref{eq:incidence}),
$\mathbf e_r=(\cos\varphi,0,\sin\varphi)$, $\mathbf e_\varphi=(-\sin\varphi,0,\cos\varphi)$. Das Tripel
$(\mathbf e_r,\mathbf e_y,\mathbf e_\varphi)$ ist rechtshändig ($\mathbf e_r\times\mathbf e_y=\mathbf e_\varphi$) und
entspricht $(\mathbf e_x,\mathbf e_y,\mathbf e_z)$: \textbf{$x\leftrightarrow r$, $z\leftrightarrow\varphi$}. In normierten
Einheiten \eqref{eq:maxwell_norm} ($\mu=1$, $\mathbf H\equiv\tilde{\mathbf H}$) lauten die Komponenten (Gegenstück zu
\eqref{eq:mx1}--\eqref{eq:mx3}):
\begin{align}
  \partial_yE_\varphi-\tfrac1r\partial_\varphi E_y&=iB_r, &
  \partial_yH_\varphi-\tfrac1r\partial_\varphi H_y&=-iD_r,\label{eq:cyl_mx1}\\
  \tfrac1r\partial_\varphi E_r-\tfrac1r\partial_r(rE_\varphi)&=iB_y, &
  \tfrac1r\partial_\varphi H_r-\tfrac1r\partial_r(rH_\varphi)&=-iD_y,\label{eq:cyl_mx2}\\
  \partial_rE_y-\partial_yE_r&=iB_\varphi, & \partial_rH_y-\partial_yH_r&=-iD_\varphi .\label{eq:cyl_mx3}
\end{align}
Herleitung: Standardrotation in $(\rho,\varphi,Z)$ mit $Z=-y$, $A_Z=-A_y$ (denn $\mathbf e_r\times\mathbf e_\varphi=-\mathbf e_y$).

\subsection{Azimutale Entwicklung und zirkulare Komponenten}
Jede Komponente wird nach $u=\sum_mu_m(r,y)\,e^{im\varphi}$ entwickelt. Hängt $\varepsilon$ nicht von $\varphi$ ab,
kommutiert der Maxwelloperator mit Drehungen um $y$; die Harmonischen entkoppeln exakt, $\partial_\varphi\to im$.
Statt $E_r,E_\varphi$ werden die \textbf{zirkularen Komponenten}
\begin{equation}
  E_\pm=E_r\pm iE_\varphi,\qquad E_x\pm iE_z=e^{\pm i\varphi}E_\pm,\qquad
  E_x=\tfrac12\big(e^{i\varphi}E_++e^{-i\varphi}E_-\big),\quad E_z=\tfrac1{2i}\big(e^{i\varphi}E_+-e^{-i\varphi}E_-\big)
  \label{eq:cyl_pm}
\end{equation}
verwendet (ebenso $H_\pm$, $D_\pm$, $B_\pm$). Da die kartesischen Komponenten glatt sind, verhält sich die
$m$-te Harmonische auf der Achse wie $E_{\pm,m}\sim r^{|m\pm1|}$, $E_{y,m}\sim r^{|m|}$; die natürlichen Basen sind
$J_{m+1}$, $J_{m-1}$ und $J_m$. Aus \eqref{eq:cyl_mx1}--\eqref{eq:cyl_mx3} folgt je Harmonischer
\begin{equation}
\begin{aligned}
  \partial_yE_+&=\big(\partial_r-\tfrac mr\big)E_y-B_+, & \partial_yE_-&=\big(\partial_r+\tfrac mr\big)E_y+B_-,\\
  \partial_yH_+&=\big(\partial_r-\tfrac mr\big)H_y+D_+, & \partial_yH_-&=\big(\partial_r+\tfrac mr\big)H_y-D_-,\\
  B_y&=\tfrac12\Big[\big(\partial_r+\tfrac{m+1}r\big)E_+-\big(\partial_r-\tfrac{m-1}r\big)E_-\Big], &
  D_y&=-\tfrac12\Big[\big(\partial_r+\tfrac{m+1}r\big)H_+-\big(\partial_r-\tfrac{m-1}r\big)H_-\Big].
\end{aligned}
\label{eq:cyl_pmeq}
\end{equation}

\subsection{Hankel-Darstellung und Propagationsgleichung}
Hankel-Transformation der Ordnung $\nu$ (für ganzzahlige $\nu$, $J_{-n}=(-1)^nJ_n$):
\begin{equation}
  \hat f(k)=\int_0^\infty f(r)\,J_\nu(kr)\,r\,dr,\qquad f(r)=\int_0^\infty\hat f(k)\,J_\nu(kr)\,k\,dk,\qquad
  \int_0^\infty f\bar g\,r\,dr=\int_0^\infty\hat f\bar{\hat g}\,k\,dk .
  \label{eq:hankel}
\end{equation}
$E_+,H_+$ werden mit $\nu=m+1$, $E_-,H_-$ mit $\nu=m-1$, $E_y,H_y,D_y$ mit $\nu=m$ transformiert. Mit den
Rekursionen
\begin{equation}
  \Big(\frac{d}{dr}-\frac\nu r\Big)J_\nu(kr)=-kJ_{\nu+1}(kr),\qquad \Big(\frac{d}{dr}+\frac\nu r\Big)J_\nu(kr)=kJ_{\nu-1}(kr)
  \label{eq:besselrec}
\end{equation}
werden alle Ableitungen in \eqref{eq:cyl_pmeq} \textbf{diagonal} in $k$ und bilden jede Komponente auf die Ordnung
der Zielkomponente ab. Mit $\kappa=\operatorname{diag}(k_j)$ (gemeinsames Gitter für alle Ordnungen,
Abschnitt~\ref{sec:cylgrid}) und den Spektren $e_\pm,h_\pm,e_y,h_y,d_y,d_\pm$:
\begin{equation}
  h_y=\tfrac12\kappa(e_++e_-),\quad d_y=-\tfrac12\kappa(h_++h_-),\quad
  \partial_ye_\pm=\mp\kappa e_y\mp h_\pm,\quad \partial_yh_+=-\kappa h_y+d_+,\quad\partial_yh_-=\kappa h_y-d_- .
  \label{eq:cyl_spec}
\end{equation}
Die Faktorisierung (Abschnitt~\ref{sec:cylfact}) liefert $d_\pm=Q_{\pm+}e_++Q_{\pm-}e_-$ und $d_y=Q_{yy}e_y$
(in M20 keine Kopplung $r$--$y$). Mit $\mathcal E=Q_{yy}^{-1}$, also $e_y=-\tfrac12\mathcal E\kappa(h_++h_-)$, und dem
Zustandsvektor $F=(e_+,e_-,h_+,h_-)^T\in\mathbb C^{4n_f}$, $n_f=N_k$:
\begin{equation}
  \frac{dF}{dy}=M F,\qquad
  M=\begin{pmatrix}0&0&\tfrac12\kappa\mathcal E\kappa-\mathbb 1&\tfrac12\kappa\mathcal E\kappa\\
  0&0&-\tfrac12\kappa\mathcal E\kappa&\mathbb 1-\tfrac12\kappa\mathcal E\kappa\\
  Q_{++}-\tfrac12\kappa^2&Q_{+-}-\tfrac12\kappa^2&0&0\\
  \tfrac12\kappa^2-Q_{-+}&\tfrac12\kappa^2-Q_{--}&0&0\end{pmatrix}.
  \label{eq:cyl_M}
\end{equation}
$M$ hat die Form $\begin{pmatrix}0&\mathcal P\\\mathcal R&0\end{pmatrix}$; der reduzierte RCWA-Pfad
\eqref{eq:rcwa_reduced} ist unverändert anwendbar. Anders als \eqref{eq:Mmu1a} enthält $M$ keine Faktoren $i$
(sie stecken in $E_\pm$).
\begin{tcolorbox}[title=Prüfpunkt (Unit-Test \texttt{test\_cyl\_propagation\_matrix})]
Homogenes isotropes Medium ($Q_{\pm\pm}=\varepsilon\mathbb 1$, $Q_{\pm\mp}=0$, $\mathcal E=\mathbb 1/\varepsilon$): je Knoten
gilt $M^2=(k^2-\varepsilon)\mathbb 1$, also $\lambda=\pm i\beta$, $\beta=\sqrt{\varepsilon-k^2}$, und \eqref{eq:cyl_Psi}
sind Eigenvektoren (Prototyp: $4\cdot10^{-16}$). Die Gleichungen \eqref{eq:cyl_pmeq}--\eqref{eq:cyl_M} sind mit dem
Fasertest V21b numerisch geprüft (Abschnitt~\ref{sec:cylfact}, Befund).
\end{tcolorbox}

\subsection{Diskretes Spektrum}\label{sec:cylgrid}
Das Spektrum wird an Knoten $k_j$ mit Gewichten $w_j$ ($j=1\dots N_k$) einer zusammengesetzten Gauß--Legendre-Quadratur
auf $[0,k_{\max}]$ dargestellt (Nyström): $f(r)\approx\sum_jw_jk_j\hat f(k_j)J_\nu(k_jr)$. Alle Ordnungen verwenden
\emph{dieselben} Knoten, damit \eqref{eq:besselrec} exakt diagonal bleibt. Festlegungen (Startwerte, in der
Konvergenzstudie V21h bestätigen):
\begin{itemize}[nosep]
  \item \textbf{Panelgrenzen} bei $0$, $\operatorname{Re}\sqrt{\varepsilon_{\text{sub}}}$, $\operatorname{Re}\sqrt{\varepsilon_{\text{sup}}}$
        (Verzweigungspunkte $\beta=0$ der Halbräume), an den Enden des deformierten Abschnitts (unten) und bei $k_{\max}$.
  \item \textbf{Verzweigungspunkte:} Die Spektren hängen über $\beta_h=\sqrt{\varepsilon_h-k^2}$ wurzelartig von $k$ ab
        (Gegenstück zur Rayleigh-Anomalie). Jedes Panel $[a,b]$ wird in Teilpanels der Breite $h$ (Dichte, unten)
        geteilt; auf einem Teilpanel $[u_a,u_a+h]$, das an einem Verzweigungspunkt beginnt bzw.\ endet, gilt
        $k=u_a+h(1-\cos t)$ bzw.\ $k=u_a+h\sin t$, $t\in[0,\pi/2]$, sodass $dk/dt$ dort wie $\sqrt{|k-k_h|}$ verschwindet.
        Ein Panel zwischen zwei Verzweigungspunkten hat mindestens zwei Teilpanels: die Substitution
        $k=\bar k+\Delta\sin t$ über $[-\pi/2,\pi/2]$ ist mit 8 Knoten für glatte Integranden um den Faktor $10^3$
        ungenauer (v2.15, Unit-Test). Die Integranden werden dadurch glatt in $t$.
  \item \textbf{Dichte:} Panelbreite $\le\pi/(2R_s)$ mit 8 Knoten (die Kerne $\int_0^{R_s}J_\nu(kr)J_{\nu'}(k'r)r\,dr$
        oszillieren in $k$ mit Periode $\approx\pi/R_s$). Damit $N_k\approx2k_{\max}R_s/\pi$ (normierte Einheiten).
        Prototyp (Faser, $R_s=2$): ab $\Delta k\approx0{,}7\approx\pi/(2{,}3R_s)$ dichteunabhängig, der Fehler ist dann
        allein durch $k_{\max}$ bestimmt.
  \item \textbf{Abschneidung:} $k_{\max}=\kappa_{\max}\max_j\operatorname{Re}\sqrt{\varepsilon_j}$ über alle Materialien,
        Standard $\kappa_{\max}=6$; die Rolle von $k_{\max}$ entspricht der Fourierordnung $N$ im kartesischen Pfad.
  \item \textbf{Komplexer Integrationsweg:} Gebundene Moden des Hintergrundstapels (Wellenleiterfilme,
        Oberflächenplasmonen) haben Pole auf der reellen Achse zwischen $k_a=\max_h\operatorname{Re}\sqrt{\varepsilon_h}$ und
        $k_b$ (Standard: $1{,}2\max_j\operatorname{Re}\sqrt{\varepsilon_j}$ über alle Schichten; für Metall/Dielektrikum
        zusätzlich $1{,}2\operatorname{Re}\sqrt{\varepsilon_m\varepsilon_d/(\varepsilon_m+\varepsilon_d)}$). Auf $[k_a,k_b]$:
        $k(t)=t-i\delta\sin\big(\pi\tfrac{t-k_a}{k_b-k_a}\big)$, $dk=k'(t)\,dt$, Standard $\delta=0{,}05(k_b-k_a)$.
        Bei $e^{-i\omega t}$ liegen die Pole verlustbehafteter Moden in der oberen Halbebene; der Weg weicht nach unten
        aus. Auf dem Weg gilt der Wurzelzweig \eqref{eq:beta_branch} unverändert ($\operatorname{Im}\beta^2>0$). Die
        Deformation ist nur nötig, wenn der Hintergrund gebundene Moden trägt (Prüfung: Eigenwerte der Hintergrund-S-Matrix
        im Intervall; sonst Weg reell). Gebundene Moden verlassen die Struktur seitlich; ihr Fluss erscheint nicht in
        oben/unten (Ausgaben).
\end{itemize}

\subsection{Galerkin-Matrizen}\label{sec:cylgal}
Für eine Funktion $g(r)$ (in einer Scheibe) und Ordnungen $\nu$ (Ausgang), $\nu'$ (Eingang):
\begin{equation}
  \hk{g}_{\nu\leftarrow\nu'}=\Big(w_lk_l\int_0^\infty g(r)\,J_\nu(k_jr)\,J_{\nu'}(k_lr)\,r\,dr\Big)_{jl},\qquad
  \hk{g}_{\nu\leftarrow\nu}=g_\infty\mathbb 1+\hk{g-g_\infty}_{\nu\leftarrow\nu}.
  \label{eq:cyl_gal}
\end{equation}
Das ist das Gegenstück zur Toeplitz-Matrix \eqref{eq:toeplitz}. Der konstante Hintergrund $g_\infty$ ist wegen der
Orthogonalität in \eqref{eq:hankel} exakt die Einheitsmatrix und wird \emph{analytisch} eingesetzt; numerisch
integriert wird nur der kompakte Rest über $[0,R_s]$. \textbf{Gemischte Ordnungen ($\nu\ne\nu'$) treten nur mit
kompakt getragenen Funktionen auf} (Abschnitt~\ref{sec:cylfact}); der Hintergrund würde dort eine nichtlokale
Ordnungsumrechnung erfordern, die so vermieden wird.
\begin{itemize}[nosep]
  \item \textbf{Gleiche Ordnung, stückweise konstant} ($g=g_i$ auf $[r_{i-1},r_i]$): Lommel-Integrale,
        \begin{equation}
          \int J_\nu(ar)J_\nu(br)\,r\,dr=\frac{r\big[bJ_\nu(ar)J_\nu'(br)-aJ_\nu'(ar)J_\nu(br)\big]}{a^2-b^2},\qquad
          \int J_\nu(ar)^2\,r\,dr=\frac{r^2}2\Big[J_\nu'(ar)^2+\Big(1-\frac{\nu^2}{a^2r^2}\Big)J_\nu(ar)^2\Big]
          \label{eq:lommel}
        \end{equation}
        (geprüft für $\nu=0,1,3,-2$ und komplexe $a,b$ auf $10^{-16}$). Für $|a-b|\,r_i<10^{-2}$ statt der ersten Formel
        Gauß-Quadratur (Auslöschung).
  \item \textbf{Gemischte Ordnungen und glatte Funktionen:} zusammengesetzte Gauß--Legendre-Quadratur in $r$ zwischen den
        Bruchstellen, Panelbreite $\le\pi/(2k_{\max})$, 16 Knoten; Auswertung als Matrixprodukt
        $T_\nu\,\operatorname{diag}(g(r_s)r_sW_s)\,T_{\nu'}^T\operatorname{diag}(wk)$ mit $T_\nu=(J_\nu(k_jr_s))$.
  \item Bessel-Funktionen komplexen Arguments (deformierter Weg). \textbf{Umsetzung (v2.15):} Miller-Rückwärtsrekursion
        für alle Ordnungen $0\dots n$ zugleich, Start oberhalb des Airy-Übergangs ($n_0=\max(n,|z|)+16(|z|/2)^{1/3}+20$),
        Normierung $e^{\mp iz}=J_0+2\sum_n(\mp i)^nJ_n$ mit $\mp=\operatorname{sign}\operatorname{Im}z$ (wächst wie $J$),
        für reelles $z$: $1=J_0+2\sum_nJ_{2n}$ (Ergebnis exakt reell). Genauigkeit $10^{-13}$ bezogen auf
        $\max(|J_n|,\sqrt{2/(\pi|z|)})$ gegen \verb|scipy.special.jv| bzw.\ mpmath (bei $|z|\approx500$ ist scipy nur auf
        $\approx10^{-14}$ relativ genau, ebenso \verb|std::cyl_bessel_j| von libstdc++ auf $2\cdot10^{-11}$). Die
        Argumente bleiben $|k|r\le k_{\max}R_s$.
\end{itemize}

\subsection{Faktorisierung}\label{sec:cylfact}
\paragraph{Laurent.} $Q_{++}=\hk{\varepsilon}_{m+1\leftarrow m+1}$, $Q_{--}=\hk{\varepsilon}_{m-1\leftarrow m-1}$,
$Q_{+-}=Q_{-+}=0$, $Q_{yy}=\hk\varepsilon_{m\leftarrow m}$. Einfach, aber an Wänden $r=\text{const}$ ist $E_r$ unstetig;
die Konvergenz ist dann erster Ordnung in $1/k_{\max}$ (Befund unten).

\paragraph{Li in $r$ (Standard für Scheiben mit Wänden $r=\text{const}$).} An einer Wand $r=\text{const}$ sind $E_\varphi$,
$E_y$ und $D_r$ stetig, $E_r$ und $D_\varphi$ springen. $E_r$ ist in $E_\pm$ mit $E_\varphi$ gemischt, eine inverse
Regel direkt auf $E_\pm$ ist daher falsch. Stattdessen werden die \emph{stetigen} Hilfsgrößen
\begin{equation}
  G_\pm=\frac{D_r}{\varepsilon_\infty}\pm iE_\varphi\qquad(\text{Ordnung }m\pm1,\ \varepsilon_\infty=\text{Hintergrund der Scheibe})
\end{equation}
eingeführt. Mit $\rho=\varepsilon_\infty/\varepsilon$ gilt exakt
\begin{equation}
  \begin{pmatrix}E_+\\E_-\end{pmatrix}=\frac12\begin{pmatrix}\rho+1&\rho-1\\\rho-1&\rho+1\end{pmatrix}\begin{pmatrix}G_+\\G_-\end{pmatrix},\qquad
  \begin{pmatrix}D_+\\D_-\end{pmatrix}=\frac12\begin{pmatrix}\varepsilon_\infty+\varepsilon&\varepsilon_\infty-\varepsilon\\\varepsilon_\infty-\varepsilon&\varepsilon_\infty+\varepsilon\end{pmatrix}
  \begin{pmatrix}G_+\\G_-\end{pmatrix}.
  \label{eq:cyl_G}
\end{equation}
Jedes Produkt ist unstetig $\times$ stetig, also gilt die Laurent-Regel. Mit $p=m+1$, $q=m-1$, $a=\tfrac12(\rho-1)$,
$b=\tfrac12(\varepsilon-\varepsilon_\infty)$ (beide kompakt getragen):
\begin{equation}
  \mathfrak A=\begin{pmatrix}\mathbb 1+\hk a_{p\leftarrow p}&\hk a_{p\leftarrow q}\\\hk a_{q\leftarrow p}&\mathbb 1+\hk a_{q\leftarrow q}\end{pmatrix},\quad
  \mathfrak B=\begin{pmatrix}\varepsilon_\infty\mathbb 1+\hk b_{p\leftarrow p}&-\hk b_{p\leftarrow q}\\-\hk b_{q\leftarrow p}&\varepsilon_\infty\mathbb 1+\hk b_{q\leftarrow q}\end{pmatrix},\quad
  \begin{pmatrix}Q_{++}&Q_{+-}\\Q_{-+}&Q_{--}\end{pmatrix}=\mathfrak B\,\mathfrak A^{-1}
  \label{eq:cyl_li}
\end{equation}
($\mathfrak A$ mit LU faktorisieren, nicht invertieren). Die inverse Regel steckt in $\mathfrak A^{-1}$. $E_y$ ist an
den Wänden tangential: $Q_{yy}=\hk\varepsilon_{m\leftarrow m}$ (Laurent). Grenzfall $\varepsilon\equiv\varepsilon_\infty$:
$\mathfrak A=\mathbb 1$, $\mathfrak B=\varepsilon_\infty\mathbb 1$ (Unit-Test). Für $m=0$ ist $J_{-1}=-J_1$; die Formeln
gelten unverändert.

\textbf{Regularität auf der Achse.} Mit konstantem $\varepsilon_\infty$ ist $G_\pm$ auf der Achse nicht regulär, wenn dort
ein anderes Material liegt (z.\,B.\ $m=1$: $G_+$ enthält einen Anteil $\sim r^0$ in der $J_2$-Basis). Eine glatte
Referenz $\varepsilon_{\text{ref}}(r)$ (logarithmisch von $\varepsilon(0)$ nach $\varepsilon_\infty$) behebt das, war im
Prototyp aber nicht besser (Tabelle). Standard ist daher konstantes $\varepsilon_\infty$; die glatte Referenz bleibt als
Option \verb|"li_r_reference": "smooth"| für die Konvergenzstudie mit Metallen.

\begin{tcolorbox}[title={Befund Prototyp (Python, \texttt{scripts/proto\_cylindrical.py}): Fasermoden als Eigenwerte von $M$}]
$y$-invariante Scheibe mit Kern $\varepsilon_1=4$, Radius $k_0a=2$, Mantel $\varepsilon_2=1$; geführte Moden sind
Eigenwerte $\pm i\beta$ von \eqref{eq:cyl_M}, Referenz: exakte Eigenwertgleichung der Stufenindexfaser. Fehler
$|\beta-\beta_{\text{exakt}}|$ ($N_k=20k_{\max}$, dichteunabhängig; \enquote{glatt}: $\varepsilon_{\text{ref}}$ auf $[a/2,a]$):
\begin{center}
\begin{tabular}{llccc}
\toprule
Mode & Faktorisierung & $k_{\max}=4$ & $8$ & $16$ \\
\midrule
TM$_{01}$ ($m=0$) & Laurent & $3{,}3\cdot10^{-2}$ & $1{,}6\cdot10^{-2}$ & $8{,}1\cdot10^{-3}$ \\
 & Li in $r$, $\varepsilon_\infty$ & $4{,}8\cdot10^{-3}$ & $5{,}6\cdot10^{-4}$ & $7{,}7\cdot10^{-5}$ \\
 & Li in $r$, glatt & $1{,}3\cdot10^{-2}$ & $2{,}0\cdot10^{-3}$ & $7{,}2\cdot10^{-5}$ \\
\midrule
HE$_{11}$ ($m=1$) & Laurent & $5{,}2\cdot10^{-3}$ & $3{,}1\cdot10^{-3}$ & $1{,}5\cdot10^{-3}$ \\
 & Li in $r$, $\varepsilon_\infty$ & $8{,}7\cdot10^{-4}$ & $3{,}6\cdot10^{-4}$ & $1{,}9\cdot10^{-5}$ \\
 & Li in $r$, glatt & $4{,}8\cdot10^{-3}$ & $5{,}1\cdot10^{-4}$ & $2{,}0\cdot10^{-5}$ \\
\midrule
Hybridmode $m=2$ & Laurent & $2{,}0\cdot10^{-2}$ & $1{,}0\cdot10^{-2}$ & $5{,}4\cdot10^{-3}$ \\
 & Li in $r$, $\varepsilon_\infty$ & $5{,}2\cdot10^{-3}$ & $4{,}9\cdot10^{-4}$ & $1{,}1\cdot10^{-4}$ \\
 & Li in $r$, glatt & $1{,}3\cdot10^{-2}$ & $1{,}7\cdot10^{-3}$ & $9{,}4\cdot10^{-5}$ \\
\bottomrule
\end{tabular}
\end{center}
Laurent konvergiert wie $1/k_{\max}$, Li in $r$ etwa wie $k_{\max}^{-3}$. TE$_{01}$ (nur $E_\varphi$) ist in allen Varianten
identisch. Die Vorzeichen in \eqref{eq:cyl_pmeq}--\eqref{eq:cyl_M} sind damit für $m=0,1,2$ bestätigt.
\end{tcolorbox}

\paragraph{Meridianprofile.} Schichten mit $y$-abhängigen Radien $r_i(y)$ (Kegelstumpf, Kugel, Linse): RCWA-Treppe mit
Li in $r$ (Standard) oder RK mit $M(y)$ aus \eqref{eq:cyl_li} bei den Radien $r_i(y)$. Bei geneigten Wänden ist die
inverse Regel für $E_r$ nur der Treppengrenzfall (Normale $\ne\mathbf e_r$). Eine FFF mit Normalenfeld in der
Meridianebene ($\mathbf N=(N_r,N_y,0)$, Kopplung der Ordnungen $m\pm1$ mit $m$ über $N_rN_y$) ist \emph{nicht} Teil von
M20; Herleitung im Plan nachtragen, falls V21c (Kugel) mit der Treppe zu langsam konvergiert.

\subsection{Homogene Medien, Fluss und Leistungen}
Je Knoten ist $M$ ein $4\times4$-Block. Mit $\mathcal B=\begin{pmatrix}\varepsilon-k^2/2&-k^2/2\\k^2/2&-\varepsilon+k^2/2\end{pmatrix}$
($\mathcal B^2=\varepsilon\beta^2\mathbb 1$):
\begin{equation}
  F=\Psi V,\qquad\Psi=\begin{pmatrix}\mathbb 1&\mathbb 1\\-\mathcal V&\mathcal V\end{pmatrix},\qquad
  \mathcal V=-\frac{i}{\beta}\mathcal B,\qquad\mathcal V^{-1}=\frac{i}{\varepsilon\beta}\mathcal B,
  \label{eq:cyl_Psi}
\end{equation}
$V=(\mathbf a;\mathbf b)$, $\mathbf a=(a_+,a_-)$ abwärts ($\lambda=-i\beta$), $\mathbf b$ aufwärts, wie \eqref{eq:Psi_iso}.
\textbf{TE} ($E_y=0$) $\Leftrightarrow e_+=e_-$, \textbf{TM} ($H_y=0$) $\Leftrightarrow e_+=-e_-$ (s/p-Ausgabe:
$s=\tfrac12(e_++e_-)$, $p=\tfrac12(e_+-e_-)$). Kontrolle $k=0$: abwärts $H_+=i\sqrt\varepsilon E_+$ aus
$\mathbf H=\mathbf k\times\mathbf E$. In homogenen Schichten gilt außerdem $e_y=-\frac{\kappa}{2\varepsilon}(h_++h_-)$.

Poynting-Fluss durch eine Ebene $y=\text{const}$ (aus $(\mathbf E\times\bar{\mathbf H})_y=E_\varphi\bar H_r-E_r\bar H_\varphi$,
$\varphi$-Integration und \eqref{eq:hankel}), summiert über die Harmonischen:
\begin{equation}
  P_y=\frac\pi2\sum_m\operatorname{Im}\sum_jw_jk_j\big(e_{+,j}\bar h_{+,j}-e_{-,j}\bar h_{-,j}\big),
  \label{eq:cyl_flux}
\end{equation}
für eine einzelne Teilwelle in einem verlustfreien Halbraum ($k_j<\sqrt\varepsilon$) in geschlossener Form
\begin{equation}
  p_j=\pm\frac\pi2\,\frac{w_jk_j}{\beta_j}\Big[\big(\varepsilon-\tfrac{k_j^2}2\big)\big(|e_+|^2+|e_-|^2\big)-k_j^2\operatorname{Re}(e_+\bar e_-)\Big]
  \ \ (+\text{ auf}, -\text{ ab}),
  \label{eq:cyl_flux_closed}
\end{equation}
(TE: $\propto2k\beta|e|^2$, TM: $\propto2k\varepsilon|e|^2/\beta$; Prototyp gegen \eqref{eq:cyl_flux} auf $10^{-16}$).
Die Knoten auf $[0,\sqrt\varepsilon_h)$ sind die Ausbreitungswinkel $\theta_j=\arcsin(k_j/\sqrt{\varepsilon_h})$;
evaneszente und deformierte Knoten tragen keinen Fluss.

\subsection{Anregung}\label{sec:cylexc}
\paragraph{Strahlen (Totalfeld).} Vorgegeben ist das tangentiale Feld in der Taillenebene $y_f$ im Superstrat; es bestimmt
das abwärts laufende Spektrum, Phasenbezug an der Oberkante: $a(k)=a_f(k)\,e^{i\beta(y_f-y_{\text{top}})}$, nur
$k<\sqrt{\varepsilon_{\text{sup}}}$ (Apertur $\le1$). Gauß-Strahl mit Taille $w_0$ und Amplitude $E_0$:
$\hat f(k)=E_0\tfrac{w_0^2}2e^{-k^2w_0^2/4}$ (Ordnung 0). Polarisation und Harmonische aus \eqref{eq:cyl_pm}:
\begin{center}
\small
\begin{tabular}{p{2.8cm}p{2.6cm}p{9.6cm}}
\toprule
Polarisation & Harmonische & Spektren \\
\midrule
linear $x$ & $m=-1$ und $m=+1$ & $m=-1$: $a_+=\hat f$ ($J_0$), $a_-=0$; \ $m=+1$: $a_-=\hat f$ ($J_0$), $a_+=0$ \\
zirkular ($E_z=\pm iE_x$) & $m=\pm1$ (eine) & wie oben, nur eine Harmonische, Amplitude $2\hat f$ \\
radial / azimutal & $m=0$ & radial (TM): $a_+=-a_-$, azimutal (TE): $a_+=a_-$; Ordnung $\pm1$, z.\,B.\ $f=r\,e^{-r^2/w_0^2}$ \\
\bottomrule
\end{tabular}
\end{center}
Bei rotationssymmetrischer Struktur folgt $m=+1$ aus $m=-1$ durch Spiegelung ($e_+\leftrightarrow e_-$ mit Vorzeichen nach
$J_{-n}=(-1)^nJ_n$); linear polarisierte Strahlen kosten daher eine Rechnung. Benutzerdefinierte Spektren: Tabelle
$(k,a_+,a_-)$ je $m$, auf die Knoten interpoliert. Hoch fokussierte Strahlen (Richards--Wolf) sind eine spätere Option.

\paragraph{Ebene Welle (Streufeld).} Hintergrundfeld wie in Abschnitt~\ref{sec:scatfield} (analytisch, Airy):
$\mathbf E_{\text{bg}}=e^{i(\alpha_0x+\gamma_0z)}\mathbf e(y)$, $k_\parallel=\sqrt{\alpha_0^2+\gamma_0^2}$, $(\alpha_0,\gamma_0)=k_\parallel(\cos\Psi,\sin\Psi)$.
Mit Jacobi--Anger $e^{ik_\parallel r\cos(\varphi-\Psi)}=\sum_ni^nJ_n(k_\parallel r)e^{in(\varphi-\Psi)}$ und
$c_\pm=e_x\pm ie_z$ (Tangentialkomponenten des Hintergrundes):
\begin{equation}
  E^{\text{bg}}_{+,m}=c_+\,i^{m+1}e^{-i(m+1)\Psi}J_{m+1}(k_\parallel r),\quad
  E^{\text{bg}}_{-,m}=c_-\,i^{m-1}e^{-i(m-1)\Psi}J_{m-1}(k_\parallel r),\quad
  E^{\text{bg}}_{y,m}=e_y\,i^me^{-im\Psi}J_m(k_\parallel r),
  \label{eq:cyl_jacobi}
\end{equation}
$\mathbf H_{\text{bg}}$ analog. Polarisationsquelle $\mathbf P=(\varepsilon-\varepsilon_{\text{bg}})\mathbf E_{\text{bg}}$ ist kompakt
getragen; ihre Spektren sind Spalten von \eqref{eq:cyl_gal} am Zusatzknoten $k'=k_\parallel$ (ohne Gewicht):
$[P_\pm]_j=\int_0^{R_s}(\varepsilon-\varepsilon_{\text{bg}})J_{m\pm1}(k_jr)E^{\text{bg}}_{\pm,m}(r)\,r\,dr$, ebenso $P_y$. Mit
$\mathbf D_s=\varepsilon\mathbf E_s+\mathbf P$ folgt aus \eqref{eq:cyl_spec} die inhomogene Gleichung (Gegenstück zu \eqref{eq:scat_ode})
\begin{equation}
  \frac{dF_s}{dy}=MF_s+\sigma,\qquad \sigma=\big(\kappa\mathcal Ep_y,\ -\kappa\mathcal Ep_y,\ p_+,\ -p_-\big)^T .
  \label{eq:cyl_scat}
\end{equation}
Erweiterte Sternprodukte, RK mit Quellterm und partikuläre RCWA-Lösungen sind die von M18.
\textbf{Konsistenz:} Liegt das Hintergrundfeld auf dem Gitter (Strahlen), wird wie in M18 $\sigma=(M-M_{\text{bg}})F_{\text{bg}}$
gebildet (dieselbe Faktorisierung wie das Totalfeld; V21e auf Rundung). \textbf{Ebene Wellen (v2.16):} Auch hier gibt es
die konsistente Quelle, obwohl das Spektrum ein Delta ist: mit $\varepsilon_{\text{bg}}=\varepsilon_\infty$ (äußerstes
Ringmaterial) und $\mathfrak A=\mathbb 1+\hk{a}$, $\mathfrak B=\varepsilon_\infty\mathbb 1+\hk{b}$ aus \eqref{eq:cyl_li}
ist $\mathfrak A E_{\text{bg}}=E_{\text{bg}}+\hat a$, also $\mathfrak A^{-1}E_{\text{bg}}=E_{\text{bg}}-\mathfrak A^{-1}\hat a$ und
\[
  (p_+;p_-)=(Q-\varepsilon_\infty)E_{\text{bg}}=\hat b-Q\,\hat a ,
\]
$\hat a,\hat b$ die Spalten von $\hk{a}$, $\hk{b}$ (mit den Vorzeichen von $\mathfrak A$, $\mathfrak B$, gemischte Ordnungen
$m\pm1$) am Zusatzknoten $k'=k_\parallel$ ohne Gewicht, angewandt auf $(c_+i^{m+1}e^{-i(m+1)\Psi},\,c_-i^{m-1}e^{-i(m-1)\Psi})$;
die Deltaanteile heben sich auf. Laurent: $(p_+;p_-)$ wie oben (dort ohnehin konsistent), $p_y$ immer Laurent.
Unit-Test: liegt $k_\parallel$ auf einem Knoten, ist $E_{\text{bg}}$ ein Gittervektor und die Identität gilt auf $10^{-12}$.
\textbf{Befund (V21c):} Mit der exakten Quelle konvergiert $\sigma_{\text{sca}}$ der Mie-Kugel nur erster Ordnung in
$k_{\max}$ (20 Scheiben: $+2{,}5\,\%$, $+1{,}5\,\%$ bei $\kappa_{\max}=4,8$; 40 Scheiben, $\kappa_{\max}=12$: $+0{,}7\,\%$), mit der
konsistenten Quelle bleibt im Wesentlichen der Treppenfehler ($+0{,}03\,\%$, $+0{,}3\,\%$).
\textbf{Harmonische:} $|m|\le m_{\text{cut}}=\lceil x+4x^{1/3}+2\rceil$, $x=k_\parallel R_s$ für die Quelle bzw.\
$\max_j\operatorname{Re}\sqrt{\varepsilon_j}\,R_s$ für die Feldauflösung (Wiscombe-Kriterium; Konvergenz in $m_{\text{cut}}$
prüfen). Normaleinfall: nur $m=\pm1$. Die Harmonischen sind unabhängig und laufen parallel über \verb|solveScan|.

\subsection{S-Matrix, Löser und Nahfeld}
S-Matrix \eqref{eq:Sdef} mit $n_h=2N_k$ (Amplituden $a_\pm$, $b_\pm$), Referenzmedium und Grenzflächen wie
Abschnitt~\ref{sec:smatrix} mit \eqref{eq:cyl_Psi}. RK (Abschnitt~\ref{sec:rk}), RCWA (Abschnitt~\ref{sec:rcwa}),
interne Amplituden \eqref{eq:internal}, Parallelität (M15), Genauigkeit (M16) unverändert. Option analog
\verb|split_polarizations|: für $m=0$ zerfällt das System exakt in TE ($e_+=e_-$) und TM ($e_+=-e_-$), je $2N_k$.
\textbf{Nahfeld:} Synthese $E_\pm(r)=\sum_jw_jk_je_{\pm,j}J_{m\pm1}(k_jr)$ (ebenso $E_y$, $H$), Summe über $m$ mit
$e^{im\varphi}$, Ausgabe in kartesischen Komponenten nach \eqref{eq:cyl_pm}. In Scheiben mit Li in $r$: FFF-konsistente
Rekonstruktion wie \eqref{eq:recon}, d.\,h.\ $G_\pm=\mathfrak A^{-1}e_\pm$ synthetisieren (stetig) und $E_\pm$, $D_\pm$
punktweise aus \eqref{eq:cyl_G}. Auf deformiertem Weg wächst $J_\nu(kr)$ wie $e^{|\operatorname{Im}k|r}$: Ausgabe nur
für $r\le r_{\max}$ mit $\delta\,r_{\max}\le\ln10^{3}$, sonst Warnung (kleineres $\delta$ wählen). Schnitte:
horizontal ($y=\text{const}$, Raster in $x$--$z$) und vertikal (Meridianebene $\varphi=\varphi_0$, also $x$--$y$ bei
$\varphi_0=0$); Ausgabeformat wie M9.

\subsection{Ausgaben}
\begin{itemize}[nosep]
  \item \textbf{Strahlen:} $R=P_{\text{auf}}/P_{\text{inc}}$, $T=P_{\text{ab,Substrat}}/P_{\text{inc}}$ aus \eqref{eq:cyl_flux},
        $A=1-R-T$ (enthält bei geführten Moden den seitlichen Abfluss; Hinweis in der Ausgabe).
  \item \textbf{Ebene Wellen:} Streuquerschnitt $\sigma_{\text{sca}}$ (µm$^2$) aus dem Streufeld oben und unten,
        Absorption $\sigma_{\text{abs}}=\tfrac1{I_{\text{inc}}}\cdot\tfrac12\int\operatorname{Im}\varepsilon\,|\mathbf E|^2\,dV$ über die
        strukturierten Bereiche (Gauß-Knoten in $r$, $y$; Summe über $m$), $\sigma_{\text{ext}}=\sigma_{\text{sca}}+\sigma_{\text{abs}}$.
        Im homogenen Hintergrund zusätzlich optisches Theorem als Kontrolle.
  \item \textbf{Winkelspektrum je $m$:} $dP/d\theta$ an den Knoten auf $[0,\sqrt{\varepsilon_h})$, s/p getrennt.
  \item \textbf{Fernfeld:} kartesisches Winkelspektrum aus der zweidimensionalen Fouriertransformation einer Harmonischen,
        $\widehat{f\,e^{in\varphi}}=2\pi(-i)^n\hat f_n(k)\,e^{in\varphi_k}$; mit \eqref{eq:cyl_pm}:
        $\widehat{E_x\pm iE_z}=2\pi\sum_m(-i)^{m\pm1}\hat e_\pm(k)\,e^{i(m\pm1)\varphi_k}$. Raster in $(\theta,\varphi_k)$.
  \item Nahfeld wie oben; Diagnose: Energiebilanz, größter Anteil der höchsten Knoten (Abschneideindikator).
  \item Nicht in M20: seitlicher Fluss gebundener Moden über Residuen (Option, Herleitung bei Bedarf).
\end{itemize}

\subsection{Schichttypen, CAD und Konfiguration}
\begin{itemize}[nosep]
  \item \verb|rings|: Radien $r_1<\dots<r_n$ und $n+1$ Materialien von innen nach außen (letztes $=$ Hintergrund);
        Gegenstück zu \verb|lamellar|.
  \item \verb|rings_profile|: Radien als Funktionen $r_i(y)$ (Kegelstumpf, Kugelkappe, benutzerdefiniert); RCWA-Treppe
        mit dem Scheibenkriterium \eqref{eq:slicecrit} oder RK.
  \item \verb|cad| mit \verb|"axisymmetric": true|: horizontale Schnitte wie M14; jede Schnittregion muss aus Kreisen um
        die Achse bestehen (Toleranz $10^{-9}R_s$), sonst Fehlercode \verb|E_NOT_AXISYMMETRIC| mit Höhe und Abweichung
        (für die Prüfschleife des Sprachmodells, Abschnitt~\ref{sec:cadcheck}). Bandzerlegung und
        Tschebyschow-Interpolation wie M14, interpoliert werden die Radien.
  \item Gemischte Stapel: homogene Schichten wie bisher; kartesische Schichttypen sind im Zylinderpfad unzulässig
        (Fehler), außer \verb|contour2d|/\verb|polygon2d|/\verb|cad| mit $P$-Symmetrie (M21).
\end{itemize}
{\footnotesize\begin{verbatim}
"cylindrical": {
  "symmetry": "axial",                       // M21: {"rotation": P}
  "harmonics": "auto",                       // oder Liste [-1, 1]
  "spectrum": {"kappa_max": 6, "dk_factor": 1.0, "deform": "auto"},  // deform: "auto" | false | {"delta": ...}
  "factorization": "li_r"                    // "laurent" | "li_r"; Option "li_r_reference": "smooth"
},
"incidence": {"type": "gaussian_beam", "waist_um": 2.0, "focus_y_um": 0.5, "polarization": "x"}
             // oder {"type": "plane_wave", "phi_deg": 30, "psi_deg": 0, "delta_deg": 0} -> Streufeld
\end{verbatim}}

\subsection{Architektur und Aufwand}
Neues Modul \verb|fg::cyl| (\verb|include/fg/cyl|, \verb|src/cyl|): \verb|SpectralGrid| (Knoten, Gewichte, Weg,
Substitutionen), \verb|HankelGalerkin| (Matrizen \eqref{eq:cyl_gal}, Lommel, Quadratur, Bessel komplex),
\verb|CylLayerModel| (\verb|rings|, \verb|rings_profile|, CAD), \verb|cyl::propagationMatrix| \eqref{eq:cyl_M},
\verb|cyl::ModeBasis| \eqref{eq:cyl_Psi}, \verb|cyl::Excitation| (Strahlen, Jacobi--Anger), \verb|cyl::Outputs|
(Flüsse, Fernfeld, Nahfeld). Der Agent prüft zuerst, wo \verb|StackSolver| Kartesisches annimmt (\verb|IndexMap|,
\eqref{eq:Psi_iso}, Effizienzen, Nahfeldsynthese), und legt diese Stellen hinter eine Schnittstelle (z.\,B.\
\verb|ModalBasis|: $M(y)$, $\Psi$ der Halbräume und der Referenz, $n_h$, Flussform). Der kartesische Pfad bleibt
bitgleich (Golden Files). Python: \verb|fg.solve| mit Zylinderkonfiguration, Ausgaben zusätzlich
\verb|far_field| und \verb|angular_spectrum|.

\textbf{Richtwerte} (normierte Einheiten, $\kappa_{\max}=6$, $n_{\max}=1{,}5$; kartesisch: Zelle $\approx3R_s$ mit PML,
$N\approx k_{\max}\,\text{Zelle}/2\pi$):
\begin{center}
\small
\begin{tabular}{lccc}
\toprule
$R_s$ bei $\lambda=1\,$µm & $N_k$ (je $m$) & $4n_f$ zylindrisch & $n_f$ kartesisch (M18) \\
\midrule
2\,µm & $\approx70$ & $\approx300$ & $\approx1\cdot10^4$ \\
10\,µm & $\approx360$ & $\approx1400$ & $\approx3\cdot10^5$ (nicht rechenbar) \\
50\,µm & $\approx1800$ & $\approx7000$ & -- \\
\bottomrule
\end{tabular}
\end{center}
Schräger Einfall vervielfacht die Zahl der Rechnungen um $\approx2m_{\text{cut}}+1$ (unabhängig, parallel).

\paragraph{Umsetzung (v2.15).} Schnittstelle \verb|detail::ModalBasis| (\verb|src/solvers/modal_basis.hpp|):
\verb|ref|, \verb|sub|, \verb|M(Schicht, y)|, \verb|modes(M)|, \verb|homogeneousModes(Schicht)|, \verb|rhoEstimate| und
die optionale Polarisation; Schichten als \verb|StackLayer| (Dicke, homogen, $y$-invariant, Scheibenzahl), von dem
\verb|LayerModel| erbt. Die Assemblierung (homogene Schichten, RCWA-Scheiben, RK-Scheiben, Sternprodukt-Kette,
Aufzeichnung) spricht nur mit der Schnittstelle; \verb|CartesianBasis| ersetzt die bisherige Systemklasse (Golden Files
bitgleich), \verb|CylBasis| (\verb|src/solvers/cyl_basis.cpp|) liefert je Harmonischer $m$ die Basen
\eqref{eq:cyl_Psi} von Superstrat und Substrat und $M$ der Ringschichten nach \eqref{eq:cyl_M} (vorab berechnet,
$y$-invariant). Kartesisch bleiben vorerst die Quellen der Streufeldformulierung und die Nahfeldrekonstruktion
(zylindrisch: ebene Wellen und Nahfeld, folgen). Festlegungen dieser Stufe:
\begin{itemize}[nosep]
  \item $R_s=\max(\text{äußerster Ringradius},\,2w_0)$ (ohne Ringe bestimmt der Strahl die Dichte).
  \item Deformierter Weg bei \verb|"deform": "auto"|, wenn eine Schicht einen größeren Realteil der Brechzahl hat als
        beide Halbräume oder ein Metall ($\operatorname{Re}\varepsilon<0$) vorkommt (hinreichend für gebundene Moden;
        unnötige Deformation ist unschädlich, V21f). $\delta$ relativ: \verb|{"delta": 0.05}| bedeutet $0{,}05(k_b-k_a)$.
  \item Lineare Polarisation: beide Harmonischen $m=\pm1$ werden gerechnet (die Spiegelung ist noch nicht genutzt; ihre
        Gleichheit ist Testkriterium).
  \item \textbf{Ebene Wellen (v2.16):} Hintergrund $=$ Stapel mit dem äußersten Ringmaterial je Ringschicht, gelöst wie
        M18 mit dem kartesischen Löser bei einer Ordnung (Airy); Quelle je Ringschicht nach \eqref{eq:cyl_scat} mit der
        konsistenten Quelle (oben) als \verb|ModalBasis|-unabhängiger Hook der Streufeld-Assemblierung
        (partikuläre RCWA-Lösungen, RK mit Quellterm, erweiterte Sternprodukte unverändert). Harmonische \verb|"auto"|:
        Normaleinfall $m=\pm1$, sonst $|m|\le m_{\text{cut}}(k_\parallel R_s)$ (darüber verschwindet die Quelle).
        Ausgabe vorerst $\sigma_{\text{sca}}$ oben/unten je Harmonischer aus \eqref{eq:cyl_flux} des Streufeldes,
        bezogen auf die Intensität der einfallenden Welle; $\sigma_{\text{abs}}$, $\sigma_{\text{ext}}$ und das optische
        Theorem folgen mit dem Nahfeld.
  \item \textbf{Befund (V21a, dielektrische Säule):} Laurent erhält die Energie exakt ($|R+T-1|\le7\cdot10^{-13}$), denn
        $\operatorname{diag}(wk)\hk{g}$ ist symmetrisch; Li in $r$ ($\mathfrak B\mathfrak A^{-1}$ nicht selbstadjungiert)
        nur bis auf den Abschneidefehler: $7\cdot10^{-4}$, $1{,}7\cdot10^{-5}$, $1{,}5\cdot10^{-5}$ bei
        $\kappa_{\max}=2,4,8$ (Sockel vermutlich die Achsenregularität von $G_\pm$, s.\,o.). $R$ konvergiert mit Li
        dennoch deutlich schneller (Laurent erster Ordnung).
\end{itemize}

\subsection{Differenzierbarkeit}
Die Rückwärtsregeln von M19 gelten für $M$, $\mathfrak A^{-1}$, S-Matrizen und Sternprodukte unverändert. Neu sind
Formableitungen der Galerkin-Matrizen nach einem Wandradius $r_i$ (Leibniz):
\begin{equation}
  \frac{\partial}{\partial r_i}\hk{g}_{\nu\leftarrow\nu'}=\big(g(r_i^-)-g(r_i^+)\big)\,\big(w_lk_l\,J_\nu(k_jr_i)J_{\nu'}(k_lr_i)\,r_i\big)_{jl},
  \label{eq:cyl_shape}
\end{equation}
dazu Materialableitungen (linear in $g$) und Ableitungen der Quelle \eqref{eq:cyl_jacobi} nach $k_\parallel$, $\Psi$.

\subsection{Tests und Abnahme (M20)}
\begin{itemize}[nosep]
  \item Unit: Prüfpunkt zu \eqref{eq:cyl_M}; \eqref{eq:lommel} gegen Quadratur ($5\cdot10^{-14}$, auch komplex; Auslöschung an der Schwelle $|a-b|\,r=10^{-2}$);
        Hintergrundidentität; \eqref{eq:cyl_li} mit $\varepsilon\equiv\varepsilon_\infty$; Hankel-Paar \eqref{eq:hankel} für
        $r^\nu e^{-r^2}$ ($10^{-12}$); Bessel komplex gegen Referenz; \eqref{eq:cyl_flux_closed} $=$ \eqref{eq:cyl_flux}.
  \item \textbf{V21a} (Schichtstapel, Gauß-Strahl, $m=\pm1$ und $m=0$): $R$, $T$ gegen knotenweise Fresnel/Airy $10^{-12}$;
        verlustfrei $R+T=1$ auf $10^{-12}$.
  \item \textbf{V21b} (Fasermoden, s.\ Befund): Eigenwerte der Scheibe gegen die exakte Eigenwertgleichung; Li in $r$
        $\le2\cdot10^{-4}$ bei $k_{\max}=8n_{\text{Kern}}$ für TM$_{01}$, HE$_{11}$, $m=2$; Laurent nachweislich erster Ordnung.
  \item \textbf{V21c} (Mie-Kugel): dielektrische Kugel ($n=1{,}5$, $k_0a=3$) im homogenen Medium, ebene Welle über
        Streufeld, Treppe mit Li in $r$: $\sigma_{\text{sca}}$, $\sigma_{\text{ext}}$ gegen Mie (Bohren--Huffman, eigene
        Referenzimplementierung im Test) nach Extrapolation in Scheibenzahl und $k_{\max}$ auf $10^{-3}$; Goldkugel qualitativ.
        \textbf{Umsetzung (v2.16):} Ordnung $p$ in $N$ aus $\kappa_{\max}=4$ ($N=10,20,40$; $p\approx1{,}84$), Richardson bei
        $\kappa_{\max}=8$ ($N=20,40$), Rest in $\kappa_{\max}$ als $\kappa^{-2}$: $-5{,}7\cdot10^{-4}$ ($\sigma_{\text{ext}}$ und
        Goldkugel mit dem Nahfeld).
  \item \textbf{V21d} (gegen kartesisch): runde Säule und rundes Kontaktloch (wie V20g) bei Normaleinfall: Streu- bzw.\
        Transmissionsleistung $=$ M18-Ergebnis innerhalb dessen Konvergenz ($10^{-3}$).
  \item \textbf{V21e} (Streufeld $=$ Totalfeld): Gauß-Strahl, Streufeld $+$ Hintergrund $=$ Totalfeld (RCWA $10^{-10}$,
        RK auf RK-Toleranz), Laurent und Li in $r$.
  \item \textbf{V21f} (Symmetrie und Weg): $m\leftrightarrow-m$ $10^{-12}$; Ergebnisse unabhängig von $\delta$ ($10^{-8}$) für
        eine Lochblende im Silberfilm; mit verlustbehaftetem Metall und $\operatorname{Im}\varepsilon\to$ klein konvergiert
        der reelle Weg gegen den deformierten (falsche Ausweichrichtung $\Rightarrow$ Sprung, Test schlägt fehl).
  \item \textbf{V21g} (Lochblende im Metallschirm): qualitativer Vergleich mit Popov/Bonod [22] (Werte vor
        Verwendung bibliographisch prüfen).
  \item \textbf{V21h} (Konvergenzstudie): $N_k$-Dichte, $k_{\max}$, $m_{\text{cut}}$ und Scheibenzahl als CSV, Tabelle in
        \verb|docs/validation.md|.
  \item \textbf{V21i} (Gradienten): Radius, Material, Dicke gegen Differenzenquotienten ($10^{-6}$).
  \item Alle bisherigen Tests bitgleich (Golden Files).
\end{itemize}

%=====================================================================
\section{M21: $P$-zählige Drehsymmetrie}\label{sec:cylP}
%=====================================================================

\subsection{Ziel}
Strukturen, die bei Drehung um $2\pi/P$ um die $y$-Achse in sich übergehen ($P\ge1$; $P=1$ ist der allgemeine Fall
eines beliebigen isolierten 3D-Objekts). Die Querschnitte $y=\text{const}$ werden wie bisher als \verb|polygon2d|,
\verb|contour2d| oder \verb|cad| beschrieben, um die Achse zentriert.

\subsection{Winkelentwicklung und Kopplung}
\begin{equation}
  \varepsilon(r,\varphi,y)=\sum_q\varepsilon_q(r,y)\,e^{iqP\varphi},\qquad
  \varepsilon_q=\frac1{2\pi}\int_0^{2\pi}\varepsilon\,e^{-iqP\varphi}d\varphi .
  \label{eq:cylP_eps}
\end{equation}
Ein Produkt koppelt die Harmonische $m$ nur mit $m-qP$. Die Harmonischen zerfallen in $P$ unabhängige \textbf{Klassen}
$m_l=m_0+lP$, $l\in\mathbb Z$, $m_0\in\{\lceil-P/2\rceil+1,\dots,\lfloor P/2\rfloor\}$; abgeschnitten wird
$|l|\le L$. Für $q\ne0$ ist $\varepsilon_q$ kompakt getragen (der Hintergrund ist rotationssymmetrisch), $\varepsilon_0$
ist Hintergrund plus kompakter Rest. Alle Galerkin-Blöcke außerhalb der Diagonale in $l$ sind daher kompakt, wie in
Abschnitt~\ref{sec:cylgal} gefordert.

\subsection{Gekoppelte Propagationsgleichung}
Zustandsvektor $F=(e_+,e_-,h_+,h_-)$ mit $e_\pm=(e^{(l)}_\pm)_{l=-L..L}$, je Block $N_k$ Knoten, $n_f=(2L+1)N_k$.
Die Ableitungsanteile hängen nach \eqref{eq:besselrec} nicht von $m$ ab (die Ordnungsbuchhaltung steckt in der Basis):
in \eqref{eq:cyl_M} wird $\kappa$ durch $\mathbb 1_{2L+1}\otimes\kappa$ ersetzt. Laurent-Blöcke:
\begin{equation}
  \big(Q_{++}\big)_{ll'}=\hk{\varepsilon_{l-l'}}_{m_l+1\leftarrow m_{l'}+1},\quad
  \big(Q_{--}\big)_{ll'}=\hk{\varepsilon_{l-l'}}_{m_l-1\leftarrow m_{l'}-1},\quad
  \big(Q_{yy}\big)_{ll'}=\hk{\varepsilon_{l-l'}}_{m_l\leftarrow m_{l'}},\quad Q_{\pm\mp}=0 .
  \label{eq:cylP_Q}
\end{equation}
Mit $\mathcal E=Q_{yy}^{-1}$ (voll, $n_f\times n_f$) gilt \eqref{eq:cyl_M} unverändert. Index: $(l,j)\mapsto lN_k+j$ über eine
Klasse \verb|HarmonicMap| (Gegenstück zu \verb|IndexMap|).
\begin{tcolorbox}[title={Befund Prototyp ($P=1$): exzentrischer Faserkern}]
Kern wie in V21b, um $x_0$ (normiert) aus der Achse verschoben, Klasse $m_0=1$, Laurent, $k_{\max}=3$, $N_k=60$. Die
Moden müssen mit denen des zentrierten Kerns übereinstimmen. Größte Abweichung des HE$_{11}$-Paares vom zentrierten
Laurent-Wert $1{,}7167536$: $x_0=0{,}3$: $8\cdot10^{-5}$ ($L=3$), $4\cdot10^{-8}$ ($L=5$); $x_0=0{,}6$: $7\cdot10^{-7}$
($L=5$), $5\cdot10^{-8}$ ($L=7$). Die Konvergenz in $L$ ist exponentiell; Kopplung, Ordnungen und Vorzeichen von
\eqref{eq:cylP_Q} sind damit bestätigt.
\end{tcolorbox}

\subsection{Winkelkoeffizienten aus Querschnitten}
Für eine Region $\mathcal R$ (Polygon, \verb|Contour2D|, CAD-Schnitt) ist $\chi_q(r)=\frac1{2\pi}\int_{\text{Bögen}}e^{-iqP\varphi}d\varphi$
über die Bögen des Kreises $|\mathbf x|=r$ in $\mathcal R$:
\begin{equation}
  \chi_0(r)=\frac{\text{Bogenlänge}}{2\pi r},\qquad
  \chi_q(r)=\frac1{2\pi}\sum_{\text{Bögen }[\varphi_a,\varphi_b]}\frac{e^{-iqP\varphi_b}-e^{-iqP\varphi_a}}{-iqP}\quad(q\ne0),
  \label{eq:cylP_chi}
\end{equation}
$\varepsilon_q=\varepsilon_{\text{bg}}\delta_{q0}+\sum_b(\varepsilon_b-\varepsilon_{\text{bg}})\chi^{(b)}_q$.
\begin{itemize}[nosep]
  \item Schnittpunkte Kreis--Rand: Polygonkanten quadratisch; rationale Bézierstücke $\mathbf C=\mathbf P(t)/w(t)$:
        Nullstellen von $|\mathbf P(t)|^2-r^2w(t)^2$ (Grad $\le2p\le18$, Bernstein-Unterteilung mit Newton-Politur). Innen/außen
        der Bögen über die Bogenmitte und den bestehenden Punkttest.
  \item \textbf{Kritische Radien:} Eckenradien, Endpunkte der Segmente, lokale Extrema von $|\mathbf C(t)|$ (Tangentialpunkte);
        dazwischen ist $\chi_q$ glatt. Quadratur in $r$ zwischen kritischen Radien mit Gauß--Legendre wie in
        Abschnitt~\ref{sec:cylgal}, adaptiv (Vergleich mit doppelter Knotenzahl, $10^{-12}$). Nahe Tangentialradien
        verhält sich $\chi_q$ wie $\sqrt{|r-r_t|}$: Panelgrenze genau dort.
  \item Liegt die Achse in der Region, gilt $\chi_q\to\delta_{q0}$ für $r\to0$.
  \item Symmetrieprüfung: Region um $2\pi/P$ gedreht $=$ Region (Randstichproben, Toleranz $10^{-9}R_s$), sonst Fehler
        (CAD: Code \verb|E_ROT_SYMMETRY|). Spiegelsymmetrie $\varphi\to-\varphi$ wird erkannt ($\chi_q$ reell und gerade) und
        für die Klassenreduktion genutzt.
\end{itemize}

\subsection{Faktorisierung}
M21 verwendet \textbf{Laurent} \eqref{eq:cylP_Q} (Pflicht). Wände sind im Allgemeinen nicht $r=\text{const}$; Li in $r$
gilt nur für den rotationssymmetrischen Anteil und ist allein nicht konsistent. FFF als \textbf{Stufe 2} (nicht Teil der
Abnahme): Normalenfeld aus Algorithmus~NF im Querschnitt, in zirkularen Komponenten
$N_\pm=e^{\mp i\varphi}(N_x\pm iN_z)$, $E_N=\tfrac12(N_-E_++N_+E_-)$, und stetige Hilfsgrößen nach dem Muster von
\eqref{eq:cyl_G} mit $(D_N/\varepsilon_\infty,\ \mathbf E_T)$; die Winkelkoeffizienten von $N_\pm N_\pm$ und $N_\pm N_\mp$
koppeln dann $e_+$ und $e_-$ über die Klassenblöcke. Herleitung und Prototyp (Fasertest mit quadratischem Kern gegen
kartesische RCWA-Moden) vor der Umsetzung im Plan nachtragen; Anlass ist V22c mit Metall.

\subsection{Anregung und Symmetrieklassen}
Strahlen auf der Achse: $m=\pm1$ (linear) bzw.\ eine davon (zirkular), $m=0$ (radial/azimutal); jede Harmonische
gehört zu genau einer Klasse $m_0\equiv m\pmod P$ (für $P=2$ liegen $m=\pm1$ in derselben Klasse und werden gemeinsam
gerechnet). Ebene Wellen: alle $|m|\le m_{\text{cut}}$, also alle $P$ Klassen mit $L\approx m_{\text{cut}}/P$; die Quelle
\eqref{eq:cyl_jacobi} wird je Klasse nach $l$ sortiert und mit den Blöcken $\hk{\varepsilon_q-\varepsilon_{\text{bg}}\delta_{q0}}$ am
Zusatzknoten $k_\parallel$ gebildet. Startwert für die Feldauflösung $L=\lceil m_{\text{cut}}/P\rceil+2$, Konvergenz in $L$ prüfen.

\subsection{Aufwand}
Je Klasse $4(2L+1)N_k$ Unbekannte, $P$ Klassen. Gegenüber M20 kostet eine Klasse $\approx(2L+1)^3$-mal so viel wie eine
Harmonische; der Gewinn gegenüber dem kartesischen Pfad wächst mit $R_s$ und $P$, weil $N_k$ linear in $R_s$ wächst
statt $n_f$ quadratisch. Für kleine Objekte (wenige Wellenlängen) ohne Rotationssymmetrie ($P=1$) ist der kartesische
Pfad (M18) meist günstiger; die Wahl bleibt beim Anwender, \verb|fgsolve --info| gibt beide Größen aus.

\subsection{CAD und Konfiguration}
\verb|"symmetry": {"rotation": P}| im Block \verb|"cylindrical"|; zulässig sind \verb|polygon2d|, \verb|contour2d|,
\verb|cad| (Schnitte wie M14, Achse bei $x=z=0$) und \verb|rings|. Die CAD-Prüfung meldet zusätzlich die tatsächlich
gefundene Symmetrieordnung (größtes $P$ mit erfüllter Drehsymmetrie), damit das Sprachmodell $P$ korrekt setzt.

\subsection{Tests und Abnahme (M21)}
\begin{itemize}[nosep]
  \item Unit: \eqref{eq:cylP_chi} für Kreis (nur $q=0$), Quadrat und regelmäßiges Sechseck gegen numerische Quadratur
        ($10^{-13}$); Bézier-Kreis $=$ exakter Kreis; Symmetrieprüfung lehnt ein gestörtes Sechseck ab.
  \item \textbf{V22a}: rotationssymmetrische Struktur über den $P$-Pfad (beliebiges $P$) $=$ M20 ($10^{-12}$).
  \item \textbf{V22b} (exzentrischer Kern, $P=1$, s.\ Befund): Moden $=$ zentrierter Kern ($10^{-7}$ bei $L=7$).
  \item \textbf{V22c} (gegen kartesisch): quadratische ($P=4$) und hexagonale ($P=6$) dielektrische Säule auf Substrat,
        Normaleinfall: Streuquerschnitt $=$ \verb|polygon2d| mit PML (M18) innerhalb dessen Konvergenz ($10^{-3}$),
        beide Laurent; Silbersäule als Konvergenzstudie (Anlass für FFF Stufe~2).
  \item \textbf{V22d} (Auswahlregeln): Blöcke zwischen Klassen exakt null; Drehung des Einfallsazimuts $\Psi$ um $2\pi/P$
        ändert $\sigma_{\text{sca}}$, $\sigma_{\text{abs}}$ nicht ($10^{-12}$), Fernfeld gedreht.
  \item \textbf{V22e} (Spiegelsymmetrie): Klassen $m_0$ und $-m_0$ ineinander überführt ($10^{-12}$).
  \item \textbf{V22f} (CAD): hexagonale Säule aus Nurbforge $=$ \verb|polygon2d|-Pfad ($10^{-10}$); \verb|E_ROT_SYMMETRY|
        bei gestörter Geometrie; gemeldete Symmetrieordnung korrekt.
  \item \textbf{V22g} (Gradienten): Kontrollpunkte eines \verb|contour2d|-Querschnitts (Formableitung über die
        Bogenenden in \eqref{eq:cylP_chi}) gegen Differenzenquotienten ($10^{-6}$).
\end{itemize}

\subsection*{Literatur zu M20/M21 (vor Verwendung bibliographisch prüfen)}
\begin{enumerate}[label={[\arabic*]},nosep,leftmargin=2em,start=20]
  \item N.~Bonod, E.~Popov, M.~Nevière, \enquote{Differential theory of diffraction by finite
        cylindrical objects}, JOSA A 22, 481--490 (2005).
  \item E.~Popov, M.~Nevière, N.~Bonod, \enquote{Factorization of products of discontinuous functions
        applied to Fourier--Bessel basis}, JOSA A 21, 46--51 (2004).
  \item E.~Popov, N.~Bonod, \enquote{Differential theory of diffraction in cylindrical coordinates},
        phys.\ stat.\ sol.\ (b) (2007), doi:10.1002/pssb.200743151.
  \item P.~Boyer, E.~Popov, M.~Nevière, G.~Renversez, \enquote{Diffraction theory: application of the fast Fourier
        factorization to cylindrical devices with arbitrary cross section lighted in conical mounting}, JOSA A 23, 1146 (2006).
  \item A.~W.~Snyder, J.~D.~Love, \emph{Optical Waveguide Theory}, Chapman \& Hall 1983 (Eigenwertgleichung der Stufenindexfaser, V21b).
  \item C.~F.~Bohren, D.~R.~Huffman, \emph{Absorption and Scattering of Light by Small Particles}, Wiley 1983 (Mie, V21c).
  \item W.~J.~Wiscombe, \enquote{Improved Mie scattering algorithms}, Appl.\ Opt.\ 19, 1505 (1980) (Abschneidekriterium $m_{\text{cut}}$).
  \item NIST Digital Library of Mathematical Functions, §10.6 (Rekursionen), §10.22 (Integrale, Lommel).
\end{enumerate}

}
\end{document}
